%
% 構文解析の仕組み
%
% B5変形判
%%
\documentclass[a4paper,9pt,openright,tombow,]{asciibook}
\usepackage{lmodern}
\usepackage[match]{luatexja-fontspec}

% 用紙サイズ、版面設計

% bxpapersize で紙のサイズ（A4）と、ターゲット（B5変型）を明示
%\usepackage[size=a4, layout={182mm,234mm}, layoutcenter]{bxpapersize}

% 印刷位置の調整
\advance\hoffset 14mm
\advance\voffset 32mm


\usepackage[
% 版面（B5変型）
  paperwidth=182mm,
  paperheight=234mm,
  % A4の中央にB5変型（182x234）を配置し、さらにその中で
  % 「上25mm、下25mm、左20mm、右20mm」の余白を取る場合の絶対マージン計算値
  % 左: 14 + 20 = 34mm / 右: 14 + 20 = 34mm
  % 上: 31.5 + 25 = 56.5mm / 下: 31.5 + 25 = 56.5mm
  inner=20mm,
  outer=20mm,
  top=25mm,
  bottom=25mm
]{geometry}

% 明朝体（リュウミン）の直接指定
\setmainjfont[
  Path = /Users/kahei/Library/Fonts/MorisawaFonts/, % フォルダのパス
  BoldFont = AP-OTF-RyuminPr6N-Bold.otf,            % 太字のファイル名
%  ItalicFont = AP-RyuminPr6N-R.otf                  % 必要なら指定
]{AP-OTF-RyuminPr6N-Regular.otf}                    % 標準のファイル名

% ゴシック体（新ゴ）の直接指定
\setsansjfont[
  Path = /Users/kahei/Library/Fonts/MorisawaFonts/,
  BoldFont = AP-OTF-ShinGoPr6N-Bold.otf
]{AP-OTF-ShinGoPr6N-Medium.otf}

\usepackage{luatexja-ruby} % ルビ
\usepackage{graphicx}
\usepackage{xcolor}
\usepackage{setspace}
\usepackage{array} % 表組みの罫線を改善する
\usepackage{listings}
% listingsの設定
\lstset{
  basicstyle=\ttfamily\small, % ここで上記2つをまとめた \ttfamily を適用
  columns=fixed,               % 文字幅を完全に固定する設定
  keepspaces=true,             % スペースを維持
  showstringspaces=false       % 文字列中のスペースに記号を出さない
%  frame=tb
}


% tikzで図を描く
\usepackage{tikz}
\usetikzlibrary{shapes,arrows,positioning,calc,automata,chains,fit,decorations.pathmorphing}


\usepackage{xurl} % \urlコマンド中で改行できるようにする
\usepackage[unicode]{hyperref}
\hypersetup{%
 bookmarksnumbered=true,%
% colorlinks=true,%
 colorlinks=false,%
 hidelinks,%
 linkcolor=blue,%
 urlcolor=blue,%
 citecolor=blue,%
 linktocpage=true,%
 setpagesize=false}

%
\usepackage[most,listings]{tcolorbox}
\tcbuselibrary{breakable}
\tcbuselibrary{skins}

\newtcolorbox{notebox}[1]{%
colframe=black, colback=white,
coltitle=black, colbacktitle=white,
boxrule=0.8pt, arc=0mm,
fonttitle=\sffamily\bfseries,
enhanced,
attach boxed title to top left={xshift=10mm,yshift=-3mm},
boxed title style={frame hidden},
title=#1}

% Note環境の定義
% 使用するときは、\begin{important}{}〜\end{important}のように
% 必ず{}をつけること。つけないと本文の1文字目が削除
% されてしまう。
\newenvironment{important}[1]{%
  \begin{tcolorbox}[fontlower=\fontsize{8pt}{12pt}\selectfont \sffamily,before skip=4.5mm,after skip=4.5mm,blanker,sidebyside,lefthand width=1cm]%
  \includegraphics[width=1.5cm]{img/aside-icons/important.pdf}%
  \tcblower
  }%
  {\end{tcolorbox}}

\newenvironment{tip}[1]{%
  \begin{tcolorbox}[fontlower=\fontsize{8pt}{12pt}\selectfont \sffamily,before skip=4.5mm,after skip=4.5mm,blanker,sidebyside,lefthand width=1cm]%
  \includegraphics[width=1.5cm]{img/aside-icons/tip.pdf}%
  \tcblower
  }%
  {\end{tcolorbox}}

\newenvironment{info}[1]{%
  \begin{tcolorbox}[fontlower=\fontsize{8pt}{12pt}\selectfont \sffamily,before skip=4.5mm,after skip=4.5mm,blanker,sidebyside,lefthand width=1cm]%
  \includegraphics[width=1.5cm]{img/aside-icons/info.pdf}%
  \tcblower
  }%
  {\end{tcolorbox}}

% 角丸のコード用環境「codescreen」を定義
\newtcblisting{codescreen}[1]{
  breakable,
  arc=3mm,                      % 角丸の半径を指定
  boxrule=0.5pt,                % 枠線の太さ
  colback=gray!5,               % 背景色
  colframe=gray!80,             % 枠線の色
  listing only,
  listing options={
    basicstyle=\ttfamily\small, % フォント設定
    breaklines=true             % 長い行の自動折り返し
  },
  title=#1,                     % タイトル（キャプション）の設定
  % タイトルが不要な場合は title=#1 の行を削除してください
}



% 下線を引く
\usepackage{lua-ul}

% 部の番号を大文字のローマ数字（Ⅰ, Ⅱ, Ⅲ...）にする設定
\renewcommand{\thepart}{\Roman{part}}

% widowの禁止
\widowpenalty=10000
\clubpenalty=10000

%
% 索引
\usepackage{makeidx}
\makeindex

%
% 本文構成
%
\begin{document}
{\pagestyle{frontheadings}
\thispagestyle{empty}
\vbox to 0mm{\vspace*{-32.5truemm}
\hbox to 0mm{
\hspace*{-22.5truemm}\includegraphics[width=186mm]{koubun_tobira.pdf}
}\vss}

\pagebreak

% \thispagestyle{frontheadings}
\vbox to200mm{\hsize120mm
\vfil
\begin{minipage}[b]{12cm}
\textgt{\textbf{商標}}
\begin{spacing}{0.8}
\begin{small}
\noindent{}本文中に記載されている社名および商品名は、一般に開発メーカーの登録商標です。\\
[-1mm]なお、本文中では\texttrademark ・\copyright ・\textregistered 表示を明記しておりません。
\end{small}
\end{spacing}
\end{minipage}}
\pagebreak
}

\pagestyle{asciiheadings}
\tableofcontents
\chapter*{序章：\\構文解析の世界へようこそ}\label{preface}\addcontentsline{toc}{chapter}{序章：構文解析の世界へようこそ}\markboth{序章}{構文解析の世界へようこそ}
\index{こうぶんかいせき@構文解析}
\index{ぱーしんぐ@Parsing}
\index{げんごしょりけい@言語処理系}
\index{こんぱいら@コンパイラ}

本章では、構文解析とはどのような技術なのかを概観します。まずその歴史をたどり、次に構文解析という営みがそもそも何をするものなのかを確かめます。続いて身近な応用例を通して、この技術がどれほど広い範囲で使われているのかを見ていきます。そして最後に、一見すると単純に思えるこの処理が実はどれほど奥深いものであるかを、ひとつの短いプログラムを題材に示します。

\section*{構文解析の歴史的背景}\label{0.1}\addcontentsline{toc}{section}{構文解析の歴史的背景}\markboth{序章}{構文解析の歴史的背景}
\index{こうぶんかいせき@構文解析!歴史}
\index{けいしきげんごりろん@形式言語理論}
\index{ちょむすきーかいそう@チョムスキー階層}
\index{ぶんみゃくじゆうぶんぽう@文脈自由文法}
\index{えるあーるほう@LR法}
\index{えるえるほう@LL法}
\index{BNF}
\index{Yacc}
\index{Lex}
\index{ANTLR}
\index{PEG}

構文解析の歴史は、コンピュータサイエンスの黎明期にまで遡ります。1950年代から1960年代にかけて、初期のコンパイラ開発者たちは、人間が書いたプログラムをコンピュータが理解できる形に変換する方法を模索していました。当時、プログラミング言語の文法を厳密に定義し、それを機械的に処理する方法は確立されていませんでした。

転機となったのは、言語学者ノーム・チョムスキーが1956年に発表した形式言語理論\footnote{Noam
  Chomsky, ``Three models for the description of language'', IRE
  Transactions on Information Theory, 2(3), 1956.
  なお、文脈自由文法のより精密な定式化は1959年の ``On certain formal
  properties of grammars'' で行われています。}でした。チョムスキーは言語を4つの階層（チョムスキー階層）に分類し、その中で「文脈自由文法」（第3章で解説）がプログラミング言語の記述に適していることが明らかになりました。この理論的基盤により、構文解析は科学的なアプローチが可能な分野となったのです。

構文解析アルゴリズムの発展において重要な貢献をした人物たちがいます。1965年、Donald
Knuth（ドナルド・クヌース）はLR法（第4章で解説）を発表しました。LR法の名前は「Left-to-right
scan, Rightmost
derivation」の略で、入力を左から右に読み進めながら最右導出を逆順にたどる、という意味を持ちます。これは効率的な構文解析を可能にする画期的なアルゴリズムでした。一方、LL法（第4章で解説）は1960年代後半にPhilip
M. Lewis II（フィリップ・ルイス）とRichard E.
Stearns（リチャード・スターンズ）によって体系化されました。こちらの名前は「Left-to-right
scan, Leftmost
derivation」の略です。LL法は1960年代初頭から各種コンパイラで使われていた再帰下降構文解析に理論的基礎を与え、より単純で理解しやすい手法として普及しました。これらのアルゴリズムは、現在でも多くの構文解析器の基礎となっています。

初期のプログラミング言語の開発においても、構文解析は重要な役割を果たしました。John
Backus（ジョン・バッカス）は1950年代にFORTRANの開発を主導し、Peter
Naur（ピーター・ナウア）とともにBNF（Backus-Naur
Form）（第1章で解説）という文法記述法を確立しました。BNFは、プログラミング言語の文法を形式的に記述するための標準的な記法となり、現在でも広く使われています。また、Grace
Hopper（グレース・ホッパー）は初期のコンパイラ開発のパイオニアとして、人間が理解しやすい言語から機械語への変換技術の基礎を築きました。

1970年代に入ると、Stephen C.
Johnson（スティーブン・ジョンソン）によってYacc（Yet Another Compiler
Compiler）が開発され、Mike Lesk（マイク・レスク）とEric
Schmidt（エリック・シュミット）によってLex（Lexical Analyzer
Generator）が開発されました。これらのツールは、文法定義から自動的に構文解析器を生成することを可能にし、コンパイラ開発の民主化に大きく貢献しました。現在でも、これらのツールの思想は多くのパーサージェネレータに受け継がれています。

現代では、ANTLR（ANother Tool for Language
Recognition）のような、より強力で使いやすいツールが登場し、構文解析はますます身近なものになっています。また、Parsing
Expression
Grammar（PEG）（第4章で解説）のような新しい形式文法も提案され、構文解析の世界は今なお進化を続けています。

\section*{構文解析とは何か}\label{0.2}\addcontentsline{toc}{section}{構文解析とは何か}\markboth{序章}{構文解析とは何か}
\index{こうぶんかいせき@構文解析!定義}
\index{ちゅうしょうこうぶんぎ@抽象構文木}
\index{AST}
\index{ParseResult}
\index{こうぶんぎ@構文木}

このテーマについては別の立場からの意見もあるかもしれませんが、筆者は構文解析を以下のように定義しています。

構文解析とは、入力された文字列が特定の文法に従っているかどうかを判断し、その構造を抽出するプロセスです。プログラミング言語においては、ソースコードを解析して文法に従っているかを真偽値で返したり、抽象構文木（AST:
Abstract Syntax
Tree）を生成することが一般的な目的です。抽象構文木については第1章の後半で詳しく説明します。

文法に従っているかを真偽値で返す場合、構文解析のためのメソッドはJavaでは以下のように表現できます。

\begin{codescreen}{}
boolean parse(String input);
\end{codescreen}

文法に従っているかを真偽値で返すだけではなく、構文解析の結果として抽象構文木を生成する場合、メソッドは以下のように表現できます。ここで、記述を簡便にするため、Java
17から正式導入されたsealed interfaceとJava
16から正式導入されたrecordを使っています。

\begin{codescreen}{}
interface Tree {}
sealed interface ParseResult 
  permits ParseResult.Success, ParseResult.Failure {
  record Success(Tree ast) implements ParseResult {}
  record Failure(String errorMessage) implements ParseResult {}
}
ParseResult parse(String input);
\end{codescreen}

インタフェース\texttt{Tree}は抽象構文木の基底クラスであり、\texttt{ParseResult}は構文解析の結果を表すクラスです。\texttt{ParseResult.Success}は構文解析が成功した場合の結果を表し、抽象構文木を含みます。一方、\texttt{ParseResult.Failure}は構文解析が失敗した場合の結果を表し、エラーメッセージを含みます。

構文解析によって得られる抽象構文木を処理することで、さまざまな操作が可能になります。たとえば、抽象構文木を使ってコードの最適化や変換、静的解析、コード生成などを行えます。

皆さんが何かしらの言語（プログラミング言語に限りません。HTMLやMarkdown、JSONなども含みます）で書かれたテキストを扱っているとき、構文解析器は必ずどこかで動いています。たとえば、JavaScriptのコードをブラウザで実行する際、ブラウザはまずJavaScriptのコードを構文解析し、抽象構文木を生成してから実行します。また、JSONを扱うAPIでは、JSON文字列を構文解析して抽象構文木に変換し、その後の処理に利用します。

\section*{構文解析の身近な応用例}\label{0.3}\addcontentsline{toc}{section}{構文解析の身近な応用例}\markboth{序章}{構文解析の身近な応用例}
\index{IDE}
\index{りんたー@リンター}
\index{ふぉーまったー@フォーマッター}
\index{Markdown}
\index{JSON}

構文解析は、皆さんが日常的に使っているツールの中で活躍しています。いくつか具体例を見てみましょう。

\begin{description}
  \item[\textbf{\sffamily IDE}の機能：] Visual Studio CodeやIntelliJ
IDEAなどの統合開発環境では、コードを書いている最中にリアルタイムで構文解析が行われています。コード補完（入力途中で候補を表示）、コードの構造表示やナビゲーション、エラーの即座の検出などは構文解析によって実現されています。
\item[リンターとフォーマッター：]
コードの品質チェックや自動整形を行うツールが昨今一般的に使われています。たとえば、ESLintやPrettierなどがあります。これらも内部で構文解析を行っています。コードを抽象構文木に変換し、その構造を解析することで、潜在的な問題を検出したり、一貫したスタイルに整形できるのです。
\item[マークダウンパーサー：]
ドキュメントを執筆する際に使われるMarkdownも、構文解析の対象です。\texttt{**太字**}や\texttt{*斜体*}、見出しやリストなどの記法を解釈し、HTMLに変換するためには、Markdownの文法に従った構文解析が必要です。
\item[設定ファイルの解析：]
\texttt{.env}ファイルの環境変数、\texttt{.gitignore}のパターン、\texttt{package.json}の依存関係など、さまざまな設定ファイルも構文解析の対象です。
\item[データベースクエリ：]
SQLクエリを実行する際、データベースエンジンはまずクエリを構文解析し、実行計画を立てます。\texttt{SELECT\ *\ FROM\ users\ WHERE\ age\ \textgreater{}\ 18}のようなクエリも、構文解析を経て初めて実行可能になります。
\item[正規表現エンジン：]
正規表現そのものも、実は小さな言語です。\texttt{/\^{}{[}a-zA-Z0-9{]}+@{[}a-zA-Z0-9{]}+\textbackslash{}.{[}a-zA-Z{]}+\$/}のようなパターンを解釈するには、正規表現の文法に従った構文解析が必要です。
\end{description}

これらの例からわかるように、構文解析は特別な分野の技術ではなく、私たちの開発を支える基盤技術なのです。

\section*{構文解析の奥深さ}\label{0.4}\addcontentsline{toc}{section}{構文解析の奥深さ}\markboth{序章}{構文解析の奥深さ}

皆さんがもっとも身近に使っている言語の1つであるJavaScriptを例にとって、構文解析の奥深さについて説明してみます。

以下のJavaScriptプログラムを構文解析することにします。

\begin{codescreen}{}
x = 1 +
  2
\end{codescreen}

このコードは、JavaScriptエンジンによって \texttt{x\ =\ 1\ +\ 2;}
と解釈され、直観的には以下のような抽象構文木が結果として返ってくるのが正しいように思われます（抽象構文木の説明は、いったんおいておきます）。

\begin{center}
\begin{tikzpicture}[
  level distance=1.5cm,
  sibling distance=2cm,
  every node/.style={circle,draw,minimum size=0.8cm},
  root/.style={fill=gray!50},
  internal/.style={fill=brown!50},
  leaf/.style={fill=green!30}
]
  \node[root] {=}
    child { node[leaf] {x} }
    child { node[internal] {+}
      child { node[leaf] {1} }
      child { node[leaf] {2} }
    };
\end{tikzpicture}
\end{center}

ところで、JavaScriptには\textgt{自動セミコロン挿入}（ASI: Automatic
Semicolon Insertion）というルールがあり、文末のセミコロンを書き忘れた場合に処理系が暗黙的にセミコロンを補ってくれます。一見便利な機能ですが、ASIのルールは「行終端文字の後に続くトークンが、現在の文法的文脈で許されない場合に限り、行終端文字の位置にセミコロンを挿入する」というもので、時として直感に反する結果を生みます。

たとえば、上記のコードと一見似ている次の例は、実は
\textbf{\sffamily ASI}\textgt{が効かない} ためそのまま \texttt{x\ =\ 1\ +\ 2}
と解釈されます。

\begin{codescreen}{}
x = 1
+ 2
\end{codescreen}

\texttt{x\ =\ 1} の次の行が \texttt{+} で始まっていても、\texttt{+}
は二項演算子として前の式の継続と見なせるため、ASIによるセミコロン挿入は起こらないのです。

一方、\texttt{return}文の直後で改行すると、ASIが効いて意図しない結果になります。

\begin{codescreen}{}
return
a + b;
\end{codescreen}

これは \texttt{return;\ a\ +\ b;} と解釈され、関数は \texttt{undefined}
を返してしまいます。\texttt{return}
の直後にどんなトークンが来ても文として続けられないため、ASIが発動するのです。さらに、次のような行頭の
\texttt{[} も典型的な落とし穴です。

\begin{codescreen}{}
var a = 1, b = 2
b
[a, b].forEach(x => console.log(x))  // b[a, b].forEach(...) と解釈される
\end{codescreen}

\texttt{b} の次に \texttt{{[}} が来ても、\texttt{b{[}...{]}}
というインデックスアクセスとして文を継続できてしまうため、ASIは発動せず、意図と異なる解析結果になります。このような複雑さは、構文解析器の実装を難しくする一因となります。ASIによる予期せぬ挙動を避けるために、行頭にセミコロンを記述するスタイル（いわゆるセミコロン前置記法）を採用する開発者もいます。

2000年以降に登場して普及した言語は、このような特徴を持っていることが多いです。たとえば、Go、Swift、Kotlin、Scalaの文法はこのような特徴を持っています。より古い言語であるRubyやPythonも同じ特徴を持っています。このように、ちょっとした違いで構文解析結果が変わってしまう例は珍しくありません。

このような文法の進化の背景には、プログラマにとって「書きやすく読みやすい」文法を提供するという意図があります。このような「改行に敏感な」文法があると、構文解析器が複雑になりがちです。プログラミング言語の設計者にとっては考慮すべき点が増えるので、このような文法を嫌う言語設計者もいますが、利用者にとってはセミコロンの入力を省略できるなどの利便性があるため、広く採用されています。

しかし「構文解析が複雑になる」というのはいったいどういうことでしょうか。ほとんどの方はピンとこないのではないでしょうか。この本では、このような問いに対して一定の答えを提示することを目指します。「構文解析の複雑さ」と一言で言っても、その要因はさまざまです。たとえば、文法が曖昧であるか（一つの文に対して複数の解釈が可能か）、解析にバックトラック（試行錯誤）が必要か、どれだけの先読みトークン（入力の一部を先に見る数）が必要か、といった要素が複雑さに関わってきます。

本書では、これらの概念に触れながら、構文解析の奥深さを探求していきます。これらの複雑さを理解する上で重要な役割を果たすのが、「文脈自由文法」という考え方です。文脈自由文法とは、簡単に言えば、プログラムの構造を定義するための一連の書き換え規則のことです。たとえば、「数値は式である」や「式は、式と演算子と式からなる」といったルールを形式的に記述するものです。詳細については第3章で詳しく解説しますが、この文脈自由文法という道具立てがあることで、構文解析のさまざまな側面を体系的に議論できるようになります。

読者の皆さんは既存の言語に歯がゆさを感じたことはないでしょうか。たとえば、昨今はJSONやYAMLを使って領域特化言語（DSL: Domain Specific Language）を提供することが一般的になっています。Amazon Web
Services（AWS）のCloudFormationやGoogle Cloud
Platform（GCP）のDeployment
Managerなどがその例です。これらのDSLは、JSONやYAMLという既存のデータ形式を利用して、特定のドメインに特化した言語を提供しています。

しかし、これらのDSLはJSONやYAMLの枠に収めるために無理をしており、どうしても「不自然」な文法になってしまいます。JSONやYAMLは世界中に普及していますから、JSONやYAMLを使ったDSLを提供する合理性はあるものの、場合によってはJSONやYAMLに縛られないDSLの文法を考案することが求められることもあります。そのためには構文解析の知識が必要です。

構文解析は単に実用的であるにとどまらず非常に「楽しい」テーマです。読者の皆さんには本書を通じて、ぜひ「構文解析の面白さ」を感じていただければと思います。

さて、序章を読み終えた皆さんは、構文解析というテーマに対してどのような印象を持たれたでしょうか。もしかしたら、「普段何気なく書いているプログラムの裏側では、こんな複雑なことが行われているのか」と驚かれたかもしれません。あるいは、「自分の手で言語のルールを定義し、それを解釈するプログラムを作るのは面白そうだ」と感じた方もいるかもしれません。あなたがこれまでに触れたことのある言語で、構文が原因で混乱した経験はありますか？もしあれば、それはなぜだったのか、この本を読み進める中でヒントが見つかるかもしれません。

次の第1章では、いよいよ具体的な構文解析の世界に足を踏み入れ、簡単な算術式を題材に、手を動かしながら構文解析の第一歩を体験していただきます。お楽しみに！


\begin{tcolorbox}[breakable,title=現代の構文解析の課題]
  構文解析技術は長い歴史を持ちますが、現代においては新たな課題に直面しています。

\begin{description}
  \item[エラーリカバリー：]
従来の構文解析器は、文法エラーに遭遇すると即座に処理を停止していました。これは、最初の文法エラーが後続の文法エラーを引き起こす上に、後続の文法エラーが実態を表していないことが多かったための措置でした。しかし、IDEでは入力途中のコードを常に解析する必要があります。そのため、エラーがあっても可能な限り解析を続行し、意味のあるエラーメッセージを提供する「エラーリカバリー」技術が重要になっています。
\item[インクリメンタル構文解析：]
ファイル全体を毎回解析し直すのは非効率です。変更された部分だけを再解析する「インクリメンタル構文解析」により、大規模なファイルでもリアルタイムな解析が可能になります。
\item[\textbf{\sffamily Language Server Protocol（LSP）}：]
Microsoftが提唱したLSPは、言語サーバーとエディタ間の通信プロトコルです。実質的にはIDEのための基盤プロトコルと言ってもよいでしょう。構文解析結果を効率的に共有することで、一つの言語サーバーが複数のエディタに対応できるようになりました。
\item[\textbf{\sffamily AI}との融合：]
GitHub CopilotやChatGPTのようなAIツールは、コードの構造を理解する必要があります。構文解析によって得られた抽象構文木は、AIがコードを\textbf{理解}するための重要な手がかりとなっています。
\end{description}

これらの課題は、構文解析が単なる「文字列から木構造への変換」以上の役割を担っていることを示しています。
\end{tcolorbox}

\include{ch01}
\chapter{JSONの構文解析}\label{ch02}
\index{JSON}
\index{こうぶんかいせき@構文解析!JSON}

第1章では構文解析に必要な基本概念について学びました。この章ではJSONという実際に使われている言語を題材に、より実践的な構文解析のやり方を学んでいきます。

\section{JSONの概要}
\index{JSON!概要}
\index{JavaScript}
\index{ECMA-404}
\index{RFC 8259}

JSON（JavaScript Object
Notation）は、WebサービスにアクセスするためのAPIで一般的に使われているデータフォーマットです。また、企業内サービス間で連携するときにもよく使われます。皆さんは何らかの形でJSONに触れたことがあるのではないかと思います。

JSONは元々はJavaScriptのオブジェクトリテラル構文を参考にして設計されたデータ形式ですが、現在はECMA-404およびIETFのRFC 8259\footnote{\url{https://datatracker.ietf.org/doc/html/rfc8259}}で標準化されており、いろいろな言語でJSONを扱うライブラリがあります。なお、JSONはJavaScriptの厳密な部分集合ではなく、両者の仕様は別々に管理されている独立した規格です（U+2028/U+2029の扱いなど、歴史的に細かな差異がありました）。また、JSONはデータ交換用フォーマットの中でも非常にシンプルであるという特徴があり、そのシンプルさ故か、同じ言語でもJSONを扱うライブラリが乱立するほどです。今のWebアプリケーション開発に携わる開発者にとってJSONは避けて通れないといってよいでしょう。

以降では簡単なJSONのサンプルを通してJSONの概要を説明します。

\subsection{オブジェクト}
\index{JSON!オブジェクト}
\index{おぶじぇくと@オブジェクト}
\index{ぺあ@ペア}

以下は、2つの名前/値のペアからなる\textgt{オブジェクト}です。

\begin{codescreen}{}
{
  "name": "Kota Mizushima",
  "age":  41
}
\end{codescreen}

このJSONは、\texttt{name}と\texttt{"Kota\ Mizushima"}という文字列の\textgt{ペア}と、\texttt{age}と\texttt{41}という数値の\textgt{ペア}からなる\textgt{オブジェクト}であることを示しています。

なお、用語については、ECMA-404の仕様書に記載されているものに準拠しています。名前/値のペアは、属性やプロパティと呼ばれることもあるので、適宜読み替えてください。
日本語で表現すると、このオブジェクトは、名前が\texttt{Kota\ Mizushima}、年齢が\texttt{41}という人物一人分のデータを表していると考えられます。オブジェクトは、\texttt{\{\}}で囲まれた、\texttt{"name":value}の対が\texttt{,}を区切り文字として続く形になります。後述しますが、\texttt{name}の部分は\textgt{文字列}である必要があります。

\subsection{配列}
\index{JSON!配列}
\index{はいれつ@配列}

別の例として、以下のJSONを見てみます。

\begin{codescreen}{}
{
  "kind":"Rectangle",
  "points": [
    {"x":0,   "y":0  },
    {"x":0,   "y":100},
    {"x":100, "y":100},
    {"x":100, "y":0  }
  ]
}
\end{codescreen}

このJSONは、

\begin{itemize}
\item
  \texttt{"kind"}と\texttt{"Rectangle"}のペア
\item
  \texttt{"points"}と\texttt{{[}...{]}}のペア
\end{itemize}

からなるオブジェクトです。さらに、\texttt{"points"}に対応する値が\textgt{配列}になっていて、その中に以下の4つの要素が含まれています。

\begin{itemize}
\item
  オブジェクト： \texttt{\{"x":0,\ \ \ "y":0\}}
\item
  オブジェクト： \texttt{\{"x":0,\ \ \ "y":100\}}
\item
  オブジェクト： \texttt{\{"x":100,\ "y":100\}}
\item
  オブジェクト： \texttt{\{"x":100,\ "y":0\}}
\end{itemize}

配列は、\texttt{{[}{]}}で囲まれた要素の並びで、区切り文字は\texttt{,}です。

このオブジェクトは、種類が四角形で、それを構成する点が\texttt{(0,\ 0),\ (0,\ 100),\ (100,\ 100),\ (100,\ 0)}からなっているデータを表現しているとみることができます。

\subsection{数値}

これまで見てきたオブジェクトと配列は複合的なデータでしたが、すでに出てきているように、JSONにはこれ以上分解できないデータもあります。先ほどから出てきている数値もそうです。数値は、

\begin{codescreen}{}
1
10
100
1000
1.0
1.5
\end{codescreen}

のような形になっており、整数または小数です。JSONの数値は10進表記ですが、構文解析後にそれを2進浮動小数点数として扱うか、10進小数として扱うかは規定されていません。

\subsection{文字列}

JSONのデータには文字列もあります。

\begin{codescreen}{}
"Hello, World"
"Kota Mizushima"
"hogehoge"
\end{codescreen}

のように、\texttt{""}で囲まれたものが文字列となります。JSONの仕様では、オブジェクトのキーは必ずダブルクォートで囲まれた文字列でなければなりません。たとえば、以下は\textbf{\sffamily JavaScript}\textgt{の}オブジェクトリテラルとしては許容されますが（\texttt{name}
が識別子の命名規則に合致するため）、\textbf{\sffamily JSON}の定義には従っていません。

\begin{codescreen}{}
{
  name: "Kota Mizushima",
  age:  41
}
\end{codescreen}

\begin{tcolorbox}[title=JSONとJavaScriptオブジェクトの違い]
JSONとJavaScriptのオブジェクトリテラルは似ていますが、いくつかの重要な違いがあります。以下の点に注意してください。

\begin{itemize}
\item
  キーはJSONでは必ずダブルクォートで囲まれた文字列です（例：\texttt{"name"}）。JavaScriptでは識別子であればダブルクォート省略可能です。
\item
  コメントはJSONには存在しません（\texttt{//} や \texttt{/*\ ...\ */}
  は不可）。
\item
  末尾カンマはJSONでは不可です（JavaScriptの一部構文では許容される場合あり）。
\item
  値の種類も異なります。\texttt{undefined}/\texttt{NaN}/\texttt{Infinity}
  などはJSONの値としては不正です。
\end{itemize}
\end{tcolorbox}

\subsection{真偽値}

JSONには、多くのプログラミング言語にある真偽値もあります。JSONの真偽値は以下のように、\texttt{true}または\texttt{false}の2通りです。

\begin{codescreen}{}
true
false
\end{codescreen}


真偽値も解釈方法は定められていませんが、ほとんどのプログラミング言語で、該当するリテラル表現があるので、おおむねそのような真偽値リテラルにマッピングされます。

\subsection{null}

多くのプログラミング言語にある要素ですが、JSONには\texttt{null}もあります。多くのプログラミング言語のJSONライブラリでは、無効値に相当する値にマッピングされますが、JSONの仕様では\texttt{null}の解釈は定められていません。\texttt{null}に相当するリテラルがあればそれにマッピングされることも多いですが、\texttt{Option}や\texttt{Maybe}といったデータ型によって\texttt{null}を表現する言語では、そのようなデータ型にマッピングされることが多いようです。

\subsection{JSONの全体像}

ここまでで、JSONで現れる6つの要素について説明しましたが、JSONで使える要素は\textgt{これだけ}です。このシンプルさが、多くのプログラミング言語でJSONが使われる要因でもあるのでしょう。JSONで使える要素について改めて並べてみます。

\begin{itemize}
\item
  オブジェクト
\item
  配列
\item
  数値
\item
  文字列
\item
  真偽値
\item
  \texttt{null}
\end{itemize}

次の節では、このJSONの\textgt{文法}が、どのような形で表現できるかについて見ていきます。

\section{JSONのBNF}
\index{JSON!BNF}
\index{BNF}
\index{とーくん@トークン}
\index{くうはく@空白}

前の節でJSONの概要について説明し終わったところで、いよいよJSONの文法について見ていきます。JSONの文法はECMA-404の仕様書に記載されていますが、ここでは、それを若干変形したBNFで表現されたJSONの文法を見ていきます。

JSONのBNFによる定義を簡略化したものは以下ですべてです。特に小数点以下の部分は煩雑になる割に本質的でないので削除しました。本来のJSONでは文字列のエスケープシーケンス（\texttt{\textbackslash{}n}、\texttt{\textbackslash{}t}、\texttt{\textbackslash{}"}など）を使えますが、本章では扱いません。これらは構文解析の本質を理解する上では必須ではないためです。

ただし、\texttt{"\""}という表記は、BNFでダブルクォートを表現するためのエスケープシーケンスであり、JSON文字列のエスケープシーケンスとは関係ないことに注意してください。

\begin{codescreen}{}
json = ws value;
value = true | false | null | number | string | object | array;
object = LBRACE RBRACE | LBRACE pair {COMMA pair} RBRACE;
pair = string COLON value;
array = LBRACKET RBRACKET | LBRACKET value {COMMA value} RBRACKET ;
string = "\"" {CHAR} "\"" ws;
number = INT ws;
true = "true" ws;
false = "false" ws;
null = "null" ws;

COMMA = "," ws;
COLON = ":" ws;
LBRACE = "{" ws;
RBRACE = "}" ws;
LBRACKET = "[" ws;
RBRACKET = "]" ws;

ws = {" " | "\t" | "\n" | "\r"} ;
CHAR = "a" | "b" | "c" | "d" | "e" | "f" | "g" | "h" | "i" | "j" | "k" | "l" | "m" |
       "n" | "o" | "p" | "q" | "r" | "s" | "t" | "u" | "v" | "w" | "x" | "y" | "z" |
       "A" | "B" | "C" | "D" | "E" | "F" | "G" | "H" | "I" | "J" | "K" | "L" | "M" |
       "N" | "O" | "P" | "Q" | "R" | "S" | "T" | "U" | "V" | "W" | "X" | "Y" | "Z" |
       "0" | "1" | "2" | "3" | "4" | "5" | "6" | "7" | "8" | "9" |
       " " | "!" | "#" | "$" | "%" | "&" | "'" | "(" | ")" | "*" | "+" | "," |
       "-" | "." | "/" | ":" | ";" | "<" | "=" | ">" | "?" | "@" | "[" | "]" |
       "^" | "_" | "`" | "{" | "|" | "}" | "~" ;
INT = ["-"] ("0" | (NONZERO {DIGIT})) ;
DIGIT = "0" | "1" | "2" | "3" | "4" | "5" | "6" | "7" | "8" | "9" ;
NONZERO = "1" | "2" | "3" | "4" | "5" | "6" | "7" | "8" | "9" ;
\end{codescreen}

これまで説明したJSONの要素と比較して見慣れない記号が出てきましたが、1つ1つ見ていきましょう。

なお、規則名の大文字・小文字の使い分けについてですが、\texttt{COMMA}、\texttt{LBRACE}のような大文字で書かれた規則は、後の字句解析の節で出てくる「トークン」に対応します。これらは単一の記号や固定の文字列を表すもので、構文解析の際には1つの単位として扱われます。一方、\texttt{json}、\texttt{value}のような小文字の規則は、より複雑な構造を表します。補助規則として、\texttt{ws}（空白）、\texttt{CHAR}（印字可能文字の集合）、\texttt{INT}（整数）、\texttt{DIGIT}（数字1桁）、\texttt{NONZERO}（非ゼロの数字）を定義しています。

\subsection{json}

一番上から読んでいきます。第1章の復習になりますが、BNFでは、

\begin{codescreen}{}
json = ws value;
\end{codescreen}

のような\textgt{規則}の集まりによって、文法を表現します。各規則はセミコロン（\texttt{;}）で終わります。\texttt{=}の左側である\texttt{json}が\textgt{左辺}で、右側（ここでは
\texttt{ws\ value}）が\textgt{本体}になります。さらに、本体の中に出てくる、他の規則を参照する部分（ここでは\texttt{value}や\texttt{ws}）を非終端記号と呼びます。非終端記号は同じBNFで定義されている規則の左辺と一致する必要があります。

この規則を日本語で表現すると、「\texttt{json}という名前の規則は、\texttt{ws}の後に\texttt{value}が続く」と読めます。\texttt{value}は、JSONの値を表しているので、\texttt{json}という規則は\texttt{ws}（空白文字）の後にJSONの値が続くものを表しています。

\subsection{object}

\texttt{object}はJSONのオブジェクトを表す規則で、定義は以下のようになっています。

\begin{codescreen}{}
object = LBRACE RBRACE | LBRACE pair {COMMA pair} RBRACE;
\end{codescreen}

\texttt{pair}の定義はのちほど出てきますので心配しないでください。

この規則は、\texttt{object}が

\begin{itemize}
\item
  ブレースで囲まれたもの（\texttt{LBRACE\ RBRACE}）である

  \begin{itemize}
  \item
    \texttt{LBRACE}はLeft-Brace（開き波カッコ）の略で\texttt{\{}を示しています
  \item
    \texttt{RBRACE}はRight-Brace（閉じ波カッコ）の略で\texttt{\}}を示しています
  \end{itemize}
\item
  \texttt{LBRACE}が来た後に、\texttt{pair}が1回出現して、さらにその後に、「\texttt{COMMA}（カンマ）とそれに続く\texttt{pair}」というペアが0回以上繰り返して出現した後、\texttt{RBRACE}が来る
\end{itemize}

\noindent のどちらかであることを表しています。ここで \texttt{\{\ ...\ \}}
は「0回以上の繰り返し」を表します（第1章のBNFの説明を参照）。

具体的なJSONを当てはめてみましょう。以下のJSONは\texttt{LBRACE\ RBRACE}にマッチします。

\begin{codescreen}{}
{}
\end{codescreen}


以下のJSONは\texttt{LBRACE\ pair\ \{COMMA\ pair\}\ RBRACE}にマッチします。

\begin{codescreen}{}
{"x":1}
{"x":1,"y":2}
{"x":1,"y":2,"z":3}
\end{codescreen}


しかし、以下のテキストは、\texttt{object}に当てはまらず、エラーになります。\texttt{\{COMMA\ pair\}}とあるように、カンマは後ろにペアを必要とするからです。

\begin{codescreen}{}
{"x":1,} // ,で終わっている
\end{codescreen}

\subsection{pair}

\texttt{pair}（ペア）は、JSONのオブジェクト内での\texttt{"x":1}に当たる部分を表現する規則です。\texttt{value}の定義については後述します。

\begin{codescreen}{}
pair = string COLON value;
\end{codescreen}

これによってペアは\texttt{:}（\texttt{COLON}）の前に文字列リテラル（\texttt{string}）が来て、その後にJSONの値（\texttt{value}）が来ることを表しています。\texttt{pair}にマッチするテキストとしては、

\begin{codescreen}{}
"x":1
"y":true
\end{codescreen}

などがあります。一方で、以下のテキストは\texttt{pair}にマッチしません（JavaScriptのオブジェクトとJSONの違いについては前述のコラム「JSONとJavaScriptオブジェクトの違い」を参照）。

\begin{codescreen}{}
x:1 // 文字列リテラルでないといけない
\end{codescreen}

\subsection{COMMA}

\texttt{COMMA}は、カンマを表す規則です。カンマそのものを表すには、単に\texttt{","}と書けばいいのですが、任意個の空白文字が続くことを表現したいため、規則\texttt{ws}（後述）を参照しています。

\begin{codescreen}{}
COMMA = "," ws;
\end{codescreen}

\subsection{array}

\texttt{array}は、JSONの値の配列を表す規則です。

\begin{codescreen}{}
array = LBRACKET RBRACKET | LBRACKET value {COMMA value} RBRACKET ;
\end{codescreen}

\texttt{LBRACKET}は開き大カッコ（\texttt{{[}}）を、\texttt{RBRACKET}は閉じ大カッコ（\texttt{{]}}）を表しています。\texttt{value}の定義については後述します。

これは、\texttt{array}が

\begin{itemize}
\item
  大カッコで囲まれたもの（\texttt{LBRACKET\ RBRACKET}）である
\item
  開き大カッコ（\texttt{LBRACKET}）が来た後に、\texttt{value}が1回現れ、さらにその後に、「\texttt{COMMA}（カンマ）とそれに続く\texttt{value}」というペアが0回以上繰り返して現れた後、閉じ大カッコが来る（\texttt{RBRACKET}）
\end{itemize}

\noindent のどちらかであることを表しています。よく見ると、先ほどの\texttt{object}と同様の構造を持っていることがわかります。

\texttt{array}についても具体的なJSONを当てはめてみましょう。以下のJSONは\texttt{LBRACKET\ RBRACKET}にマッチします。

\begin{codescreen}{}
[]
\end{codescreen}


また、以下のJSONは\texttt{LBRACKET\ value\ \{COMMA\ value\}\ RBRACKET}にマッチします。

\begin{codescreen}{}
[1]
[1, 2]
[1, 2, 3]
["foo"]
\end{codescreen}

しかし、以下のテキストは、\texttt{array}に当てはまらず、エラーになります。\texttt{\{COMMA\ value\}}とあるように、カンマは必ず後ろに\texttt{value}を必要とするからです。

\begin{codescreen}{}
[1,] // ,で終わっている
\end{codescreen}
\clearpage
\subsection{value}

\texttt{value}はJSONの値を表現する規則です。

\begin{codescreen}{}
value = true | false | null | number | string | object | array;
\end{codescreen}

これは、JSONの値は、

\begin{itemize}
\item
  真（\texttt{true}）
\item
  偽（\texttt{false}）
\item
  ヌル（\texttt{null}）
\item
  数値（\texttt{number}）
\item
  文字列（\texttt{string}）
\item
  オブジェクト（\texttt{object}）
\item
  配列（\texttt{array}）
\end{itemize}

\noindent のいずれかでなければいけないことを示しています。JSONを普段使っている皆さんにはおなじみでしょう。

\subsection{true}

\texttt{true}は、真を表すリテラルを表す規則です。

\begin{codescreen}{}
true = "true" ws;
\end{codescreen}

文字列 \texttt{true} が真を表すということでそのままですね。

\subsection{false}

\texttt{false}は、偽を表すリテラルを表す規則です。構造的には、\texttt{true}と同じです。

\begin{codescreen}{}
false = "false" ws;
\end{codescreen}

\subsection{null}

\texttt{null}は、ヌルリテラルを表す規則です。構造的には、\texttt{true}や\texttt{false}と同じです。

\begin{codescreen}{}
null = "null" ws;
\end{codescreen}

JSON仕様は\texttt{null}の意味付けを規定しません。多くの実装ではホスト言語の\texttt{null}相当の値に対応づけますが、\texttt{null}を表現するためのリテラルがない言語では、\texttt{Option}や\texttt{Maybe}といったデータ型に対応づけることもあります。

\subsection{number}

\texttt{number}は、数値リテラルを表す規則です。

\begin{codescreen}{}
number = INT ws;
\end{codescreen}

整数（\texttt{INT}）に続いて、\texttt{ws}が来るのが\texttt{number}であるということを表現しています。簡単にするため、本章では数値を整数に限定し、小数部や指数部は扱いません（実装拡張の話題は章末の演習を参照）。

\subsection{string}

\texttt{string}は文字列リテラルを表す規則です。

\begin{codescreen}{}
string = "\"" {CHAR} "\"" ws;
\end{codescreen}

\texttt{"}で始まって、\texttt{CHAR}で定義される文字が0個以上続いて、
\texttt{"}
で終わります。

\subsection{JSONのBNFまとめ}

JSONのBNFは、非常に少数の規則だけで表現することができます。読者の中には、あまりにも簡潔過ぎて驚かれた方もいるのではないでしょうか。しかし、これだけ単純であるにも関わらず、JSONのBNFは\textgt{再帰的に定義されている}ため、非常に複雑な構造も表現できます。たとえば、

\begin{itemize}
\item
  一要素の配列があり、その要素はオブジェクトであり、キー\texttt{"a"}に対応する要素の中に配列があって、その配列は空配列である
\end{itemize}

\noindent といったことも、JSONでは以下のように表現できます。

\begin{codescreen}{}
[{"a":[]}]
\end{codescreen}


再帰的な規則は構文解析において非常に重要な要素なので、これから本書を読み進める上でも念頭に置いてください。

\section{JSONの抽象構文木}
\index{JSON!抽象構文木}
\index{ちゅうしょうこうぶんぎ@抽象構文木}
\index{AST}

JSONの定義と、文法について見てきました。構文解析器を実装する前に、まずJSONの抽象構文木（AST:
Abstract Syntax Tree）をJavaでどのように表現するかを定義しましょう。

抽象構文木は、第1章でも説明したとおり、プログラムの構造を表現するためのデータ構造です。重要な点は、抽象構文木では元の文字列に含まれていた「構文上のノイズ」が取り除かれることです。たとえば

\begin{itemize}
\item
  空白文字（\texttt{ws}）は抽象構文木には含まれません
\item
  カンマ（\texttt{,}）やコロン（\texttt{:}）などの区切り文字も、構造として表現されるため個別のノードにはなりません
\item
  括弧（\texttt{\{\}}、\texttt{{[}{]}}）も同様に、オブジェクトや配列という構造で表現されます
\end{itemize}

以下の\texttt{JsonAst}の定義を見ると、JSONの各要素をJavaのレコードクラス（record）やインタフェースとして表現していることがわかります。これは、JSONという「文字列」を、Javaのオブジェクトという「構造化されたデータ」に変換するための定義です。

\begin{codescreen}{}
public interface JsonAst {
    // value
    sealed interface JsonValue permits 
        JsonNull, JsonTrue, JsonFalse, 
        JsonNumber, JsonString, 
        JsonObject, JsonArray {}

    // NULL
    record JsonNull() implements JsonValue {
        @Override
        public String toString() {
            return "null";
        }
    }

    // TRUE
    record JsonTrue() implements JsonValue {
        @Override
        public String toString() {
            return "true";
        }
    }

    // FALSE
    record JsonFalse() implements JsonValue {
        @Override
        public String toString() {
            return "false";
        }
    }

    // NUMBER
    record JsonNumber(double value) implements JsonValue {
        @Override
        public String toString() {
            return "JsonNumber(" + value + ')';
        }
    }

    // STRING
    record JsonString(String value) implements JsonValue {
        @Override
        public String toString() {
            return "JsonString(\"" + value + "\")";
        }
    }

    // object
    record JsonObject(List<Pair<JsonString, JsonValue>> properties) 
        implements JsonValue {
        @Override
        public String toString() {
            return "JsonObject{" + properties + '}';
        }
    }

    // array
    record JsonArray(List<JsonValue> elements) 
        implements JsonValue {
        @Override
        public String toString() {
            return "JsonArray[" + elements + ']';
        }
    }
}
\end{codescreen}


各クラスがBNFの左辺に対応しているのがわかるでしょうか。ただし、前述のとおり、\texttt{ws}や区切り文字（\texttt{,}、\texttt{:}）などは抽象構文木には現れません。これらは構文解析の過程で消費されますが、最終的なデータ構造には含まれないのです。

\section{JSONの構文解析器}
\index{JSON!構文解析器}
\index{こうぶんかいせきき@構文解析器}

この節では、BNFを元に、JSONのデータを\textgt{構文解析}するプログラムを実装していきます。以下のようなインタフェース\texttt{JsonParser}インタフェースを実装したクラスを「JSONの構文解析器」と考えることにします。

\begin{codescreen}{}
package parser;
interface JsonParser {
    public ParseResult<JsonAst.JsonValue> parse(String input);
}
\end{codescreen}

クラス\texttt{ParseResult\textless{}T\textgreater{}}は以下のようなジェネリックなクラスになっています。\texttt{value}は解析結果の値です。これは任意の型をとり得るので、\texttt{T}としています。また、\texttt{input}は「構文解析の対象となる文字列」を表します。

\begin{codescreen}{}
public record ParseResult<T>(T value, String input) {}
\end{codescreen}

インタフェース\texttt{JsonParser}は\texttt{parse()}メソッドだけを持ちます。\texttt{parse()}メソッドは、文字列\texttt{input}を受け取り、\texttt{ParseResult\textless{}JsonAst.JsonValue\textgreater{}}型を返します。

\subsection{構文解析の戦略}

実装に入る前に、これから作るJSONパーサーの解析戦略について説明しておきます。この章の実装は、PEG（Parsing
Expression Grammar）に基づくトップダウン解析を採用します。

\begin{enumerate}
\item
  トップダウン解析:
  BNFの規則を上から順番に適用していく方式を採用します。たとえば、\texttt{value}規則から始めて、それぞれの選択肢を試していきます。
\item
  バックトラック:
  解析に失敗した場合、前の状態に戻って別の選択肢を試します。これにより、複数の可能性がある場合でも正しい解析結果を見つけられます。
\item
  順序付き選択（Ordered Choice）: BNFの \texttt{\textbar{}}
  で区切られた選択肢を、左から右へ順番に試します。たとえば、\texttt{value\ =\ true\ \textbar{}\ false\ \textbar{}\ null\ \textbar{}\ number\ \textbar{}\ string\ \textbar{}\ object\ \textbar{}\ array}
  の場合、まず\texttt{true}を試し、失敗したら\texttt{false}を試す、という具合です。この順序は重要で、より具体的なパターンを先に配置することで、正しい解析を保証します。実装の\texttt{parseValue()}メソッドでも、BNFの定義順序と同じ順番で各選択肢を試すようにしています。
\item
  例外による失敗の表現:
  構文解析の失敗は、\texttt{ParseException}という例外で表現します。これにより、深くネストした解析処理から簡潔に失敗を伝播させられます。
\end{enumerate}

本章の実装はPEGの順序付き選択とバックトラックという特徴を持っています。そのため、クラス名は\texttt{PegJsonParser}としています（PEGの詳細は第4章で解説します）。

\subsection{構文解析器の全体像}

それでは、JSONの構文解析器の実装を見ていきましょう。構文解析を行う各メソッドについては省略して、後ほど解説することにします。以下が\texttt{PegJsonParser}クラスの全体像です。

\begin{codescreen}{}
package parser;

import java.util.ArrayList;
import java.util.List;

public class PegJsonParser implements JsonParser {
    private int cursor;
    private String input;

    private static class ParseException extends RuntimeException {
        public ParseException(String message) {
            super(message);
        }
    }

    public ParseResult<JsonAst.JsonValue> parse(String input) {
        this.input = input;
        this.cursor = 0;
        var value = parseValue();
        return new ParseResult<>(value, input.substring(this.cursor));
    }

    /**
     * 指定された文字列リテラルが現在のカーソル位置にあることを確認し、
     * マッチした場合はカーソルを進めます。
     * @param literal 認識すべき文字列リテラル
     * @throws ParseException リテラルがマッチしない場合
     */
    private void recognize(String literal) {
        if(input.substring(cursor).startsWith(literal)) {
            cursor += literal.length();
        } else {
            String substring = input.substring(cursor);
            int endIndex = cursor + 
              (literal.length() > substring.length() ? 
                substring.length() : literal.length());
            throw new ParseException(
                "expected: " + literal + 
                ", actual: " + input.substring(cursor, endIndex)
            );
        }
    }


    /**
     * 空白文字（スペース、タブ、改行、キャリッジリターン）を
     * スキップし、カーソルを次の非空白文字まで進めます。
     */
    private void skipWhitespace() {
        OUTER:
        while(cursor < input.length()) {
            char currentCharacter = input.charAt(cursor);
            switch (currentCharacter) {
                // 空白文字
                case ' ':       // Space
                case '\t':      // Horizontal Tab
                case '\n':      // Line Feed
                case '\r':      // Carriage Return
                    cursor++;
                    continue OUTER;
                default:
                    break OUTER;
            }
        }
    }

    /**
     * JSONの値（value）を解析します。
     * BNF: value = true | false | null | number | string | object | array
     * @return 解析されたJSON値
     * @throws ParseException 有効なJSON値が見つからない場合
     */
    private JsonAst.JsonValue parseValue() {
        // 後で解説
    }

  
    /**
     * true値を解析します。
     * BNF: true = "true" ws
     * @return JSONのtrue値
     * @throws ParseException "true"がマッチしない場合
     */
    private JsonAst.JsonTrue parseTrue() {
        // 後で解説
    }

    /** BNF: false = "false" ws */
    private JsonAst.JsonFalse parseFalse() {
        // 後で解説
    }

    /** BNF: null = "null" ws */
    private JsonAst.JsonNull parseNull() {
        // 後で解説
    }

    /** BNF: LBRACE = '{' ws */
    private void parseLBrace() {
        // 後で解説
    }

    /** BNF: RBRACE = '}' ws */
    private void parseRBrace() {
        // 後で解説
    }

    /** BNF: LBRACKET = '[' ws */
    private void parseLBracket() {
        // 後で解説
    }

    /** BNF: RBRACKET = ']' ws */
    private void parseRBracket() {
        // 後で解説
    }

    /** BNF: COMMA = ',' ws */
    private void parseComma() {
        // 後で解説
    }

    /** BNF: COLON = ':' ws */
    private void parseColon() {
        // 後で解説
    }

    /** BNF: string = ('""' | '"' {CHAR} '"') ws */
    private JsonAst.JsonString parseString() {
        // 後で解説
    }

    /** BNF: number = INT ws */
    private JsonAst.JsonNumber parseNumber() {
        // 後で解説
    }

    /** BNF: pair = string COLON value */
    private Pair<JsonAst.JsonString, JsonAst.JsonValue> parsePair() {
        // 後で解説
    }

    /** BNF: object = LBRACE RBRACE | LBRACE pair {COMMA pair} RBRACE */
    private JsonAst.JsonObject parseObject() {
        // 後で解説
    }

    /** BNF: array = LBRACKET RBRACKET | LBRACKET value {COMMA value} RBRACKET */
    public JsonAst.JsonArray parseArray() {
        // 後で解説
    }
}
\end{codescreen}

このクラス\texttt{PegJsonParser}で重要なことは、クラスがフィールドとして以下を保持していることです。

\begin{codescreen}{}
public class PegJsonParser implements JsonParser {
    private int cursor;
    private String input;
    // ...
}
\end{codescreen}

構文解析器の実装スタイルには、大きく分けて「関数型（不変データ・純粋関数を重視）」と「手続き型（内部状態と更新を用いる）」があります。どちらも正しく実装すれば、同じ入力に対して同じ解析結果を返しますが、表現と状態管理の方法が異なります。本書では説明のしやすさから手続き型の実装を採用します。

また、\texttt{skipWhitespace()}メソッドと\texttt{recognize()}メソッドを定義して、空白文字の読み飛ばしと、特定の文字列の認識を行います。さらに、構文解析中にエラーが発生した場合は\texttt{ParseException}を投げることで、どこでどのような問題が発生したかを呼び出し元に伝えます。

\texttt{cursor}フィールドは現在の読み取り位置を、\texttt{input}フィールドは解析対象の文字列を保持します。エラーが発生した場合は\texttt{ParseException}を投げることで、どこでどのような問題が発生したかを呼び出し元に伝えます。

次からはいよいよ構文解析の各メソッドを見ていきます。

\subsection{valueの構文解析メソッド}

まずは、\texttt{parseValue()}メソッドから始めましょう。

\begin{codescreen}{}
    private JsonAst.JsonValue parseValue() {
        int backup = cursor;
        try {
            return parseTrue();
        } catch (ParseException e) {
            cursor = backup;
        }

        try {
            return parseFalse();
        } catch (ParseException e) {
            cursor = backup;
        }
        // parseNull(), parseNumber(), parseString(), 
        // parseObject(), parseArray() と続く...
    }
\end{codescreen}

\texttt{parseValue()}の実装は、BNFの\texttt{value\ =\ true\ \textbar{}\ false\ \textbar{}\ null\ \textbar{}\ number\ \textbar{}\ string\ \textbar{}\ object\ \textbar{}\ array}という定義に対応しています。

\texttt{\textbar{}}は「または」を意味するので、「valueは、trueまたはfalseまたはnullまたは……」という意味になります。このコードでは、それぞれの可能性を順番に試していきます。

\begin{enumerate}
\item
  まず\texttt{parseTrue()}を試し、失敗したらカーソルを元に戻す
\item
  次に\texttt{parseFalse()}を試し、失敗したらカーソルを元に戻す
\item
  以下同様に続く……
\end{enumerate}

このように、失敗したら元の位置に戻って別の可能性を試すことを「バックトラック」と呼びます。

ちなみに、この順番は重要です。\texttt{value}の例では右辺がそれぞれ排他的なので問題ありませんが、順番を変えると結果が変わってしまうことがあります。

\subsection{trueの構文解析メソッド}

\texttt{true}の構文解析は、次のような \texttt{parseTrue()}
メソッドとして定義します。

\begin{codescreen}{}
private JsonAst.JsonTrue parseTrue() {
    recognize("true");
    skipWhitespace();
    return new JsonAst.JsonTrue();
}
\end{codescreen}

このメソッドで行っていることを見ていきましょう。このメソッドでは、入力である\texttt{input}の現在位置が\texttt{"true"}という文字列で始まっているかをチェックします。もしそうなら、\textbf{\sffamily JSON}\textgt{の}\textbf{\sffamily true}を表す\texttt{JsonAst.JsonTrue}のインスタンスを返します。もし、先頭が\texttt{"true"}でなければ、構文解析は失敗なので例外を発生させますが、これは\texttt{recognize()}メソッドの中で行われています。\texttt{recognize()}の内部では、入力の現在位置と与えられた文字列を照合して、マッチしない場合例外を投げます。

次に、\texttt{skipWhitespace()}メソッドを呼び出して、「空白の読み飛ばし」を行っています。

\subsection{falseの構文解析メソッド}

\texttt{false}の構文解析は、次のような \texttt{parseFalse()}
メソッドとして定義します。

\begin{codescreen}{}
private JsonAst.JsonFalse parseFalse() {
    recognize("false");
    skipWhitespace();
    return new JsonAst.JsonFalse();
}
\end{codescreen}

これも、\texttt{parseTrue()}とほぼ同じですので、特に説明の必要はないでしょう。

\subsection{nullの構文解析メソッド}

\texttt{null}の構文解析は、次のような \texttt{parseNull()}
メソッドとして定義します。

\begin{codescreen}{}
private JsonAst.JsonNull parseNull() {
    recognize("null");
    skipWhitespace();
    return new JsonAst.JsonNull();
}
\end{codescreen}

これも、\texttt{parseTrue()}や\texttt{parseFalse()}とほぼ同じです。固定の文字列を解析するという点で、これらのメソッドは共通の処理になっています。

\subsection{数値の構文解析メソッド}

数値の構文解析は、次のような \texttt{parseNumber()}
メソッドとして定義します。

\begin{codescreen}{}
    private JsonAst.JsonNumber parseNumber() {
        // 本書の実装では、JSONの数値型を整数のみに限定して扱います。
        int start = cursor;
        char ch = 0;
        // オプショナルなマイナス記号
        if (cursor < input.length() && input.charAt(cursor) == '-') {
            cursor++;
        }

        // 数字の連続 (0, または 1-9 で始まり数字が続く)
        int digitsStart = cursor;
        if (cursor < input.length()) {
            ch = input.charAt(cursor);
            if (ch == '0') {
                cursor++;
            } else if ('1' <= ch && ch <= '9') {
                cursor++;
                while (cursor < input.length()) {
                    ch = input.charAt(cursor);
                    if (!('0' <= ch && ch <= '9')) break;
                    cursor++;
                }
            } else {
                // 数字で始まらない場合はエラー
                throw new ParseException(
                    "expected: digit actual: " + 
                    (cursor < input.length() ? input.charAt(cursor) : "EOF")
                );
            }
        } else {
             throw new ParseException("expected: digit actual: EOF");
        }

        if (digitsStart == cursor) {
             // 数字が1つも読まれなかった場合（"-"のみなど）
             throw new ParseException(
                "expected: digit after '-' actual: " + 
                (cursor < input.length() ? input.charAt(cursor) : "EOF")
             );
        }
        
        String numberStr = input.substring(start, cursor);
        try {
            double value = Double.parseDouble(numberStr);
            skipWhitespace();
            return new JsonAst.JsonNumber(value);
        } catch (NumberFormatException e) {
            throw new ParseException("invalid number format: " + numberStr);
        }
    }
\end{codescreen}

\texttt{parseNumber()}
メソッド（PegJsonParser内）では、入力文字列から数値部分を読み取り、\texttt{Double.parseDouble}
を用いて数値に変換しています。この実装は、ECMA-404で定義されるJSONの数値型の完全な仕様（小数部、指数部を含む）には対応しておらず、整数のみを扱えるように単純化されています。

\subsection{文字列の構文解析メソッド}

文字列の構文解析は、次のような \texttt{parseString()}
メソッドとして定義します。

\begin{codescreen}{}
    private JsonAst.JsonString parseString() {
        if(cursor >= input.length()) {
            throw new ParseException("expected: \"" + " actual: EOF");
        }
        char ch = input.charAt(cursor);
        if(ch != '"') {
            throw new ParseException("expected: \"" + "actual: " + ch);
        }
        cursor++;
        var builder = new StringBuilder();
        var closed = false;
        OUTER:
        while(cursor < input.length()) {
            ch = input.charAt(cursor);
            switch(ch) {
                case '\\':
                    throw new ParseException(
                        "escape sequences are not supported in this parser"
                    );
                case '"':
                    cursor++;
                    closed = true;
                    break OUTER;
                default:
                    builder.append(ch);
                    cursor++;
                    break;
            }
        }

        if(!closed) {
            throw new ParseException("expected: " + "\"" + " actual: EOF");
        } else {
            skipWhitespace();
            return new JsonAst.JsonString(builder.toString());
        }
    }
\end{codescreen}

\texttt{while}文の中が若干複雑になっていますが、1つ1つ見ていきます。

まず、最初の部分では、

\begin{codescreen}{}
        if(cursor >= input.length()) {
            throw new ParseException("expected: \"" + " actual: EOF");
        }
        char ch = input.charAt(cursor);
        if(ch != '"') {
            throw new ParseException("expected: \"" + "actual: " + ch);
        }
        cursor++;
\end{codescreen}

\begin{itemize}
\item
  入力が終端に達していないこと
\item
  入力の最初が\texttt{"}であること
\end{itemize}

\noindent をチェックしています。文字列は当然ながら、ダブルクォートで始まりますし、入力の終端にも達していないはずですから、それらの条件が満たされなければ例外が投げられるわけです。

\texttt{while}文の中では、各文字を読み込んで文字列を構築していきます。switch文の中で最も重要なのは\texttt{default}ケースで、ここで通常の文字を\texttt{StringBuilder}に追加して文字列を構築しています。ダブルクォート（\texttt{"}）が現れたら文字列の終端として処理し（このとき、閉じ引用符を読んだことを変数\texttt{closed}に記録します）、バックスラッシュ（\texttt{\textbackslash{}}）が現れた場合はエラーとして処理します。実際のJSONパーサーではエスケープシーケンスの処理が必要ですが、構文解析の本質を理解する上では必須ではないため、本書では省略しています。

\texttt{while}文が終わったあとで、

\begin{codescreen}{}
        if(!closed) {
            throw new ParseException("expected: " + "\"" + " actual: EOF");
        } else {
            skipWhitespace();
            return new JsonAst.JsonString(builder.toString());
        }
\end{codescreen}

\noindent というチェックを入れることによって、閉じのダブルクォートを読んで文字列が正しく終端したことを確認した後、空白を読み飛ばしています。閉じ引用符を読む前に入力が尽きた場合は\texttt{closed}が\texttt{false}のままなので、例外が投げられます。

\subsection{配列の構文解析メソッド}

配列の構文解析は、次のような \texttt{parseArray()}
メソッドとして定義します。

\begin{codescreen}{}
    public JsonAst.JsonArray parseArray() {
        int backup = cursor;
        try {
            parseLBracket();
            parseRBracket();
            return new JsonAst.JsonArray(new ArrayList<>());
        } catch (ParseException e) {
            cursor = backup;
        }

        parseLBracket();
        List<JsonAst.JsonValue> values = new ArrayList<>();
        var value = parseValue();
        values.add(value);
        try {
            while (true) {
                backup = cursor;
                parseComma();
                value = parseValue();
                values.add(value);
            }
        } catch (ParseException e) {
            cursor = backup;
            parseRBracket();
            return new JsonAst.JsonArray(values);
        }
    }
\end{codescreen}

この\texttt{parseArray()}は多少複雑になります。まず、先頭に\texttt{"{[}"}が来るかチェックする必要があります。これをコードにすると、以下のようになります。

\begin{codescreen}{}
parseLBracket();
\end{codescreen}

\texttt{parseLBracket()}は以下のように定義されています。

\begin{codescreen}{}
 private void parseLBracket() {
    recognize("[");
    skipWhitespace();
}
\end{codescreen}

\texttt{recognize()}で、現在の入力位置が\texttt{{[}}と一致しているかチェックをした後、空白を読み飛ばしています。

\texttt{{[}}の次には任意の\texttt{JsonValue}または\texttt{"{]}"}が来る可能性があります。このとき、まず最初に、\texttt{{]}}が来ると\textgt{仮定}するのがポイントです。

\begin{codescreen}{}
        int backup = cursor;
        try {
            parseLBracket();
            parseRBracket();
            return new JsonAst.JsonArray(new ArrayList<>());
        } catch (ParseException e) {
            cursor = backup;
        }
\end{codescreen}

仮定が成り立たなかった場合、\texttt{ParseException}がthrowされるはずですから、それをcatchして、バックアップした位置に巻き戻します。

\texttt{{]}}が来るという仮定が成り立たなかった場合、再び最初に\texttt{{[}}が出現して、その次に来るのは任意の\texttt{JsonValue}ですから、以下のようなコードになります。

\begin{codescreen}{}
        parseLBracket();
        List<JsonAst.JsonValue> values = new ArrayList<>();
        var value = parseValue();
        values.add(value);
\end{codescreen}

ローカル変数\texttt{values}は、配列の要素を格納するためのものです。

配列の中で、最初の要素が読み込まれた後、次に来るのは、\texttt{,}か\texttt{{]}}のどちらかですが、ひとまず、\texttt{,}が来ると仮定して\texttt{while}ループで

\begin{codescreen}{}
parseComma();
value = parseValue();
values.add(value);
\end{codescreen}

\noindent を繰り返します。この繰り返しは、1回ごとに必ず入力を1以上進めます。失敗したときは、テキストが正しいJSONなら、\texttt{{]}}が来るはずなので、

\begin{codescreen}{}
parseRBracket();
return new JsonAst.JsonArray(values);
\end{codescreen}

\noindent とします。もし、テキストが正しいJSONでない場合、\texttt{parseRBracket()}から例外が投げられるはずですが、その例外は\textgt{より上位の層が適切にリカバーしてくれると期待して}放置します。JSONのような再帰的な構造を解析するとき、このような、「自分の呼び出し元が適切にやってくれるはず」（なぜなら、自分はその呼び出し元で適切にcatchしているのだから）という考え方が重要になります。

多少複雑になりましたが、\texttt{parseArray()}の定義が、EBNFにおける表記

\begin{codescreen}{}
array = LBRACKET RBRACKET | LBRACKET value {COMMA value} RBRACKET ;
\end{codescreen}

\noindent に対応していることがわかるでしょうか
。読み方のポイントは、\texttt{\textbar{}}の後を、例外をキャッチした後の処理ととらえることです。

\subsection{オブジェクトの構文解析メソッド}

オブジェクトの構文解析は、次のような \texttt{parseObject()}
メソッドとして定義します。

\begin{codescreen}{}
    /** BNF: object = LBRACE RBRACE | LBRACE pair {COMMA pair} RBRACE */
    private JsonAst.JsonObject parseObject() {
        int backup = cursor;
        try {
            parseLBrace();
            parseRBrace();
            return new JsonAst.JsonObject(new ArrayList<>());
        } catch (ParseException e) {
            cursor = backup;
        }

        parseLBrace();
        List<Pair<JsonAst.JsonString, JsonAst.JsonValue>> members =
            new ArrayList<>();
        var member = parsePair();
        members.add(member);
        try {
            while (true) {
                backup = cursor;
                parseComma();
                member = parsePair();
                members.add(member);
            }
        } catch (ParseException e) {
            cursor = backup;
            parseRBrace();
            return new JsonAst.JsonObject(members);
        }
    }
\end{codescreen}

この定義を見て、ひょっとしたら、

「あれ？これ、\texttt{parseArray()}とほとんど同じでは」

\noindent と気づかれた読者の方がいるかも知れません。実際、\texttt{parseObject()}がやっていることは\texttt{parseArray()}と類似しています。

最初に、

\begin{codescreen}{}
parseLBrace();
parseRBrace();
return new JsonAst.JsonObject(new ArrayList<>());
\end{codescreen}

\noindent としている箇所は、\texttt{\{\}}という形の空オブジェクトを読み取ろうとしていますが、これは、空配列\texttt{{[}{]}}を読み取るコードとほぼ同じです。

続くコードも、対応する記号が\texttt{\{\}}か\texttt{{[}{]}}の違いこそあるものの、基本的に同じです。唯一の違いは、オブジェクトの各要素は、\texttt{name:value}というペアなため、\texttt{parseValue()}の代わりに\texttt{parsePair()}を呼び出しているところくらいです。

そして、\texttt{parsePair()}は以下のように定義されています。

\begin{codescreen}{}
    private Pair<JsonAst.JsonString, JsonAst.JsonValue> parsePair() {
        var key = parseString();
        parseColon();
        var value = parseValue();
        return new Pair<>(key, value);
    }
\end{codescreen}

これはEBNFにおける以下の定義にそのまま対応しているのがわかるでしょう。

\begin{codescreen}{}
pair = string COLON value;
\end{codescreen}

\subsection{構文解析における再帰}

配列やオブジェクトの構文解析メソッドを見るとわかりますが、

\begin{itemize}
\item
  \texttt{parseArray()\ →\ parseValue()\ →\ parseArray()}
\item
  \texttt{parseArray()\ →\ parseValue()\ →\ parseObject()}
\item
  \texttt{parseObject()\ →\ parseValue()\ →\ parseObject()}
\item
  \texttt{parseObject()\ →\ parseValue()\ →\ parseArray()}
\end{itemize}

\noindent のような再帰呼び出しが起こり得ることがわかります。このような再帰呼び出しでは、各ステップで必ず構文解析が1文字以上進むため、JSONがどれだけ深くなっても（スタックが溢れない限り）うまく構文解析ができるのです。

\subsection{構文解析とPEG}

このようにしてJSONの構文解析器を実装できました。

実は、ここで作った構文解析器には、以下のような特徴があります。

\begin{enumerate}
\item
  \textgt{バックトラック}：失敗したら元の位置に戻って別の可能性を試す
\item
  \textgt{順序付き選択}：\texttt{\textbar{}}で区切られた選択肢を左から順番に試す
\item
  \textgt{文字列を直接解析}：特別な前処理なしに、入力文字列をそのまま解析する
\end{enumerate}

このような構文解析の手法を\textbf{\sffamily PEG（Parsing Expression
Grammar}\textgt{、解析表現文法）}と呼びます。PEGは2004年に提案された比較的新しい手法で、プログラミング言語のような「曖昧さがない」言語の解析に適しています。最近ではPython（バージョン3.9以降）もPEGベースの構文解析器を使っています\footnote{PythonにおけるPEGパーサー採用の経緯と理由は
  \href{https://peps.python.org/pep-0617/}{PEP 617 -- New PEG parser for
  CPython} で詳しく説明されています。}。

厳密にはPEG自体はBNFと異なる文法の定義方法（形式文法）ですが、PEG自体が「どのように解析されるか」を定義するため、ここではPEG自体を構文解析の手法として扱います。この点については第4章で詳しく解説します。

PEGは直感的でシンプルなので、最初に学ぶのに適しています。ただし、従来から使われている別の構文解析手法も重要です。次の節では、その伝統的な手法について解説します。

\section{古典的な構文解析器}
\index{こうぶんかいせきき@構文解析器!古典的}
\index{じくかいせきき@字句解析器}
\index{じくかいせき@字句解析}
\index{とーくん@トークン}

前節では、PEGという手法を使って構文解析器を作りました。しかし、LL/LR系（例：LL(1)、LR(1)）に代表される伝統的な構文解析の手法では、少し違ったアプローチをとります。

伝統的な手法では、構文解析を以下の2つのステップに分けます。

\begin{enumerate}
\item
  \textgt{字句解析（トークナイザー）}：文字列を「トークン」という単位に分割する
\item
  \textgt{構文解析（パーサー）}：トークンの列から構造を組み立てる
\end{enumerate}

たとえば、以下の英文があったとします。

\begin{codescreen}{}
We are parsers.
\end{codescreen}

我々は構文解析器であるというジョーク的な文ですが、それはさておき、この文は

\begin{codescreen}{}
[We, are, parsers]
\end{codescreen}

\noindent という3つのトークン（単語）に分解すると考えるのが字句解析の発想法です。

古典的な構文解析の世界では、字句解析が必須とされていましたが、それは後の章で説明される構文解析アルゴリズムの都合に加えて、空白のスキップという処理を字句解析で行えるからでもあります。

前節で出てきたJSONの構文解析器では\texttt{skipWhitespace()}の呼び出しが頻出していましたが、字句解析器を使う場合、空白を読み飛ばす処理を先に行うことで、構文解析器では空白の読み飛ばしという作業をしなくてよくなります。

この点はトレードオフがあって、空白に関する規則がある言語では、字句解析という前処理を安易に入れると逆に扱いづらくなることすらあります。ともあれ、字句解析という前処理を通すことには一定のメリットがあるのは確かです。

以下では字句解析器を使った構文解析器の全体像を示します。ここでは、コアとなるアイデアに絞って説明します。完全なソースコードは本書のサポートサイト（\url{https://github.com/kahei/parser_book}）を参照してください。

\subsection{JSONの字句解析器}

JSONの字句解析器（トークナイザー）の基本構造は次のようになります。

\begin{codescreen}{}
public class SimpleJsonTokenizer implements JsonTokenizer {
    private final String input;
    private int index;

    // すべての入力を一度にトークン化
    public List<Token> tokenizeAll() {
        List<Token> tokens = new ArrayList<>();
        while (moveNext()) {
            tokens.add(current());
        }
        tokens.add(new Token(Token.Type.EOF, null));
        return tokens;
    }

    // 次のトークンを読み取る
    private boolean moveNext() {
        skipWhitespace();
        
        if (index >= input.length()) {
            return false;
        }
        
        char ch = input.charAt(index);
        switch (ch) {
            case '"':
                return tokenizeStringLiteral();
            case '{':
                accept("{", Token.Type.LBRACE, "{");
                return true;
            case '[':
                accept("[", Token.Type.LBRACKET, "[");
                return true;
            // 他のトークンも同様に処理...
        }
    }
    
    // 文字列トークンの読み取り（エスケープシーケンス非対応）
    private boolean tokenizeStringLiteral() {
        if(input.charAt(index) != '"') return false;
        index++;
        var builder = new StringBuilder();
        while(index < input.length()) {
            char ch = input.charAt(index);
            if(ch == '"') {
                fetched = new Token(Token.Type.STRING, builder.toString());
                index++;
                return true;
            }
            builder.append(ch);
            index++;
        }
        return false;
    }
}
\end{codescreen}

このコードは\texttt{SimpleJsonTokenizer}クラスの基本構造を示しています。主な特徴は以下のとおりです。

\begin{itemize}
\item
  \texttt{tokenizeAll()}：

  \begin{itemize}
  \item
    入力文字列全体を一度に処理してトークン列を生成
  \item
    \texttt{moveNext()}を繰り返し呼び出してトークンを収集
  \item
    最後に\texttt{EOF}（End Of File）トークンを追加
  \end{itemize}
\item
  \texttt{moveNext()}：

  \begin{itemize}
  \item
    まず空白文字をスキップ
  \item
    現在の文字を見て、トークンの種類を判定
  \item
    switch文で各文字に応じた処理を実行

    \begin{itemize}
    \item
      \texttt{"} → 文字列リテラルの処理
    \item
      \texttt{\{}、\texttt{{[}} → 括弧トークンの生成
    \item
      その他の文字も同様にパターンマッチング
    \end{itemize}
  \end{itemize}
\item
  \texttt{tokenizeStringLiteral()}メソッド：

  \begin{itemize}
  \item
    ダブルクォートで囲まれた文字列を読み取る
  \item
    閉じクォートが見つかるまで文字を収集
  \item
    エスケープシーケンス非対応
  \end{itemize}
\end{itemize}

この設計により、字句解析の責任が明確になり、構文解析器は純粋にトークンの構造解析に専念できます。

\subsection{JSONの構文解析器}

この字句解析器を使った構文解析器の基本構造を示します。トークナイザーが一度にすべてのトークンを生成し、パーサーはそのトークン列を処理します。

\begin{codescreen}{}
public class SimpleJsonParser implements JsonParser {
    private List<Token> tokens;
    private int index;

    public ParseResult<JsonAst.JsonValue> parse(String input) {
        // 入力を一度に完全にトークン化
        var tokenizer = new SimpleJsonTokenizer(input);
        this.tokens = tokenizer.tokenizeAll();
        this.index = 0;
        var value = parseValue();
        return new ParseResult<>(value, "");
    }
    
    private Token current() {
        if (index < tokens.size()) {
            return tokens.get(index);
        }
        return new Token(Token.Type.EOF, null);
    }
    
    private boolean moveNext() {
        if (index < tokens.size() - 1) {
            index++;
            return true;
        }
        return false;
    }

    // トークン列から抽象構文木を構築
    private JsonAst.JsonValue parseValue() {
        // 中身は後で解説
    }

    // オブジェクトの解析（トークンベース）
    private JsonAst.JsonObject parseObject() {
        // 中身は後で解説
    }
    
    // 配列の解析（トークンベース）
    private JsonAst.JsonArray parseArray() {
        // 中身は後で解説
    }
    // 他のメソッドも同様に実装
}
\end{codescreen}

このアプローチの重要な点は、\textgt{字句解析と構文解析の責任が明確に分離されている}ことです。\texttt{tokenizeAll()}メソッドが最初に入力全体をトークン列に変換し、パーサーはそのトークン列を処理します。これにより

\begin{enumerate}
\item
  関心の分離:
  字句解析器はトークンの認識に、構文解析器は構造の構築に専念できます
\item
  デバッグの容易さ:
  トークン列を事前に確認できるため、問題の切り分けが簡単です
\item
  実装の簡潔さ:
  各コンポーネントが単一の責任を持つため、コードがシンプルになります
\end{enumerate}

構文解析器の\texttt{parseXXX()}メソッドを見ると、文字列の代わりにトークン列を処理していることがわかります。また、この構文解析器には空白の読み飛ばしに関する処理が入っていません。これは、字句解析器が空白を処理済みだからです。

PEG版と異なり、途中で失敗したら後戻り（バックトラック）するという処理も存在しません。トークン列が事前に確定しているため、より決定的な解析が可能になります。

以下では構文解析のための各メソッドの詳細を説明します。残りのコードは、本書のサポートサイト（\url{https://github.com/kahei/parser_book}）に掲載しています。

\subsection{valueの構文解析メソッド（トークンベース）}

\texttt{parseValue()}メソッドは、JSONの値を解析するためのメソッドです。これは、BNFで定義された\texttt{value\ =\ true\ \textbar{}\ false\ \textbar{}\ null\ \textbar{}\ number\ \textbar{}\ string\ \textbar{}\ object\ \textbar{}\ array}に対応しています。実装は以下のようになります。

\begin{codescreen}{}
    private JsonAst.JsonValue parseValue() {
        var token = current();
        switch(token.type) {
            case TRUE:
                return parseTrue();
            case FALSE:
                return parseFalse();
            case NULL:
                return parseNull();
            case INTEGER:
                return parseNumber();
            case STRING:
                return parseString();
            case LBRACE:
                return parseObject();
            case LBRACKET:
                return parseArray();
        }
        throw new RuntimeException("cannot reach here");
    }
\end{codescreen}

\texttt{current()}メソッドは、現在のトークンを取得するためのメソッドで、トークン列から現在の位置のトークンを返します。\texttt{switch}文では、先頭のトークンの種類に応じて適切な解析メソッドを呼び出しています。バックトラックが必要ないため、各トークンの種類に応じて直接対応するメソッドを呼び出す形になっています。

\subsection{parseObject}

\texttt{parseObject()}メソッドは、規則\texttt{object}に対応するメソッドで、JSONのオブジェクトリテラルに対応するものを解析するメソッドでもあります。実装を示すと以下のようになります。

\begin{codescreen}{}
    private JsonAst.JsonObject parseObject() {
        if(current().type != Token.Type.LBRACE) {
            throw new ParseException(
                "expected `{`, actual: " + current().value
            );
        }

        moveNext();
        if(current().type == Token.Type.RBRACE) {
            return new JsonAst.JsonObject(new ArrayList<>());
        }

        List<Pair<JsonAst.JsonString, JsonAst.JsonValue>> members = 
            new ArrayList<>();
        var pair= parsePair();
        members.add(pair);

        while(moveNext()) {
            if(current().type == Token.Type.RBRACE) {
                return new JsonAst.JsonObject(members);
            }
            if(current().type != Token.Type.COMMA) {
                throw new ParseException(
                    "expected: `,`, actual: " + current().value
                );
            }
            moveNext();
            pair = parsePair();
            members.add(pair);
        }

        throw new ParseException("unexpected EOF");
    }
\end{codescreen}

まず、最初のif文で、次のトークンが\texttt{\{}であることを確認した後に、

\begin{itemize}
\item
  その次のトークンが\texttt{\}}であった場合：空オブジェクトを返す
\item
  それ以外の場合： \texttt{parsePair()} を呼び出し、
  \texttt{string:value} のようなペアを解析した後、以下のループに突入

  \begin{itemize}
  \item
    次のトークンが\texttt{\}}の場合、集めたペアのリストを引数として、\texttt{JsonAst.JsonObject()}オブジェクトを作って返す
  \item
    それ以外で、次のトークンが\texttt{,}でない場合、構文エラーを投げて終了
  \item
    それ以外の場合：次のトークンをフェッチして来て、\texttt{parsePair()}を呼び出して、ペアを解析した後、リストにペアを追加
  \end{itemize}
\end{itemize}

\noindent のような動作を行います。実際のJSONのオブジェクトと対応付けてみると、より理解が進むでしょう。

\subsection{parseArray}

\texttt{parseArray()}メソッドは、規則\texttt{array}に対応するメソッドで、JSONの配列リテラルに対応するものを解析するメソッドでもあります。実装を示すと以下のようになります。

\begin{codescreen}{}
    /** BNF: array = LBRACKET RBRACKET | LBRACKET value {COMMA value} RBRACKET */
    private JsonAst.JsonArray parseArray() {
        if(current().type != Token.Type.LBRACKET) {
            throw new ParseException(
                "expected: `[`, actual: " + current().value
            );
        }

        moveNext();
        if(current().type == Token.Type.RBRACKET) {
            return new JsonAst.JsonArray(new ArrayList<>());
        }

        List<JsonAst.JsonValue> values = new ArrayList<>();
        var value = parseValue();
        values.add(value);

        while(moveNext()) {
            if(current().type == Token.Type.RBRACKET) {
                return new JsonAst.JsonArray(values);
            }
            if(current().type != Token.Type.COMMA) {
                throw new ParseException(
                    "expected: `,`, actual: " + current().value
                );
            }
            moveNext();
            value = parseValue();
            values.add(value);
        }

        throw new ParseException("unexpected EOF");
    }
\end{codescreen}

まず、最初のif文で、次のトークンが\texttt{{[}}であることを確認した後に、

\begin{itemize}
\item
  その次のトークンが\texttt{{]}}であった場合：空の配列
  (\texttt{JsonAst.JsonArray}) を返す
\item
  それ以外の場合： \texttt{parseValue()} を呼び出し、
  \texttt{value}を解析した後、以下のループに突入

  \begin{itemize}
  \item
    次のトークンが\texttt{{]}}の場合、集めた\texttt{values}のリストを引数として、\texttt{JsonAst.JsonArray()}オブジェクトを作って返す
  \item
    それ以外で、次のトークンが\texttt{,}でない場合、構文エラーを投げて終了
  \item
    それ以外の場合：次のトークンをフェッチして来て、\texttt{parseValue()}を呼び出して、\texttt{value}を解析した後、リストに\texttt{value}を追加
  \end{itemize}
\end{itemize}

\noindent のような動作を行います。実際のJSONの配列と対応付けてみると、より理解が進むでしょう。

\texttt{parseArray()}のコードを読めばわかるように、ほとんどのコードは、\texttt{parseObject()}と共通のものになっています。もしこれが気になるようであれば、共通部分をくくりだすこともできます。他の割愛したメソッドも同様に、トークンを読み取って期待される構文を解析するという流れになっています。

\subsection{字句解析器と構文解析器の連携}

字句解析器を使った構文解析の流れを、\texttt{\{"key":\ "value"\}}
を例に説明します。

\begin{enumerate}
\item
  \textgt{字句解析フェーズ}:
  \texttt{tokenizeAll()}が入力全体を一度にトークン列に変換

\begin{codescreen}{}
入力: {"key": "value"}
↓
トークン列: [LBRACE, STRING("key"), COLON, STRING("value"), RBRACE, EOF]
\end{codescreen}

\item
  \textgt{構文解析フェーズ}:
  パーサーがトークン列を走査しながら抽象構文木を構築

\begin{itemize}
\item
  index == 0: \texttt{LBRACE}を見て\texttt{parseObject()}を呼び出し
\item
  index == 1: \texttt{STRING("key")}を読み取り
\item
  index == 2: \texttt{COLON}を確認
\item
  index == 3: \texttt{STRING("value")}を\texttt{parseValue()}で処理
\item
  index == 4: \texttt{RBRACE}で終了を確認
\item
  最終的にオブジェクトの抽象構文木を生成
\end{itemize}
\end{enumerate}

この方式ではトークン列が事前に確定しているため、以下の利点があります。

\begin{itemize}
\item
  パーサーは純粋に構造の解析に集中できる
\item
  トークン列をログ出力してデバッグが容易
\end{itemize}

\section{PEGベースの構文解析器と伝統的な構文解析器の違い}
\index{PEG}
\index{こうぶんかいせきき@構文解析器!PEGベース}
\index{じくかいせき@字句解析}

前節のPEGベースの構文解析器と、この節の字句解析器を使った構文解析器の主な違いは以下のとおりです。

\begin{enumerate}
\item
  関心の分離: 字句解析と構文解析が明確に分離されている
\item
  空白の処理:
  字句解析器が空白を処理するため、構文解析器は空白を意識しない
\item
  バックトラック: 字句解析器を使う方式では通常バックトラックを行わない
\item
  性能: 一般的に字句解析器を使う方式の方が高速
\end{enumerate}

PEGベースの構文解析器にはいいとこなしのように見えますが、PEGベースの構文解析器には、以下のような利点もあります。

\begin{enumerate}
\item
  シンプルな実装:
  PEGは直感的で、構文解析のロジックが大幅に簡潔に表現できる
\item
  再帰的な構造の自然な表現:
  再帰的な文法を自然に扱えるため、特にネストされた構造の解析が容易
\end{enumerate}

PEGベースの構文解析器は、トークン列を生成する必要がないため、文字列補間のような、一見トークン化が難しいケースでも、構文解析器の中で直接文字列を解析できるという利点があります。

\section{まとめ}

この章では、JSONの構文解析や字句解析を実際に作ってみることを通して、構文解析の基礎について学んでもらいました。特に、

\begin{itemize}
\item
  JSONの概要
\item
  JSONのBNF
\item
  JSONの構文解析器（PEG版）
\item
  古典的な構文解析器
\item
  JSONの字句解析器
\item
  JSONの構文解析器
\end{itemize}

\noindent といった順番で、JSONの定義から入って、PEGによるJSONパーサー、字句解析器を使った構文解析器の作り方について学んでもらいました。この書籍中で使ったJSONはECMA-404で定義されている正式なJSONのサブセットになっています。たとえば、浮動小数点数が完全に扱えないという制限がありますが、構文解析器全体から見ればささいなことなので、この章を理解できれば、JSONの構文解析について理解できたと思って構いません。

次の章では、文脈自由文法（Context-Free Grammar,
CFG）の考え方について学んでもらいます。文脈自由文法は、現在使われているほとんどの構文解析アルゴリズムの基盤となっている概念であって、CFGの理解なくしては、その後の構文解析の理解もおぼつかないからです。

逆に、CFGの考え方さえわかってしまえば、個別の構文解析アルゴリズム自体は、それほど難しいとは感じられなくなって来るかもしれません。

\begin{center}\rule{0.5\linewidth}{0.5pt}\end{center}
\clearpage
\textgt{演習問題}

\begin{enumerate}
\item
  コメントのサポート:

\begin{itemize}
\item
  JSONのBNF定義を拡張し、\texttt{//}
  から行末までの単一行コメントと、\texttt{/*} から \texttt{*/}
  までの複数行コメントをサポートするようにしてください。

  \begin{itemize}
  \item
    \texttt{PegJsonParser} と \texttt{SimpleJsonTokenizer}
    の両方を修正し、これらのコメントを正しく無視するように実装してください。
  \item
    ヒント: \texttt{PegJsonParser} では \texttt{skipWhitespace}
    にコメントスキップのロジックを追加するか、各解析メソッドの適切な箇所でコメントを読み飛ばす処理を挟みます。\texttt{SimpleJsonTokenizer}
    では \texttt{moveNext} の \texttt{switch}
    文にケースを追加し、そこからコメントの種別を判定して読み飛ばす処理を実装します。
  \end{itemize}
\end{itemize}

\item
  数値型の拡張:

\begin{itemize}
\item
  \texttt{PegJsonParser} の \texttt{parseNumber}
  メソッドと、\texttt{SimpleJsonTokenizer} の \texttt{tokenizeNumber}
  メソッドを修正し、ECMA-404仕様に準拠した数値型（小数部、指数部
  \texttt{e} または \texttt{E}
  を含む）を正しく解析できるようにしてください。

  \begin{itemize}
  \item
    \texttt{JsonAst.JsonNumber} の \texttt{value} フィールドの型を
    \texttt{double} から \texttt{java.math.BigDecimal}
    に変更し、精度が失われないように対応してください。
  \item
    テストケースとして、\texttt{123}, \texttt{-0.5}, \texttt{1.2e3},
    \texttt{0.4E-1} のような多様な数値表現を試してみてください
  \end{itemize}
\end{itemize}
\end{enumerate}

\chapter{文脈自由文法の世界}\label{ch03}
\index{ぶんみゃくじゆうぶんぽう@文脈自由文法}
\index{CFG}
\index{けいしきげんご@形式言語}

第2章では、JSONの構文解析器を記述することを通して、構文解析のやり方を学びました。構文解析器についても、PEG型の構文解析器および字句解析器を使った2通りを作ってみることで、構文解析器といってもいろいろな書き方があるのがわかってもらえたのではないかと思います。

この第3章では、現代の構文解析を語る上で必須である、文脈自由文法という概念について学ぶことにします。「文脈自由文法」というと、一見、堅くて難しそうな印象を持つ方も多いかもしれません。

しかし、実は皆さんはすでに文脈自由文法を使っているのです。Javaの\texttt{if}文やメソッド定義、JSONの入れ子構造など、プログラミングで日常的に扱っている「構造」はすべて文脈自由文法で表現されています。この章では、そんな身近な例から始めて、徐々に文脈自由文法の概念を理解していきましょう。

\section{身近な例から始める文脈自由文法}
\index{ぶんみゃくじゆうぶんぽう@文脈自由文法!例}
\index{いれここうぞう@入れ子構造}
\index{HTML}
\index{JSON}

まず、皆さんが普段書いているJavaコードを見てみましょう。以下はシンプルなif文です。

\begin{codescreen}{}
if (x > 0) {
    System.out.println("正の数です");
}
\end{codescreen}

このif文の構造を言葉で説明すると「\texttt{if}の後に条件式を括弧で囲み、その後に文のブロックがくる」となります。さらに、if文の中にif文を書くこともできます。

\begin{codescreen}{}
if (x > 0) {
    if (x > 100) {
        System.out.println("100より大きい");
    }
}
\end{codescreen}

この「入れ子にできる構造」こそが、文脈自由文法の本質なのです。

この入れ子構造は、プログラミング以外の場面でも頻繁に現れます。

\begin{itemize}
\item
  \textbf{\sffamily HTML}\textgt{の要素}：

\begin{codescreen}{}
<div>
    <p>段落の中に<strong>強調</strong>があり、
       さらに<a href="...">リンク</a>も含まれる</p>
</div>
\end{codescreen}
\item
  \textgt{数式の括弧}：

\begin{codescreen}{}
((2 + 3) × (4 - 1)) ÷ 5
\end{codescreen}
\item
  \textbf{日本語の引用}：

\begin{codescreen}{}
彼は「『明日は晴れる』と言った」らしい。
\end{codescreen}
\item
  \textbf{\sffamily JSON}\textgt{のオブジェクト}：

\begin{codescreen}{}
{
  "user": {
    "profile": {
      "name": "太郎"
    }
  }
}
\end{codescreen}
\end{itemize}

\enlargethispage{\baselineskip}これらの例に共通するのは、ある構造の中に同じような構造を含むことができ、その深さに制限がないという点です。このような構造を正確に記述し、解析するために文脈自由文法が必要になるのです。この問題は、次節で述べる「括弧の対応」という根本的な問題に帰着します。

\section{括弧の対応という根本問題}
\index{かっこのたいおう@括弧の対応}
\index{Dyck言語}
\index{いれここうぞう@入れ子構造}

形式言語で最も基本的で重要な構造の1つが「括弧の対応」です。大抵のプログラミング言語（JSONやYAMLのようなものも含む）は開き括弧と閉じ括弧に対応するルールを持っています。たとえ括弧そのものでなくても、\texttt{begin/end}や\texttt{\{\}}などが任意の深さだけ入れ子にできる場合は同じことです。

かなり原始的なプログラミング言語ですら、数式を入力できる機能があり、かつ数式の中で括弧を使うことができる以上、プログラミング言語の構文解析において括弧の対応は避けて通れない問題です。

さて、ここで括弧の対応とは何かを考えてみましょう。括弧の対応が取れているとは、\texttt{(}および\texttt{)}からなる文字列において、開き括弧と閉じ括弧が正しくペアになっていることを意味します。たとえば、以下は正しい括弧の対応です。

\begin{codescreen}{}
()          // 1組の括弧
(())        // 入れ子になった括弧
(()())      // 入れ子と並列の組み合わせ
()()()      // 並列に並んだ括弧
\end{codescreen}

一方、以下は正しくない例です。

\begin{codescreen}{}
)(          // 順序が逆
(()         // 閉じ括弧が不足
())         // 開き括弧が不足
\end{codescreen}

このような「\texttt{(}および\texttt{)}からなる文字列で、括弧の釣り合いが取れた文字列の並び」を表すものをDyck（ディック）言語と呼びます。「Dyck」は数学者の名前に由来しています。この、一見単純に見える問題が、構文解析において極めて重要な位置を占めているのです。

\section{BNFで括弧の構造を表現する}
\index{BNF}
\index{さいき@再帰}
\index{くうもじれつ@空文字列}
\index{いぷしろん@ε}

この括弧の対応をどのように文法として表現すればよいでしょうか？まずはすでに皆さんに学んでもらったBNFを使って考えてみましょう。

括弧の構造には以下の2つのパターンがあります。

\begin{enumerate}
\item
  空文字列（括弧なし）
\item
  \texttt{(} + 内側の括弧構造 + \texttt{)} + 続きの括弧構造
\end{enumerate}

これをBNFで書くと以下のようになります。

\begin{codescreen}{}
P = "(" P ")" P | "";
\end{codescreen}

\texttt{P}は「括弧のパターン」を表します。この定義は再帰的になっていることに注目してください。\texttt{P}の定義の中に\texttt{P}自身が現れています。これこそが、入れ子構造を表現する鍵なのです。

このBNFは上記の2つのパターンと直接対応しています。

\begin{itemize}
\item
  \texttt{"("\ P\ ")"\ P}：開き括弧、内側の構造、閉じ括弧、続きの構造（パターン1）
\item
  \texttt{""}：空文字列（パターン2）
\end{itemize}

\section{BNFから文脈自由文法へ}
\index{BNF}
\index{ぶんみゃくじゆうぶんぽう@文脈自由文法}
\index{せいせいきそく@生成規則}
\index{ひしゅうたんきごう@非終端記号}
\index{しゅうたんきごう@終端記号}
\index{かいしきごう@開始記号}

前節のBNFはすでに文脈自由文法の一種です。ただし、文脈自由文法の議論をする際には、より標準的な記法を使います。段階的に変換してみましょう。

\begin{itemize}
\item
  ステップ1: 選択を分離

まず、\texttt{\textbar{}}で区切られた選択肢を別々の規則に分けます。

\begin{codescreen}{}
P = "(" P ")" P;
P = "";
\end{codescreen}

同じ\texttt{P}が2つの規則で定義されていますが、これは「Pは2つのパターンのどちらかになる」という意味です。

\item
  ステップ2: 記号の変更

次に、記法を数学的な標準形に近づけます。

\begin{itemize}
\item
  \texttt{=} を \texttt{→}（生成規則を表す矢印）に変更
\item
  \texttt{""} を \texttt{ε}（イプシロン：空文字列を表す記号）に変更
\item
  文字を囲む引用符を削除
\end{itemize}

\begin{codescreen}{}
P → ( P ) P
P → ε
\end{codescreen}
\end{itemize}

これで文脈自由文法の標準的な記法になりました！

\subsection{用語の整理}

ここで重要な用語を整理しておきましょう。

\begin{itemize}
\item
  生成規則：\texttt{P\ →\ (\ P\ )\ P} のような矢印で結ばれた規則
\item
  非終端記号：\texttt{P}のような、さらに展開される記号（変数のようなもの）
\item
  終端記号：\texttt{(}や\texttt{)}のような、これ以上展開されない記号（実際の文字）
\item
  開始記号：文法の起点となる非終端記号（この例では\texttt{P}）
\end{itemize}

つまり、文脈自由文法とは以下のように表現できます。

\begin{itemize}
\item
  生成規則の集まり

  \begin{itemize}
  \item
    各生成規則は「非終端記号 → 記号の並び」の形
  \item
    記号の並びは終端記号と非終端記号の組み合わせ（空文字列εも可）
  \end{itemize}
\end{itemize}

これだけのシンプルなルールで、プログラミング言語の複雑な構文を表現できるのです。

\section{実例で理解する「言語」の概念}
\index{げんご@言語}
\index{けいしきげんご@形式言語}
\index{もじれつ@文字列}
\index{しゅうごう@集合}

前の節で定義したDyck言語の文法をもう一度見てみましょう。ここでは、文法の入口を明示するために、開始記号\texttt{D}とそれをPに置き換える規則\texttt{D → P}を先頭に加えています。

\begin{codescreen}{}
D → P
P → ( P ) P
P → ε
\end{codescreen}

この文法は括弧の対応を表すもの、というふうにぼかしてきました。実のところ、これまで出てきた文法はすべて「形式言語」を定義するものです。でも、ちょっと待ってください。「形式言語」って何でしょうか？

私たちが普段「プログラミング言語」と呼んでいるJavaやPythonと、この3行の規則で表される「括弧の言語」は、どう関係しているのでしょうか？実は、どちらも同じ「形式言語」という概念で理解できるのです。

\subsection{「形式言語」とは}

私たちが普段使う「言語」という言葉は、自然言語（日本語や英語など）を指すことが多いですが、理論の世界では「形式言語」という概念が重要です。形式言語とは、\textgt{文字列の集合}として定義される言語のことです。

ここで言う集合とは数学の集合のことです。形式言語の世界では文字列をひたすら集めていった、その集合が「言語」と呼ばれるのです。

単に「aabcaa」のような意味のない文字列の羅列に感じるかもしれませんが、重要なのは「どの文字列がその言語に属するか」という\textgt{規則}です。たとえば

\begin{itemize}
\item
  \textgt{パスワードの言語}：「8文字以上で、英数字と記号を含む文字列の集合」
\item
  \textgt{電話番号の言語}：「数字とハイフンからなる特定のパターンに従う文字列の集合」
\item
  \textgt{プログラミング言語}：「その言語の文法に従う正しいプログラムの集合」
\end{itemize}

つまり、形式言語は単なる文字列の寄せ集めではなく、明確な規則に基づいて「その言語に属する/属さない」が判定できる文字列の集合なのです。

さて、この観点から「Java言語」や「JavaScript言語」というとき、それは何を指しているのでしょうか？

\subsection{言語を「文字列の集合」として考える}

形式言語の立場からは、プログラミング言語を\textgt{その言語で書ける正しいプログラムすべての集合}として定義することができます。「正しいプログラム」といっても、構文解析ができるもの、型検査を通るもの、実行時にエラーが出ないものなど、いくつかの基準がありますが、ここでは「構文解析が通るプログラム」と考えます。

具体例で考えてみましょう。以下はすべて正しいJavaScriptプログラムです。

\begin{codescreen}{}
console.log("Hello, World!");
\end{codescreen}

\begin{codescreen}{}
console.log(3);
\end{codescreen}

\begin{codescreen}{}
console.log(3 + 5);
\end{codescreen}

これらを集めていくと、JavaScriptという言語（以後\texttt{JS}と表記）は次のような文字列の集合として表現できます。

\begin{codescreen}{}
JS = {
  "console.log(\"Hello, World!\");",
  "console.log(3);",
  "console.log(3 + 5);",
  "const x = 10;",
  "function add(a, b) { return a + b; }",
  ...（無限に続く）
}
\end{codescreen}

\enlargethispage{\baselineskip}JSで書ける正しいプログラムは無限にあるので、この集合は\textgt{無限集合}（正確には可算無限個の集合）になります。

\subsection{Dyck言語も集合として理解する}

同様に、Dyck言語（括弧の対応が取れた文字列）も集合として表現できます。

\begin{codescreen}{}
Dyck = {
  "()",
  "(())",
  "((()))",
  "(()())",
  "()()",
  "()()()",
  ...（無限に続く）
}
\end{codescreen}

Dyck言語も無限集合です。基本的に私たちが扱いたい言語＝集合は、ほとんどが無限集合です。なぜなら、基本的にはその言語のテキストがいくら長くなってもいいようにしたいからです。言語であらかじめ「長さの制限」を設けることは、実用的ではありません。

\subsection{集合として言語を扱うメリット}

言語を集合として扱うと、数学の集合論で使う記号が使えるようになります。これには実用的なメリットがあります。

まず、基本的な記号の意味を説明します。

\begin{itemize}
\item
  \texttt{∈}：「含まれる」を意味します（要素と集合の関係）
\item
  \texttt{∉}：「含まれない」を意味します（要素と集合の関係）
\item
  \texttt{⊂}：「部分集合」を意味します（集合と集合の関係）
\item
  \texttt{∩}：「共通部分」を意味します（両方の集合に含まれる要素）
\item
  \texttt{∪}：「和集合」を意味します（どちらかの集合に含まれる要素）（第6節で使用）
\end{itemize}

重要な違いとして、\texttt{∈}は要素（文字列）と集合（言語）の関係を表し、\texttt{⊂}は集合と集合の関係を表します。たとえば、以下のようになります。

\begin{codescreen}{}
"()" ∈ Dyck        // 「文字列"()"」が「言語Dyck」に含まれる
{""} ⊂ Dyck        // 「空文字列だけからなる集合」は「言語Dyck」の部分集合
\end{codescreen}

\begin{itemize}
\item
  例1：ある文字列が言語に含まれるかの判定

\begin{codescreen}{}
"()" ∈ Dyck         // "()"はDyck言語に含まれる
")(" ∉ Dyck         // ")("はDyck言語に含まれない
\end{codescreen}
\item
  例2：言語の後方互換性

Java 8がJava 5の後方互換であることを、集合の包含関係で表現できます。

\begin{codescreen}{}
Java5 ⊂ Java8     // Java 5で書けるプログラムはすべてJava 8でも書ける
\end{codescreen}
\item
  例3：言語の共通部分

たとえば、「JSでもTSでも有効なプログラム」は以下のように表現できます。

\begin{codescreen}{}
JS ∩ TS = {
  "console.log(123);",
  "const x = 5;",
  "function f() { return 1; }",
  ...
}
\end{codescreen}
\end{itemize}

このように、言語を集合として扱うことで、言語間の関係を明確に表現できるようになります。これは単なる理論ではなく、言語設計や互換性の議論において非常に役立ちます。

\section{文脈自由文法の数学的定義}
\index{ぶんみゃくじゆうぶんぽう@文脈自由文法!数学的定義}
\index{せいせいきそく@生成規則}
\index{ひしゅうたんきごう@非終端記号}
\index{しゅうたんきごう@終端記号}
\index{かいしきごう@開始記号}

ここまでふわっとした説明をしてきましたが、文脈自由文法の定義はもう少し厳密に行えます。本書では厳密な数学的扱いよりも実用的な理解を重視するため、あくまで厳密にはこう定義されるという話であって、この定義を直接使うことはありません。興味のない方は読み飛ばしても構いません。

文脈自由文法（Context-Free Grammar, CFG）は、形式的には4つ組
\(G = (V, \Sigma, P, S)\) として定義されます。

\begin{itemize}
\item
  \(V\)：非終端記号（Non-terminal symbols）の有限集合
\item
  \(\Sigma\)：終端記号（Terminal
  symbols）の有限集合（\(V \cap \Sigma = \emptyset\)）
\item
  \(P\)：生成規則（Production rules）の有限集合

  \begin{itemize}
  \item
    各規則は \(A \rightarrow \alpha\)
    の形（\(A \in V\)、\(\alpha \in (V \cup \Sigma)^*\)）
  \end{itemize}
\item
  \(S\)：開始記号（Start symbol）（\(S \in V\)）
\end{itemize}

たとえば、Dyck言語の文法は数学的には以下のように表せます。

\begin{itemize}
\item
  \(V = \{P\}\)
\item
  \(\Sigma = \{(, )\}\)
\item
  \(P = \{P \rightarrow (P)P, P \rightarrow \varepsilon\}\)
\item
  \(S = P\)
\end{itemize}

ここで \(\varepsilon\) は空文字列を表します。

文法 \(G\) が生成する言語 \(L(G)\) は、開始記号 \(S\)
から生成規則を繰り返し適用して導出できるすべての終端記号列の集合として定義されます。

\[L(G) = \{w \in \Sigma^* \mid S \Rightarrow^* w\}\]

ここで、\(\Rightarrow^*\) は「0回以上の導出ステップ」を表します。

もう少し噛み砕いて説明すると

\begin{enumerate}
\item
  開始記号 \(S\) から出発する
\item
  生成規則を適用して、非終端記号を置き換えていく
\item
  すべてが終端記号になったら、それが言語 \(L(G)\) の要素の1つ
\end{enumerate}

たとえば、Dyck言語の文法では以下のようになります。

\begin{itemize}
\item
  \(P \Rightarrow (P)P \Rightarrow ()P \Rightarrow ()\)
  という導出により、\texttt{()} が \(L(G)\) に含まれる
\item
  \(P \Rightarrow (P)P \Rightarrow ((P)P)P \Rightarrow (()P)P \Rightarrow (())P \Rightarrow (())\)
  により、\texttt{(())} も \(L(G)\) に含まれる
\item
  このようにして、すべての「正しく対応した括弧列」が \(L(G)\)
  の要素となる
\end{itemize}

この数学的定義を「きっちりと」理解することは、理論的には重要ですが、構文解析器の実装では、直感的な理解で問題ありません。BNFなどの記法は数学的定義と本質的に同じものを、より読みやすい形で表現しています。

\section{正規表現の限界と文脈自由言語}
\index{せいきひょうげん@正規表現}
\index{せいきげんご@正規言語}
\index{ぶんみゃくじゆうげんご@文脈自由言語}
\index{おーとまとん@オートマトン}
\index{ゆうげんおーとまとん@有限オートマトン}

ここまでで文脈自由文法の基本的な概念と言語を集合として扱う考え方を学びました。BNFと文脈自由文法の表現能力が実質的に等価であることはなんとなく理解できたかと思います。しかし、ここで1つ重大な疑問が浮かび上がります。それは、私たちにとってなじみ深い正規表現との関係です。正規表現も文脈自由文法と同じように言語＝文字列の集合を表現するためのツールですが、正規表現と文脈自由文法にはどのような違いがあるのでしょうか？

\subsection{身近な正規表現から考える}

皆さんは日常的に正規表現を使っているでしょう。ファイル検索、テキスト処理、入力値の検証など、さまざまな場面で活躍しています。

たとえば、Javaで電話番号の形式を簡易チェックするとき、以下のような正規表現を使うことがあります。

\begin{codescreen}{}
String phonePat = "\\d{3}-\\d{4}-\\d{4}";  // 例：090-1234-5678
if (phone.matches(phonePat)) {
    // 有効な電話番号
}
\end{codescreen}

あるいは、郵便番号の簡易チェックでは以下のような正規表現を使います。

\begin{codescreen}{}
String postalCodePat = "\\d{3}-\\d{4}";  // 例：123-4567
if (postalCode.matches(postalCodePat)) {
    // 有効な郵便番号
}
\end{codescreen}

正規表現は、特定のパターンにマッチする文字列を簡潔に表現できる強力なツールです。また、正規表現も無限集合を扱えます。上記の正規表現はあくまで有限個のパターンしか表せませんが、以下のパターンでは自然数を無限に表現できます。

\begin{codescreen}{}
String natPat = "0|[1-9][0-9]*";  // 0以上の自然数
if (number.matches(natPat)) {
    // 有効な自然数
}
\end{codescreen}

このとき、\texttt{natPat}は以下のような無限集合を表します。

\begin{codescreen}{}
{'0', '1', '2', '3', ...}
\end{codescreen}

注意してほしいのは、これはあくまで文字列としての自然数の集合であり、数値としての自然数ではないということです。正規表現は文字列のパターンを表現するためのツールであり、数値そのものを扱うわけではありません。

\subsection{正規表現でできないこと}

しかし、正規表現には決定的な限界があります。それは\textgt{括弧の対応が取れているかチェックできない}ということです。

試しに、すでにでてきたDyck言語を正規表現で判定することを考えてみましょう。

\begin{codescreen}{}
OK:  (), (()), (()(()))
NG:  )(, ((), ())
\end{codescreen}

どんなに工夫しても、任意の深さの括弧の対応を正規表現で表現することはできません。なぜでしょうか？次からは正規表現の基本的な仕組みとその限界について詳しく見ていきましょう。

\subsection{正規表現の構成要素を詳しく見る}

正規表現は、実は非常にシンプルな3つの基本演算から構成されています。現代の正規表現エンジンは多くの便利な記法を提供していますが、理論的にはすべて以下の1-3の基本演算に帰着できます（ただし、後方参照など「非正規表現」とでも呼ぶべき機能を除く）。

\subsubsection{連接（Concatenation）}

2つの正規表現を続けて書くことを\textgt{連接}と呼びます。

\begin{codescreen}{}
正規表現: ab
マッチする文字列: "ab"
マッチしない文字列: "a", "b", "ba", "abc"
\end{codescreen}

連接は最も基本的な演算で、「次に」という順序関係を表現します。

\subsubsection{選択（Alternation）}

\texttt{\textbar{}}記号を使って、複数の選択肢を表現します。

\begin{itemize}
\item
  正規表現: \texttt{a\textbar{}b}
\item
  マッチする文字列: \texttt{"a"}, \texttt{"b"}
\item
  マッチしない文字列: \texttt{"ab"}, \texttt{"c"}, \texttt{""}
\end{itemize}

より複雑な例：

\begin{itemize}
\item
  正規表現: \texttt{(ab\textbar{}cd)}
\item
  マッチする文字列: \texttt{"ab"}, \texttt{"cd"}
\item
  マッチしない文字列: \texttt{"ac"}, \texttt{"bd"}, \texttt{"abcd"}
\end{itemize}

\subsubsection{繰り返し（Kleene Star）}

\texttt{*}記号は、直前の要素を0回以上繰り返すことを意味します。この演算は数学者Stephen
Cole Kleene（スティーブン・クリーネ）にちなんで「Kleene
Star」と呼ばれます。

\begin{itemize}
\item
  正規表現: \texttt{a*}
\item
  マッチする文字列: \texttt{""}, \texttt{"a"}, \texttt{"aa"},
  \texttt{"aaa"}, \ldots{}
\item
  マッチしない文字列: \texttt{"b"}, \texttt{"ab"}
\end{itemize}

\subsubsection{派生的な演算子}

現代の正規表現では、より便利にするために多くの派生的な演算子が追加されています。

\begin{codescreen}{}
a+       1回以上の繰り返し（aa*と同等）
a?       0回または1回（a|εと同等）
a{n}     ちょうどn回の繰り返し
a{n,m}   n回以上m回以下の繰り返し
[a-z]    文字クラス（a|b|c|...|zと同等）
.        任意の1文字
\end{codescreen}

これらはすべて基本の3演算で表現できます。たとえば代表的な例をあげると以下のようになります。

\begin{itemize}
\item
  \texttt{a+} は \texttt{aa*}
\item
  \texttt{a?} は \texttt{a\textbar{}ε}（εは空文字列）
\item
  \texttt{{[}abc{]}} は \texttt{a\textbar{}b\textbar{}c}
\end{itemize}

\subsection{正規表現の裏側にある仕組み}

\enlargethispage{\baselineskip}皆さんは正規表現を使っているとき、その裏側でどのような処理が行われているか考えたことはありますか？実は、正規表現の限界を理解するには、その背後にある「オートマトン」という仕組みを知ることが役立ちます。「オートマトン」という名前は難しそうですが、要は「文字列を一文字ずつ読みながら状態を変えていく機械」のことです。

\subsubsection{オートマトンを自動販売機で考える}

オートマトンは「有限個の状態を持つ機械」です。難しく聞こえるかもしれませんが、自動販売機を例にするとわかりやすくなります。ここで、すべての商品は100円とし、投入可能な上限金額は300円とします。

\begin{itemize}
\item
  初期状態：お金が入っていない

  \begin{itemize}
  \item
    商品ボタンを押す → 初期状態のままで、お金が不足している旨を表示
  \item
    100円を投入 → 100円状態へ遷移
\end{itemize}
\item
  100円状態

  \begin{itemize}
  \item
    商品ボタンを押す → 商品を出して初期状態へ遷移
  \item
    商品ボタンを押さず100円投入 → 200円状態へ遷移
  \end{itemize}
\item
  200円状態

  \begin{itemize}
  \item
    商品ボタンを押す → 商品と、お釣り100円をだして初期状態へ遷移
  \item
    商品ボタンを押さず100円投入 → 300円状態へ遷移
  \end{itemize}
\item
  300円状態

  \begin{itemize}
  \item
    商品ボタンを押す → 商品と、お釣り200円をだして初期状態へ遷移
  \item
    商品ボタンを押さず100円投入 → 300円状態のまま、100円を返却
  \end{itemize}
\end{itemize}

図にすると以下のようになります。

\begin{center}
\begin{tikzpicture}[
  ->,>=stealth',
  node distance=3cm,
  initial text=start,
  every state/.style={thick, fill=gray!10},
  accepting/.style={double distance=2pt}
]

% 状態の定義
\node[state,initial] (q0) {初期状態};
\node[state] (q1) [right of=q0] {100円};
\node[state] (q2) [right of=q1] {200円};
\node[state] (q3) [right of=q2] {300円};

% 遷移の定義
% 100円投入による遷移
\path (q0) edge[bend left] node[above] {100円} (q1);
\path (q1) edge[bend left] node[above] {100円} (q2);
\path (q2) edge[bend left] node[above] {100円} (q3);
\path (q3) edge[loop above] node {100円/返} (q3);

% 商品ボタンによる遷移
\path (q0) edge[loop below] node {ボ/不} (q0);
\path (q1) edge[bend right=-30] node[below] {ボ/商} (q0);
\path (q2) edge[bend right=40] node[below] {ボ/（商，釣）} (q0);
\path (q3) edge[bend right=50] node[below] {ボ/（商，釣）} (q0);

\end{tikzpicture}
\end{center}

この図の見方を説明します。

\begin{itemize}
\item
  \textgt{ノード（状態）}：

  \begin{itemize}
  \item
    円で表現される各ノードが「状態」を表します
  \item
    「start」と書かれた矢印が指す「初期状態」から開始します
  \item
    各状態には現在の投入金額が記載されています
  \end{itemize}
\item
  \textgt{矢印（遷移）}：

  \begin{itemize}
  \item
    矢印は状態間の「遷移」を表します
  \item
    矢印上のラベルは「入力/出力」の形式です

    \begin{itemize}
    \item
      「100円」：100円硬貨を投入
    \item
      「ボ/商」：商品ボタンを押すと商品が出る
    \item
      「ボ/不」：商品ボタンを押すがお金が不足
    \item
      「ボ/（商，釣）」：商品ボタンを押すと商品とお釣りが出る
    \item
      「100円/返」：100円を投入するが満杯なので返却される
    \end{itemize}
  \end{itemize}
\item
  \textgt{動作の例}：

  \begin{enumerate}
  \item
    初期状態から100円を2回投入 → 200円状態へ
  \item
    200円状態で商品ボタンを押す → 商品とお釣りが出て初期状態へ戻る
  \end{enumerate}
\end{itemize}

この図は、自動販売機の「状態遷移」を完全に表現しています。どの状態からどの入力を受けたらどの状態に遷移するかが一目でわかります。

このように、オートマトンは「状態」と「入力」によって遷移を決定します。自動販売機では、状態はお金の投入状況で、入力はお金の投入や商品ボタンの押下です。

\subsubsection{簡単な例：ab*を認識するオートマトン}

このオートマトン（有限状態機械）を使うと正規表現で表現される言語を認識できます（正規表現で認識可能な言語はすべてオートマトンで表現可能と言い換えてもいいです）。ここでは、正規表現の例として
\texttt{ab*} を認識するオートマトンを考えてみましょう。

この正規表現は以下のような文字列にマッチします。

\begin{itemize}
\item
  \texttt{a}（bが0個）
\item
  \texttt{ab}（bが1個）
\item
  \texttt{abb}（bが2個）
\item
  \texttt{abbb}（bが3個）
\item
  \ldots{}
\end{itemize}

この正規表現を認識するオートマトンは以下のようになります。

\begin{center}
\begin{tikzpicture}[
  >=stealth',
  node distance=3cm,
  state/.style={circle, draw, minimum size=10mm},
  accept/.style={state, double}
]
  % 状態の定義
  \node[state] (s0) {0};
  \node[accept] (s1) [right of=s0] {1};

  % 遷移の定義
  \draw[->] (s0) -- node[above] {$a$} (s1);
  \draw[->] (s1) edge[loop above] node[above] {$b$} ();

  % 開始状態の矢印
  \draw[->] ([xshift=-1cm]s0.west) -- node[above] {開始} (s0);
\end{tikzpicture}
\end{center}

\begin{itemize}
\item
  状態0（開始状態）：最初の状態
\item
  状態1（受理状態）：この状態で入力が終わればマッチ成功
\item
  矢印：文字を読んだときの状態遷移
\end{itemize}

このオートマトンは以下のように動作します。

\begin{itemize}
\item
  入力\verb|"a"|：状態0→状態1（\textgt{受理}）
\item
  入力\verb|"ab"|：状態0→状態1→状態1（\textgt{受理}）
\item
  入力\verb|"abb"|：状態0→状態1→状態1→状態1（\textgt{受理}）
\item
  入力\verb|"b"|：状態0で\verb|"b"|を読めない（\textgt{非受理}）
\end{itemize}

\subsubsection{2種類のオートマトン}

オートマトンには大きく分けて2種類あります。難しい名前ですが、それぞれの特徴を簡単に説明します。

\begin{itemize}
\item
  \textbf{\sffamily NFA（Non-deterministic Finite
  Automaton}\textgt{、非決定性有限オートマトン）}

  \begin{itemize}
  \item
    「非決定性」とは、同じ文字を読んだときに複数の選択肢があることを指します
  \item
    文字を読まずに状態を移動できる（ε遷移）
  \item
    正規表現から比較的簡単に作れる
  \end{itemize}
\item
  \textbf{\sffamily DFA（Deterministic Finite Automaton}\textgt{、決定性有限オートマトン）}

  \begin{itemize}
  \item
    「決定性」とは、各状態で各文字に対して次の状態が一意に決まることを指します
  \item
    文字を読まずに移動することはできない
  \item
    実装が簡単で、実行速度が速い
  \end{itemize}
\end{itemize}

この2つのオートマトンは正規表現（正確には後で説明する正規言語）を認識するモデルであり、能力的には等価です。どちらも正規表現を認識できますが、実装の複雑さや速度に違いがあります。

ちなみに、これまで見てきたオートマトンはすべて状態遷移が一意に決まるDFAです。DFAは状態遷移が一意に決まるため、実装が非常にシンプルで高速です。

一方で、NFAはそのままでは実装が複雑で遅くなりますが、正規表現からの変換が容易であるという利点があります。

\subsection{なぜ正規表現では括弧の対応が扱えないのか}

ここまで、正規表現がオートマトンという仕組みで動作することを学びました。では、なぜオートマトン（正規表現）では括弧の対応がチェックできないのでしょうか？その理由を理解するために、簡単な例から考えてみましょう。

括弧の対応を確認するには、以下のような処理が必要です。

\begin{itemize}
\item
  \texttt{(}を読んだら、「今までに読んだ開き括弧の数」を1増やす
\item
  \texttt{)}を読んだら、「対応する開き括弧があるか」確認して、1減らす
\item
  最後に、カウントが0であることを確認する
\end{itemize}

ここで問題になるのは、括弧は何重にでもネスト（入れ子）できるということです。たとえば以下のようなケースを扱える必要があります。

\begin{itemize}
\item
  \texttt{(())}：2重
\item
  \texttt{((()))}：3重
\item
  \texttt{(((())))}：4重
\item
  \ldots 無限に続く
\end{itemize}

\subsubsection{オートマトンの限界：有限の状態しか持てない}

オートマトンは「有限個の状態」しか持てません。仮に10個の状態を持つオートマトンを作ったとしても、11重にネストした括弧は扱えません。どんなに状態数を増やしても、それを超える深さのネストが存在します。

これを具体的に理解するために、簡単なオートマトンで考えてみましょう。仮に「最大3つまでの開き括弧を数える」オートマトンを作ると次のようになります。

\begin{center}
\begin{tikzpicture}[
  >=stealth',
  node distance=2.5cm,
  state/.style={circle, draw, minimum size=10mm},
  accept/.style={state, double},
  every edge/.style={draw, ->, >=stealth'}
]
  % 状態の定義
  \node[accept] (s0) {0};
  \node[state] (s1) [right of=s0] {1};
  \node[state] (s2) [right of=s1] {2};
  \node[state] (s3) [right of=s2] {3};
  \node[state] (err) [below of=s1] {エラー};

  % 開き括弧による遷移
  \draw (s0) edge[bend left] node[above] {(} (s1);
  \draw (s1) edge[bend left] node[above] {(} (s2);
  \draw (s2) edge[bend left] node[above] {(} (s3);
  \draw (s3) edge[bend right=-40] node[right] {(} (err);

  % 閉じ括弧による遷移
  \draw (s1) edge[bend left] node[below] {)} (s0);
  \draw (s2) edge[bend left] node[below] {)} (s1);
  \draw (s3) edge[bend left] node[below] {)} (s2);
  \draw (s0) edge node[left] {)} (err);

  % 開始状態の矢印
  \draw[->] ([xshift=-1cm]s0.west) -- (s0);
\end{tikzpicture}
\end{center}

このオートマトンは3つまでの開き括弧なら正しく数えられますが、4つ目の開き括弧がくるとエラー状態に遷移してしまいます。状態数を100個に増やしても、101個の開き括弧には対応できません。

\subsection{実用上の意味}

この違いは実用上極めて重要です。以下のような違いがあるからです。

\begin{itemize}
\item
  正規表現で十分な例：

  \begin{itemize}
  \item
    URLの検証
  \item
    電話番号のフォーマットチェック
  \item
    単純なトークンの切り出し
  \item
    メールアドレスの簡単な検証
  \end{itemize}
\item
  文脈自由文法が必要な例：

  \begin{itemize}
  \item
    プログラミング言語の構文解析
  \item
    JSONやYAML、XMLのパース
  \item
    数式の評価（括弧を含む）
  \item
    メールアドレスの完全な検証
  \end{itemize}
\end{itemize}

\subsection{メールアドレスと正規表現の注意点}

メールアドレスの検証は簡単だと思われがちですが、完全な仕様（RFC
5322）は非常に複雑で、厳密にはコメントの入れ子などを許可するため、正規表現では表現しきれません。たとえば、以下のような入れ子構造をもったメールアドレスも仕様上は有効です。

\begin{codescreen}{}
john.doe@(comment)example.com
john.doe@(comment1(comment2))example.com
\end{codescreen}

実用上はaddr-spec（\texttt{local-part@domain}の形式）の簡略化されたパターンを正規表現で検証するのが一般的ですが、完全な仕様準拠が必要な場合は、専用のパーサーを使用するべきです。

メールアドレスの検証ですら、完全な仕様準拠を目指すなら文脈自由文法の理解が必要になるのです。

\subsection{なぜ「文脈自由」なのか？}

ここまで文脈自由文法について学んできましたが、なぜ「文脈自由」という名前なのでしょうか？これは、対となる概念である「文脈依存文法」と比較すると理解しやすくなります。

\bigskip
\noindent\textgt{文脈自由の例}：

\begin{codescreen}{}
A → abc
\end{codescreen}

この規則では、Aはどこに現れても常に\texttt{abc}に置き換えられます。周りの文脈（前後の文字）に関係なく適用できるので「文脈自由」です。自由というと一見「好きにできる」という印象を受けますが、ここの「自由」は、たとえば、アルコールフリー＝アルコールが含まれない、という意味での「自由」です。形式言語を扱う研究者の中にはContext-free Grammarを「文脈自由文法」と訳したのは誤解を招きやすいので、あまりよくなかったのではないか、という人もいます。

\bigskip
\noindent\textgt{文脈依存の例}：

\begin{codescreen}{}
xAy → xabcy
\end{codescreen}

この規則では、Aが\texttt{x}と\texttt{y}に挟まれている時だけ\texttt{abc}に置き換えられます。文脈（周りの文字）に依存しているので「文脈依存」です。

プログラミング言語の構文のほとんどは文脈自由文法で表現できます。たとえば、\texttt{if}文の中に\texttt{if}文を書けるのは、外側の文脈に関係なく同じ規則が適用できるからです。一方、「変数は使用前に宣言されていなければならない」のような規則は文脈依存であり、構文解析とは別の段階（意味解析）で扱われます。

以下で形式言語の階層をまとめてみます。

\begin{itemize}
\item
  正規言語（RL: Regular Language）：

  \begin{itemize}
  \item
    正規表現で表現できる\textgt{言語の集合}
  \item
    例：電話番号、識別子、単純なパターン
  \end{itemize}
\item
  文脈自由言語（CFL: Context-Free Language）：

  \begin{itemize}
  \item
    文脈自由文法で表現できる\textgt{言語の集合}
  \item
    例：プログラミング言語の構文、JSON、XML
  \item
    正規言語をすべて含む（RL ⊂ CFL）
  \end{itemize}
\item
  文脈依存言語（CSL: Context-Sensitive Language）:

  \begin{itemize}
  \item
    文脈依存文法で表現できる\textgt{言語の集合}
  \item
    例：$a^nb^nc^n$ （n個のa、b、cがこの順で並ぶ集合）
  \item
    文脈自由言語をすべて含む（CFL ⊂ CSL）
  \end{itemize}
\item
  帰納的可算言語（RE: Recursively Enumerable Language）:

  \begin{itemize}
  \item
    （汎用）コンピュータが出力できるすべての記号列の集合
  \item
    JavaやPythonなどのプログラミング言語の計算能力はこのレベル
  \item
    文脈依存言語をすべて含む（CSL ⊂ RE）
  \end{itemize}
\end{itemize}

階層を図にすると次のようになります。この図は言語の表現力の階層を表しており、内側の言語クラスは外側の言語クラスに完全に含まれます（包含関係）。

\begin{center}
\begin{tikzpicture}[scale=0.9, transform shape]
  % 最外側の楕円 (RE)
  \draw[fill=purple!10,draw=purple!60,thick] (0,0) ellipse (6cm and 4cm);
  \node at (5.0,0) {\textcolor{black!60}{\large RE}};

  % 文脈依存言語 (CSL)
  \draw[fill=blue!10,draw=blue!60,thick] (0,0) ellipse (4.5cm and 3cm);
  \node at (3.5,-0.7) {\textcolor{black!60}{\large CSL}};

  % 文脈自由言語 (CFL)
  \draw[fill=green!10,draw=green!60,thick] (0,0) ellipse (3.25cm and 2.125cm);
  \node at (2.2,-1.2) {\textcolor{black!60}{\large CFL}};

  % 正規言語 (RL)
  \draw[fill=orange!10,draw=orange!60,thick] (0,0) ellipse (2cm and 1.25cm);
  \node at (0,0) {\textcolor{black!60}{\large RL}};
\end{tikzpicture}
\end{center}

プログラミング言語の\textgt{構文}は文脈自由文法で記述できますが、プログラミング言語自体の\textgt{計算能力}はチューリング完全（最上位）です。この違いを理解することが重要です。

\section{文法から文字列を作る：導出の仕組み}
\index{どうしゅつ@導出}
\index{さいさどうしゅつ@最左導出}
\index{さいうどうしゅつ@最右導出}
\index{かいせきぎ@解析木}

文脈自由文法は「生成規則」の集まりでした。では、この「生成」とは何でしょうか？

Dyck言語の文法をもう一度見てみましょう。

\begin{codescreen}{}
D → P
P → ( P ) P
P → ε
\end{codescreen}

これらの規則は「置き換えルール」として読めます。

\begin{itemize}
\item
  \texttt{D\ →\ P}：「Dを見たらPに置き換えてよい」
\item
  \texttt{P\ →\ (\ P\ )\ P}：「Pを見たら( P ) Pに置き換えてよい」
\item
  \texttt{P\ →\ ε}：「Pを見たら空文字列に置き換えてよい」
\end{itemize}

\subsection{実際に文字列を生成してみる}

\texttt{()}という文字列を生成する過程を追ってみましょう。

\begin{codescreen}{}
D               // 開始記号から始める
→ P            // D → P を適用
→ ( P ) P      // P → ( P ) P を適用
→ ( ε ) P     // 最初のPに P → ε を適用
→ ( ) P        // εは空文字列なので消える
→ ( ) ε       // 2番目のPに P → ε を適用  
→ ( )          // εは空文字列なので消える
\end{codescreen}

このように、規則を順番に適用して文字列を作ることを\textgt{導出}と呼びます。生成すると言い換えてもよいでしょう。

\subsection{複数の導出方法}

同じ文字列を導出する方法は複数あります。簡単な文法で考えてみましょう。

\begin{codescreen}{}
S → AB
A → a
B → b
\end{codescreen}

この文法から\texttt{ab}を導出するとき、どちらから先に展開するかで2通りの方法があります。

\begin{itemize}
\item
  方法1（左から展開）：
\begin{codescreen}{}
S → AB → aB → ab
\end{codescreen}
\item
  方法2（右から展開）：
\begin{codescreen}{}
S → AB → Ab → ab
\end{codescreen}
\end{itemize}

\subsection{最左導出と最右導出}

このように好きな順番で非終端記号を展開できるのですが、導出の過程を統一するために2つの標準的な方法が定義されています。

\begin{itemize}
\item
  \textgt{最左導出}：常に一番左の非終端記号を展開
\item
  \textgt{最右導出}：常に一番右の非終端記号を展開
\end{itemize}

簡単な例で見てみましょう。以下の文法は「1個以上のa」の後に「1個以上のb」が続く文字列を表します。

\begin{codescreen}{}
S → AB
A → aA | a
B → bB | b
\end{codescreen}

\texttt{aabb}を導出する場合

\begin{itemize}
\item
  \textgt{最左導出：}
\begin{codescreen}{}
S  
→ AB      (S → AB)
→ aAB     (A → aA)  
→ aaB     (A → a)
→ aabB    (B → bB)
→ aabb    (B → b)
\end{codescreen}

\item
  \textgt{最右導出：}
\begin{codescreen}{}
S
→ AB      (S → AB)
→ AbB     (B → bB)
→ Abb     (B → b)  
→ aAbb    (A → aA)
→ aabb    (A → a)
\end{codescreen}
\end{itemize}

\subsection{解析木との関係}

どちらの導出方法でも、最終的に同じ\textgt{解析木}が得られます。

\begin{center}
\begin{tikzpicture}[
  level distance=1.5cm,
  level 1/.style={sibling distance=3cm},
  level 2/.style={sibling distance=1.5cm},
  every node/.style={circle, draw, minimum size=8mm}
]
  \node {S}
    child {node {A}
      child {node[circle, draw=none] {a}}
      child {node {A}
        child {node[circle, draw=none] {a}}
      }
    }
    child {node {B}
      child {node[circle, draw=none] {b}}
      child {node {B}
        child {node[circle, draw=none] {b}}
      }
    };
\end{tikzpicture}
\end{center}

最左導出は木を左から右へ構築し、最右導出は右から左へ構築するイメージです。ここで得られた解析木は具体的な文字列の情報を含んでおり、抽象構文木と対比する形で具象構文木と呼ばれることもあります。

\subsection{なぜ2つの導出方法が重要か}

これらの導出方法は、構文解析の2つのアプローチと関連があります。

\begin{itemize}
\item
  \textgt{下向き構文解析}（トップダウン）：文法の開始記号から始めて、入力に合うよう規則を適用していく
\item
  \textgt{上向き構文解析}（ボトムアップ）：入力から始めて、段階的に規則を適用して開始記号まで到達する
\end{itemize}

第2章で作ったJSONパーサーは下向き構文解析の一種でした。ただし、導出と構文解析の対応関係は複雑で、最左導出と下向き構文解析が常に一対一で対応するわけではありません。次章では、これらの手法の実際の動作についてより詳しく学んでいきます。

\section{まとめ}

この章では、文脈自由文法について学びました。最初は難しく感じたかもしれませんが、実は私たちが普段書いているプログラムと深く関わっている概念です。この章で学んだことを整理します。

\begin{enumerate}
\item
  文脈自由文法の基礎

  \begin{itemize}
  \item
    Javaのif文やJSONの構造など、プログラミングの「入れ子構造」は文脈自由文法で表現される
  \item
    BNFから文脈自由文法への変換は記法の違いに過ぎない
  \item
    生成規則、非終端記号、終端記号という基本要素で構成される
  \end{itemize}
\item
  言語を集合として理解する

  \begin{itemize}
  \item
    プログラミング言語は「正しいプログラムの集合」として定義できる
  \item
    集合論の記法を使って言語間の関係を厳密に議論できる
  \item
    後方互換性なども集合の包含関係として表現可能
  \end{itemize}
\item
  正規表現とオートマトン

  \begin{itemize}
  \item
    正規表現は正規言語を表現するための強力なツール
  \item
    NFA（非決定性有限オートマトン）とDFA（決定性有限オートマトン）の違い
  \item
    正規表現は有限の状態しか持てないため、括弧の対応などの複雑な構造を扱えない
  \end{itemize}
\item
  言語の階層

  \begin{itemize}
  \item
    正規表現（正規言語）\textless{}
    文脈自由文法（文脈自由言語）\textless{} \ldots{} 帰納的可算言語
  \item
    プログラミング言語の構文解析には文脈自由文法が必要
  \end{itemize}
\item
  導出の仕組み

  \begin{itemize}
  \item
    生成規則を適用して文字列を生成する過程が導出
  \item
    最左導出と最右導出という2つの標準的な方法がある
  \item
    これらは下向き/上向き構文解析に対応する
  \end{itemize}
\end{enumerate}

文脈自由文法の理解は、以下の場面で役立ちます。

\begin{itemize}
\item
  \textgt{構文解析器の実装}：第4章以降で学ぶさまざまな構文解析アルゴリズムの基礎
\item
  \textgt{言語設計}：新しいDSLやプログラミング言語を設計する際の指針
\item
  \textgt{エラーメッセージの理解}：コンパイラのエラーメッセージがなぜそう言っているかの理解
\item
  \textgt{ツールの選択}：正規表現で十分か、パーサーが必要かの判断
\end{itemize}

次章ではこの文脈自由文法を元に、実際の構文解析アルゴリズムについて詳しく見ていきます。

\chapter{構文解析アルゴリズム古今東西}\label{ch04}
\index{こうぶんかいせきあるごりずむ@構文解析アルゴリズム}
\index{Dyck言語}
\index{ぶんみゃくじゆうぶんぽう@文脈自由文法}

第3章で学んだ文脈自由文法は、構文解析の理論的な基礎です。特にDyck言語を例に、括弧の対応という根本的な問題を通して、再帰的な構造を文法でどう表現するかを学びました。この章では、いよいよその文法を実際に解析するための「アルゴリズム」について深く掘り下げていきます。

\section{本章で学ぶこと}
\index{よそくがたしたむきこうぶんかいせき@予測型下向き構文解析}
\index{したむきこうぶんかいせき@下向き構文解析}
\index{うわむきこうぶんかいせき@上向き構文解析}
\index{しふとかんげんこうぶんかいせき@シフト還元構文解析}
\index{ばっくとらっく@バックトラック}
\index{PEG}
\index{Packrat Parsing}

「構文解析アルゴリズム」と聞くと難しそうに感じるかもしれません。しかし、実は皆さんはすでに2つの構文解析アルゴリズムを実装しているのです。

第2章を思い出してください。最初に実装した\texttt{PegJsonParser}は、\textbf{\sffamily PEG（Parsing
Expression
Grammar）}という手法の素朴な実装でした。次に実装した\texttt{SimpleJsonParser}は、実は\textbf{\sffamily LL(1)}（「エルエルワン」と読みます）と呼ばれる手法に近い\textgt{再帰下降構文解析器}だったのです。

この章では、これらのアルゴリズムがどのような仕組みで動いているのか、他にどのようなアルゴリズムが存在するのかを体系的に学びます。具体的には

\begin{enumerate}
\item
  \textgt{予測型下向き構文解析}\textbf{\sffamily （Predictive Top-down
  Parsing）}：文法の開始記号から始めて、入力文字列に向かって解析を進める方法

  \begin{itemize}
  \item
    先読みを使って適用する規則を予測する
  \item
    LL(1)
  \end{itemize}
\item
  \textgt{上向き構文解析}\textbf{\sffamily （Bottom-up / Shift-reduce
  Parsing）}：入力文字列から始めて、文法の開始記号に向かって解析を進める方法

  \begin{itemize}
    \item
    シフトと還元を使って解析を進める
  \item
    LR(0)、SLR(1)、LR(1)、LALR(1)
  \end{itemize}
\item
  \textgt{下向き構文解析+バックトラック}：先読みの代わりにバックトラックを使う方法

  \begin{itemize}
    \item
    PEG（Parsing Expression Grammar）、Packrat Parsing
  \end{itemize}
\end{enumerate}

それぞれのアルゴリズムには得意・不得意があり、実際のプログラミング言語の構文解析器では、言語の特性に応じて最適なアルゴリズムが選ばれています。

\subsection{本章の読み方ガイド}

本章は構文解析の主要アルゴリズムを網羅的に扱うため、内容がやや重くなっています。特にLR系（LR(0)、SLR(1)、LR(1)、LALR(1)）は、概念の積み重ねが多く、初読でつまずく読者が一定数います。そこで、本章をどのように読むかについて、二つのおすすめルートを示します。

まず構文解析の全体像をつかみたい方には、「下向きと上向きの2つの世界観」「LL(1)」「シフト還元構文解析の基本」「PEG」「Packrat Parsing」を必読として読み、LR(0)〜LALR(1)の詳細は後回しにすることをおすすめします。各節の冒頭に置いた「要点」と章末の比較表を眺めるだけでも、本章の主旨はおおむねつかめるはずです。LRの実装は構文解析器生成系（第5章で詳述）に任せるのが普通なので、ひとまずは「LRという上向き手法の系譜があり、Yacc/BisonはLALR(1)を使う」程度に押さえておけば次に進めます。

仕組みまで踏み込みたい方は、上から順番に読むのが自然です。LR(0)→SLR(1)→LR(1)→LALR(1)は、それぞれが前の手法の限界を克服する形で発展しているので、なぜ次の手法が必要だったのかを意識すると理解が深まります。

LR系の各節は、後で実務で必要になったときに参照しやすい構成にしてあります。

\begin{itemize}
\item Yaccでshift/reduce conflictのメッセージに遭遇した
\item 自分でパーサージェネレータの中身を読みたい
\end{itemize}

そんなときに、該当する節だけ読み返して理解できるように、各節を可能な限り独立した「リファレンス」として書いています。初読で全てを理解できなくても、まったく問題ありません。

\section{構文解析器生成系との関係}
\index{こうぶんかいせききせいせいけい@構文解析器生成系}
\index{ぱーさーじぇねれーた@パーサージェネレータ}
\index{Yacc}
\index{Bison}
\index{JavaCC}
\index{ANTLR}
\index{ALL(*)}
\index{LL(*)}

各アルゴリズムには、対応する\textgt{構文解析器生成系}（パーサージェネレータ）が存在することが多いです。以下、よく知られているものを列挙しますが、それ以外にも多数の構文解析器生成系が存在します。

\begin{itemize}
\item
  \textbf{\sffamily Yacc/Bison}（以降、Yaccと表記します）：

  \begin{itemize}
  \item
    LALR(1)を採用、C言語向け
  \end{itemize}
\item
  \textbf{\sffamily JavaCC}：

  \begin{itemize}
  \item
    LL(k)を採用（デフォルトはk=1）、 Java向け
  \end{itemize}
\item
  \textbf{\sffamily ANTLR}：

  \begin{itemize}
  \item
    \texttt{ALL(*)}\footnote{ANTLR
      4が採用する解析戦略。必要なだけ先読みを伸ばす\texttt{LL(*)}を動的に評価（Adaptive
      LL）する方式で、固定長先読みの\texttt{LL(k)}を実用的に一般化したものです。詳細は
      Terence Parr, Sam Harwell, Kathleen Fisher, ``Adaptive LL(*)
      Parsing: The Power of Dynamic Analysis'', OOPSLA '14, 2014
      を参照。\texttt{LL(*)} の元論文は Parr \& Fisher, ``LL(*): The
      Foundation of the ANTLR Parser Generator'', PLDI '11, 2011。}を採用（拡張されたLL系を採用）、 多言語対応
  \end{itemize}
\item
  \textgt{各種PEGパーサージェネレータ}：PEG\footnote{Bryan Ford, ``Parsing
    expression grammars: a recognition-based syntactic foundation'',
    POPL '04, 2004. PEGの形式的な定義と意味論を与えた原典です。}を採用
\end{itemize}

これらのツールについては第5章で詳しく説明しますが、本章でアルゴリズムを理解することで、各ツールがなぜそのアルゴリズムを選んだのかが見えてくるはずです。

\section{下向き構文解析と上向き構文解析 --- 2つの世界観}
\index{したむきこうぶんかいせき@下向き構文解析}
\index{うわむきこうぶんかいせき@上向き構文解析}
\index{Top-down parsing}
\index{Bottom-up parsing}

構文解析アルゴリズムの世界には、大きく分けて2つの「世界観」があります。\textgt{下向き}\textbf{\sffamily （Top-down）}と\textgt{上向き}\textbf{\sffamily （Bottom-up）}です。この2つは、同じ構文解析という仕事を、まったく逆のアプローチで実現します。

\subsection{なぜ2つのアプローチが必要なのか？}

第3章で学んだDyck言語を例に考えてみましょう。文法は以下でした。

\begin{lstlisting}
  D → P
  P → ( P ) P
  P → ε
\end{lstlisting}

入力文字列 \texttt{(())} を解析する際、2つの考え方があります。

\begin{enumerate}
\item
  \textgt{下向き（予測的）}：「Dから始まって、どうすれば\texttt{(())}が導出できるか？」
\item
  \textgt{上向き（還元的）}：「\texttt{(())}から始めて、どうすればDにたどり着けるか？」
\end{enumerate}

これは、迷路を解くときに「入口から出口を探す」か「出口から入口を探す」かの違いに似ています。どちらも正解にたどり着きますが、効率や適用できる問題が異なるのです。

\section{下向き構文解析の基本原理}

下向き構文解析は、文法の開始記号から出発し、入力文字列に向かって「上から下へ」解析を進めます。

\subsection{Dyck言語で学ぶ予測型下向き構文解析}

第3章で学んだDyck言語（括弧の対応が取れた文字列を表す言語）を例に、具体的に動作を追ってみましょう。説明のために、文法に入力の開始・終了を表す\texttt{\$}（ドル記号）を追加します。

\begin{lstlisting}
  D -> $ P $
  P -> ( P ) P
  P -> ε
\end{lstlisting}

この文法が、入力文字列 \texttt{(())} を受理するかどうかを判定してみましょう。

\subsection{スタックを使った解析の追跡}

ナイーヴな予測型下向き構文解析では、\textgt{スタック}を使って解析の進行状況を管理します。スタックには、「今どの規則のどの位置を解析しているか」という情報を記録します。

ドット記号「・」（中黒）を使って、規則内の現在位置を表すことにします。たとえば、以下のようになります。

\begin{itemize}
\item
  \texttt{D\ -\textgreater{}\ ・\$\ P\ \$} ：まだ何も解析していない状態
\item
  \texttt{D\ -\textgreater{}\ \$・P\ \$}
  ：\texttt{\$}を解析済み、次は\texttt{P}を解析する
\item
  \texttt{D\ -\textgreater{}\ \$\ P\ \$・} ：すべて解析完了
\end{itemize}

このドット記号は、Javaでいえばデバッガーでステップ実行する時の「現在の実行位置」のようなものだと考えてください。

では、実際に解析を開始しましょう。スタックは右が一番上で左に行くほど下になる、つまり最右がスタックトップであることに注意してください。

\begin{lstlisting}
  入力: $ ( ( ) ) $
  スタック: [ D -> ・$ P $ ]
\end{lstlisting}

\begin{enumerate}
\item
  最初の記号 \texttt{\$}
  を入力から読み込みます。スタックトップの規則が期待する記号と一致するので、ドットを進めます。

\begin{lstlisting}
  入力: ( ( ) ) $
  スタック: [ D -> $・P $ ]
\end{lstlisting}

\item
  非終端記号 \texttt{P}
  を解析する必要があります。ここで「予測」（あるいは先読み）、が必要になります。

非終端記号とは、第1章で学んだように「まだ展開の余地がある記号」のことでした。Javaでいえばメソッド呼び出しのようなものです。

次の入力文字（先読み文字）は \texttt{(}
です。\texttt{P}の規則は2つあります。

\begin{itemize}
\item
  \texttt{P\ -\textgreater{}\ (\ P\ )\ P} ：\texttt{(}で始まる
\item
  \texttt{P\ -\textgreater{}\ ε} ：空文字列（何もない）
\end{itemize}

ここで、どちらを選ぶかを事前に決める必要があります。ここでは、先読み文字が\texttt{(}の時は前者、それ以外のときは後者を選ぶこととします。

というわけで、先読み文字が\texttt{(}なので、\texttt{P\ -\textgreater{}\ (\ P\ )\ P}を選択します。この規則をスタックに追加します。

\begin{lstlisting}
  入力: ( ( ) ) $
  スタック: [D -> $・P $, P -> ・( P ) P]
\end{lstlisting}

\item
  入力の先頭の \texttt{(}
  を読み込み、スタックトップの規則のドットを進めます。

\begin{lstlisting}
  入力: ( ) ) $
  スタック: [D -> $・P $, P -> (・ P ) P]
\end{lstlisting}

\item
  非終端記号 \texttt{P} を解析するため、再び先読み文字を見ます。

次の先読み文字は \texttt{(} です。先ほどと同様に
\texttt{P\ -\textgreater{}\ (\ P\ )\ P}
を選択して、この規則をスタックに追加します。

\begin{lstlisting}
  入力: ( ) ) $
  スタック: [D -> $・P $, P -> ( ・ P ) P, P -> ・ ( P ) P]
\end{lstlisting}

\item
  入力の先頭の \texttt{(}
  を読み込み、スタックトップの規則のドットを進めます。

\begin{lstlisting}
  入力: ) ) $
  スタック: [D -> $・P $, P -> ( ・P ) P, P -> ( ・ P ) P]
\end{lstlisting}

\item
  非終端記号 \texttt{P} を解析するため、先読み文字を見ます。

次の先読み文字は \texttt{)} です。今度は、\texttt{P\ -\textgreater{}\ ε}
を選択します。空文字列にマッチするので、ドットを進めるだけで、入力からは何も消費しません。

\begin{lstlisting}
  入力: ) ) $
  スタック: [D -> $・P $, P -> ( ・P ) P, P -> ( P ・ ) P]
\end{lstlisting}

\item
  入力の先頭の \texttt{)}
  を読み込み、スタックトップの規則のドットを進めます。

\begin{lstlisting}
  入力: ) $
  スタック: [D -> $・P $, P -> ( ・P ) P, P -> ( P ) ・ P]
\end{lstlisting}

\item
  非終端記号 \texttt{P} を解析するため、先読み文字を見ます。

次の先読み文字は \texttt{)} です。ここで、\texttt{P\ -\textgreater{}\ ε}
を選択します。空文字列にマッチするので、ドットを進めるだけで、入力からは何も消費しません。

\begin{lstlisting}
  入力: ) $
  スタック: [D -> $・P $, P -> ( ・P ) P, P -> ( P ) P ・]
\end{lstlisting}

\item
  スタックトップの規則のドットが最後まで進んだので、スタックからその規則を除去すると同時に、新たなスタックトップのドットを進めます。ちょうどメソッド呼び出しが終わったのでそこから返って続きから再開するようなものです。

\begin{lstlisting}
  入力: ) $
  スタック: [D -> $・P $, P -> ( P ・) P]
\end{lstlisting}

\item
  入力の先頭の \texttt{)}
  を読み込み、スタックトップの規則のドットを進めます。

\begin{lstlisting}
  入力: $
  スタック: [D -> $・P $, P -> ( P ) ・ P]
\end{lstlisting}

\item
  非終端記号 \texttt{P} を解析するため、先読み文字を見ます。

次の先読み文字は \texttt{\$}
です。ここで、\texttt{P\ -\textgreater{}\ ε}
を選択します。空文字列にマッチするので、ドットを進めるだけで、入力からは何も消費しません。

\begin{lstlisting}
  入力: $
  スタック: [D -> $ ・ P $, P -> ( P ) P ・]
\end{lstlisting}

\item
  スタックトップの規則のドットが最後まで進んだので、スタックからその規則を除去すると同時に、新たなスタックトップのドットを進めます。

\begin{lstlisting}
  入力: $
  スタック: [D -> $ P ・ $]
\end{lstlisting}

\item
  入力の先頭の \texttt{\$}
  を読み込み、スタックトップの規則のドットを進めます。

\begin{lstlisting}
  入力:
  スタック: [D -> $ P $ ・]
\end{lstlisting}

\item
  ここでスタックトップの規則のドットが最後まで進んだので、スタックからその規則を除去します。スタックが空になり、入力もすべて消費されたので、受理されます。

\begin{lstlisting}
  入力:
  スタック: []
\end{lstlisting}
\end{enumerate}

このように、「先読み文字を見て、次に適用すべき規則を予測する」のが予測型下向き構文解析の特徴です。

\subsection{予測型下向き構文解析のアルゴリズム}

予測型下向き構文解析の動作パターンをまとめると以下のようになります。

\begin{enumerate}
\item
  先読み文字の取得：

  \begin{itemize}
  \item
    現在の入力位置から1文字先を見る
  \end{itemize}
\item
  非終端記号の展開：

  \begin{itemize}
  \item
    スタックトップの・の次が非終端記号の場合、先読み文字に基づいて適用する規則を「予測」
  \item
    選択した規則をスタックに追加
  \item
    適切な規則がない場合はエラー
  \end{itemize}
\item
  終端記号のマッチング：

  \begin{itemize}
  \item
    スタックトップの・の次が終端記号の場合、入力文字と比較
  \item
    一致すれば入力を消費してドットを進める
  \item
    一致しなければエラー
  \end{itemize}
\item
  規則の完了：

  \begin{itemize}
  \item
    規則の最後まで進んだら、スタックからその規則を除去
  \item
    一つ上の規則の解析を続行
  \end{itemize}
\item
  受理判定：

  \begin{itemize}
  \item
    入力がすべて消費され、スタックが空になれば受理
  \item
    それ以外は拒否
  \end{itemize}
\end{enumerate}

\section{予測型再帰下降構文解析 --- 下向き構文解析のJava実装}
\index{よそくがたさいきかこうぶんかいせき@予測型再帰下降構文解析}
\index{さいきかこうぶんかいせき@再帰下降構文解析}
\index{さいき@再帰}

スタックを使った下向き構文解析の動作を理解したところで、これをJavaで実装してみましょう。多くの場合、スタックを明示的に使わずに\textgt{再帰呼び出し}を使って実装できます。第2章で実装した\texttt{SimpleJsonParser}は、予測型再帰下降構文解析の一種で、実装もほぼ同じです。

\begin{codescreen}{}
// 文法規則：
// D -> P
// P -> ( P ) P
// P -> ε（イプシロン：空文字列）
public class Dyck {
    private final String input;
    private int pos;

    public Dyck(String input) {
        this.input = input;
        this.pos = 0;
    }

    public boolean parse() {
        boolean accepted = D();
        if(accepted && pos == input.length()) {
            return true;
        } else {
            return false;
        }
    }

    private boolean D() {
        return P();
    }

    private boolean P() {
        // 先読み文字が '(' の場合
        if (pos < input.length() && input.charAt(pos) == '(') {
            // P -> ( P ) P
            pos++; // '(' を読み進める
            if (!P()) return false;
            if (pos < input.length() && input.charAt(pos) == ')') {
                pos++; // ')' を読み進める
                return P();
            } else {
                return false;
            }
        } else {
            // P -> ε
            return true;
        }
    }
}
\end{codescreen}

\subsection{コードの解説}

この実装では、各非終端記号に対応するメソッドを作成しています。

\begin{itemize}
\item
  \texttt{D()}メソッド：文法の開始記号\texttt{D}に対応

  \begin{itemize}
  \item
    規則 \texttt{D\ -\textgreater{}\ P}
    をそのまま実装：\texttt{P()}を呼び出すだけ
  \end{itemize}
\item
  \texttt{P()}：非終端記号\texttt{P}に対応

  \begin{itemize}
  \item
    先読み文字が\texttt{(}の場合、 \texttt{P\ -\textgreater{}\ (\ P\ )\ P}
    を試す
  \item
    それ以外の場合は \texttt{P\ -\textgreater{}\ ε}
    を適用：常に\texttt{true}を返す
  \end{itemize}
\item
  文字の消費：\texttt{pos}をインクリメントすることで実現
\item
  バックトラックのない予測型：一度選択した規則で失敗したら、即座に\texttt{false}を返す
\end{itemize}

\subsection{文法規則とコードの対応関係}

BNF規則とJavaコードの対応を見てみましょう。

\begin{lstlisting}
  文法規則： D -> P
  Javaコード： private boolean D() { return P(); }

  文法規則： P -> ( P ) P | ε
  Javaコード： private boolean P() {
      if (先読み文字 == '(') {
          // P -> ( P ) P を適用
      } else {
          // P -> ε を適用
      }
  }
\end{lstlisting}

\begin{itemize}
\item
  各非終端記号にメソッドが対応
\item
  非終端記号の参照はメソッド呼び出しに変換
\item
  選択はif文で実現
\item
  連接は順次実行で実現
\end{itemize}

\noindent と対応しています。

\subsection{「再帰」「下降」という名前の由来}

この実装方法が「再帰下降」と呼ばれるのは以下の理由によるものです。

\begin{enumerate}
\item
  再帰：メソッドが自分自身または他のメソッドを呼び出す
\item
  下降：文法の開始記号から「下」の規則へと解析が進む
\end{enumerate}

実際に\texttt{parse("(())")}を実行すると、以下のような呼び出しの連鎖が発生します。

\begin{lstlisting}
  parse() -> D() -> P() -> P() -> P() -> ...
\end{lstlisting}

この呼び出しスタックが、先ほどスタックで説明した状態と対応しているのです。

\section{上向き構文解析の基本原理}
\index{うわむきこうぶんかいせき@上向き構文解析}
\index{しふとかんげんこうぶんかいせき@シフト還元構文解析}
\index{しふと@シフト}
\index{かんげん@還元}
\index{すたっく@スタック}

下向き構文解析を理解したところで、次は\textgt{上向き構文解析}を見ていきましょう。こちらは下向きとは全く逆のアプローチを取ります。

\subsection{シフト還元構文解析の基本アイデア}

ナイーヴな上向き構文解析は\textgt{シフト還元構文解析}とも呼ばれます。この名前は、2つの基本操作からきています。

\begin{itemize}
\item
  シフト（Shift）：入力から1文字読み込んでスタックに積む
\item
  還元（Reduce）：スタック上の記号列が規則の右辺と一致したら、左辺に置き換える
\end{itemize}

下向き構文解析が「文法から入力へ」と進むのに対し、上向き構文解析は「入力から文法へ」と進みます。入力文字列を徐々に「縮めて」いって、最終的に開始記号にたどり着けるかを確認するのです。

\subsection{シフト還元の基本的な動作}

シフト還元構文解析の基本的な流れは以下のとおりです。

\begin{enumerate}
\item
  入力をスタックに積んでいく（シフト操作）
\item
  \textgt{スタックの中身を見て、文法規則にマッチするパターンがあれば還元操作を行う}
\item
  \textgt{還元とシフトを繰り返し、最終的に開始記号だけが残れば成功}
\end{enumerate}

この単純な操作の繰り返しで、複雑な構文も解析できます。テトリスのブロックを積み上げていくようなイメージです。入力文字列を1文字ずつ積み上げ、規則にマッチしたらその部分を消すことで、最終的に開始記号だけが残るようにします。

重要なのは、「いつシフトして、いつ還元するか」の判断です。素朴な手法では、まず還元できるかを確認し、できなければシフトするという戦略をとります。

\subsection{例で学ぶシフト還元構文解析}

同じDyck言語を例にしますが、ナイーヴな上向き構文解析では空文字列規則（ε規則）が扱いづらいので、等価な別の文法を使います。

\begin{lstlisting}
  D -> $ P $
  D -> $ ε$
  P -> P X
  P -> X
  X -> ( P )
  X -> ()
\end{lstlisting}

この文法の特徴は以下のとおりです。

\begin{itemize}
\item
  \texttt{X}は内側の括弧ペアを表す
\item
  \texttt{P}は括弧ペアの列を表す
\item
  \texttt{D}は全体を表す
\end{itemize}

では、入力文字列 \texttt{(())}
に対してシフト還元構文解析を行ってみましょう。

\begin{enumerate}
\item
  開始記号\texttt{\$}をシフト

\begin{lstlisting}
  入力: ( ( ) ) $
  スタック: [ $ ]
\end{lstlisting}

\item
  最初の\texttt{(}をシフト

\begin{lstlisting}
  入力: ( ) ) $
  スタック: [ $, ( ]
\end{lstlisting}

\item
  さらに\texttt{(}をシフト

\begin{lstlisting}
  入力: ) ) $
  スタック: [ $, (, ( ]
\end{lstlisting}

\item
  さらに\texttt{)}をシフト

\begin{lstlisting}
  入力: ) $
  スタック: [ $, (, (, ) ]
\end{lstlisting}

\item
  スタックの末尾\texttt{(,\ )}が規則\texttt{X\ -\textgreater{}\ ()}の右辺と一致。2つの記号を\texttt{X}に\textgt{還元}

\begin{lstlisting}
  入力: ) $
  スタック: [ $, (, X ]
\end{lstlisting}

\item
  スタックの末尾\texttt{X}が規則\texttt{P\ -\textgreater{}\ X}の右辺と一致。1つの記号を\texttt{P}に\textgt{還元}（この実装はシフトする前に、可能な還元をすべて試みます）

\begin{lstlisting}
  入力: ) $
  スタック: [ $, (, P ]
\end{lstlisting}

\item
  次の\texttt{)}をシフト

\begin{lstlisting}
  入力: $
  スタック: [ $, (, P, ) ]
\end{lstlisting}

\item
  スタックの末尾\texttt{(,\ P,\ )}が規則\texttt{X\ -\textgreater{}\ (\ P\ )}の右辺と一致。3つの記号を\texttt{X}に\textgt{還元}

\begin{lstlisting}
  入力: $
  スタック: [ $, X ]
\end{lstlisting}

\item
  スタックの末尾\texttt{X}が規則\texttt{P\ -\textgreater{}\ X}の右辺と一致。1つの記号を\texttt{P}に\textgt{還元}

\begin{lstlisting}
  入力: $
  スタック: [ $, P ]
\end{lstlisting}

\item
  \texttt{\$}をシフト

\begin{lstlisting}
  入力:
  スタック: [ $, P, $ ]
\end{lstlisting}

\item
  スタックの末尾\texttt{\$,\ P,\ \$}が規則\texttt{D\ -\textgreater{}\ \$\ P\ \$}の右辺と一致。3つの記号を\texttt{D}に\textgt{還元}

\begin{lstlisting}
  入力:
  スタック: [ D ]
\end{lstlisting}
\end{enumerate}

入力が空になり、スタックに開始記号\texttt{D}だけが残ったので、受理されます。

\subsection{シフト還元構文解析のアルゴリズム}

シフト還元構文解析の動作パターンをまとめると

\begin{enumerate}
\item
  シフト操作：

\begin{itemize}
\item
  入力から1文字読み込んでスタックに積む
\item
  常に左から右へ順番に読む
\end{itemize}

\item
  還元操作：

\begin{itemize}
\item
  スタック上端の記号列が、ある規則の右辺と一致したら
\item
  その記号列を取り除いて、規則の左辺の非終端記号を積む
\end{itemize}

\item
  アクションの選択：

\begin{itemize}
\item
  シフトと還元のどちらを行うかを決定する必要がある
\item
  この決定方法がLR(0)、SLR(1)、LR(1)などの違いになる
\end{itemize}

\item
  受理判定：

\begin{itemize}
\item
  入力をすべて読み、スタックが開始記号だけになれば受理
\item
  それ以外は拒否
\end{itemize}
\end{enumerate}

\subsection{下向きと上向きの違い}

ここまでの例からわかるように

\begin{itemize}
\item
  \textgt{下向き}：\texttt{D}から始めて\texttt{(())}を「生成」しようとする
\item
  \textgt{上向き}：\texttt{(())}から始めて\texttt{D}に「還元」しようとする
\end{itemize}

上向き構文解析の還元操作は、第3章で学んだ「最右導出」の逆操作に相当します。つまり、導出の過程を逆再生しているのです。

\section{シフト還元構文解析のJava実装}

シフト還元構文解析をJavaで実装してみましょう。予測型下向き構文解析とは異なり、明示的にスタックを使い、規則をデータ構造として扱います。

\subsection{必要なデータ構造}

まず、記号（終端記号・非終端記号）を表すクラスと、文法規則を表すクラスが必要です。

\begin{codescreen}{}
// 記号を表すインターフェース
interface Element {
    char value();
}

// 終端記号（実際の文字：'(', ')', '$'など）
record Terminal(char value) implements Element {}

// 非終端記号（文法記号：'D', 'P', 'X'など）
record NonTerminal(char value) implements Element {}
\end{codescreen}

次に、文法規則を表す\texttt{Rule}クラスを定義します。

\begin{codescreen}{}
import java.util.List;
import java.util.ArrayList;

public record Rule(char lhs, List<Element> rhs) {
    // 可変長引数コンストラクタ（便利のため）
    public Rule(char lhs, Element... rhs) {
        this(lhs, List.of(rhs));
    }
    
    // スタックの上端がこの規則の右辺と一致するか判定
    public boolean matches(List<Element> stack) {
        if (stack.size() < rhs.size()) return false;
        
        // スタックの上からrhs.size()個の要素を比較
        for (int i = 0; i < rhs.size(); i++) {
            Element elementInRule = rhs.get(i);
            Element elementInStack = stack.get(
                stack.size() - rhs.size() + i
            );
            if (!elementInRule.equals(elementInStack)) {
                return false;
            }
        }
        return true;
    }
}
\end{codescreen}

\subsection{シフト還元構文解析器の実装}

上記のデータ構造を使って、実際のシフト還元構文解析器を実装します。

\begin{codescreen}{}
import java.util.List;
import java.util.ArrayList;

public class DyckShiftReduce {
    private final String input;
    private int position;
    private final List<Rule> rules;
    private final List<Element> stack = new ArrayList<>();

    public DyckShiftReduce(String input) {
        this.input = input;
        this.position = 0;
        
        // 文法規則の定義
        this.rules = List.of(
            // D -> $ P $
            new Rule('D', 
                new Terminal('$'), 
                new NonTerminal('P'), 
                new Terminal('$')
            ),
            // D -> $ $ (空文字列の場合)
            new Rule('D', 
                new Terminal('$'), 
                new Terminal('$')
            ),
            // P -> P X
            new Rule('P', 
                new NonTerminal('P'), 
                new NonTerminal('X')
            ),
            // P -> X
            new Rule('P', 
                new NonTerminal('X')
            ),
            // X -> ( P )
            new Rule('X', 
                new Terminal('('), 
                new NonTerminal('P'), 
                new Terminal(')')
            ),
            // X -> ()
            new Rule('X', 
                new Terminal('('), 
                new Terminal(')')
            )
        );
    }

    public boolean parse() {
        // 開始記号$をスタックに積む
        stack.add(new Terminal('$'));
        
        // メインループ：シフトと還元を繰り返す
        while (true) {
            // まず還元を試みる
            if (!tryReduce()) {
                // 還元できない場合、シフトを試みる
                if (position < input.length()) {
                    char c = input.charAt(position);
                    stack.add(new Terminal(c));
                    position++;
                } else {
                    // シフトもできないので終了
                    break;
                }
            }
        }
        
        // 終端記号$をスタックに積む
        stack.add(new Terminal('$'));
        
        // 最後に可能な限り還元を繰り返す
        while (tryReduce()) {
            // 還元ができなくなるまで続ける
        }
        
        // スタックが[D]のみになったら受理
        return stack.size() == 1 && 
               stack.get(0).equals(new NonTerminal('D'));
    }

    private boolean tryReduce() {
        for (Rule rule : rules) {
            if (rule.matches(stack)) {
                // マッチしたら右辺の長さ分スタックから削除
                for (int i = 0; i < rule.rhs().size(); i++) {
                    stack.remove(stack.size() - 1);
                }
                // 左辺の非終端記号をスタックに追加
                stack.add(new NonTerminal(rule.lhs()));
                return true; // 還元成功
            }
        }
        return false; // 還元できる規則がなかった
    }
}
\end{codescreen}

\subsection{実装のポイント}

この実装の重要な点は以下のとおりです。

\begin{enumerate}
\item
  シフトより還元を優先：
  まず\texttt{tryReduce()}で還元を試み、できない場合のみシフト
\item
  単純な還元ルール：
  すべての規則を順番にチェックし、最初にマッチしたものを使う
\item
  明示的なスタック操作：
  \texttt{List\textless{}Element\textgreater{}}をスタックとして使い、要素の追加・削除を明示的に行う
\end{enumerate}

この実装は最も単純なシフト還元構文解析で、実用的なパーサーでは洗練されたアルゴリズム（LR(0)、SLR(1)、LR(1)、LALR(1)）が使われます。

\section{下向き構文解析と上向き構文解析の比較}

ここまで下向き構文解析と上向き構文解析の具体例を見てきました。両者にはそれぞれ得意不得意があります。

\subsection{下向き構文解析の利点と欠点}

\textgt{利点：}

\begin{enumerate}
\item
  直感的な実装

\begin{itemize}
\item
  文法規則とメソッドが1対1に対応
\item
  手書きの構文解析器が書きやすい
\item
  多くのプログラミング言語のコンパイラが手書き再帰下降を採用
\end{itemize}

\item
  文脈依存の扱いやすさ

\begin{itemize}
\item
  メソッドの引数で情報を渡せる
\item
  解析中に状態を管理しやすい
\item
  エラー回復やエラーメッセージの生成が容易
\end{itemize}

\item
  デバッグのしやすさ

\begin{itemize}
\item
  通常の関数呼び出しなのでデバッガで追える
\item
  構文解析の過程が理解しやすい
\end{itemize}
\end{enumerate}

\textgt{欠点：}

\begin{enumerate}
\item
  左再帰の問題

たとえば、以下のBNFは上向き型だと普通に解析できますが、工夫なしに下向き型で実装すると無限再帰に陥ってスタックオーバーフローします。

\begin{lstlisting}
  A -> A "a" | ε // 左再帰を含む文法の例 (εは空文字列)
\end{lstlisting}

この文法は「\texttt{A}は、\texttt{A}の後に\texttt{a}が続くか、何もないか」という意味ですが、\texttt{A}を解析するためにまず\texttt{A}を解析する必要があるので無限ループになります。

このような問題を下向き型で解決する方法も存在します。たとえば、以下のように等価な右再帰の文法に書き換えることで除去できます。

\begin{lstlisting}
  A -> "a" A | ε
\end{lstlisting}

この変換により、下向き構文解析で問題となる無限再帰を避けられます。ただし、文法の書き換えは常に簡単とは限りません。
\end{enumerate}

\subsection{上向き構文解析の利点と欠点}

\textgt{利点：}

\begin{enumerate}
\item
  左再帰を自然に扱える
\item
  より広いクラスの文法を扱える
\item
  効率的な構文解析表による高速化が可能
\end{enumerate}

\textgt{欠点：}

\begin{enumerate}
\item
  実装が複雑
\item
  文脈依存の処理が難しい
\end{enumerate}

たとえば、それまでの文脈に応じて構文解析のルールを切り替えたくなることがあります。最近の言語によく搭載されている文字列補間などはその最たる例です。

たとえば、\texttt{"}の中は文字列リテラルとして特別扱いされますが、その中で\texttt{\#\{}が出てきたら（Rubyの場合）、通常の式を構文解析するときのルールに戻る必要があります。

このように、文脈に応じて適用するルールを切り替えるのは下向き型が得意です。もちろん、上向き型でも実現できないわけではありません。実際、Rubyの構文解析器はYaccの定義ファイルから生成されるようになっていますが、Yaccが採用しているのは代表的な上向き構文解析法である\texttt{LALR(1)}です。

\section{LL(1) --- 代表的な下向き構文解析アルゴリズム}
\index{LL(1)}
\index{えるえるほう@LL法}
\index{さきよみ@先読み}
\index{FIRST集合}
\index{FOLLOW集合}
\index{nullable}
\index{きょうつうせっとうじ@共通前置辞}
\index{ひだりさいき@左再帰}

ここからは、具体的な構文解析アルゴリズムについて詳しく見ていきましょう。まずは下向き構文解析の代表格である\textbf{\sffamily LL(1)}から始めます。

\subsection{LL(1)とは？}

\textbf{\sffamily LL(1)}という名前は以下の意味を持ちます：

\begin{enumerate}
\item
  最初のL：\textbf{\sffamily L}eft-to-right（左から右へ入力を読む）
\item
  2番目のL：\textbf{\sffamily L}eftmost derivation（最左導出を行う）
\item
  (1)：1トークン先読み
\end{enumerate}

つまり「1トークン先を見て、次に適用すべき規則を一意に決定できる」という制約を持つ下向き構文解析法です。

この「LL」という名前の由来ですが、構文解析の研究が盛んだった1960年代に、さまざまな手法を分類するために考案された命名規則です。同様に、後で出てくる「LR」も「Left-to-right,
Rightmost derivation」を表します。

\subsection{LL(1)の直感的な理解}

身近な例で考えてみましょう。Javaのif文とwhile文の解析を例にします。

\begin{codescreen}{}
if (age < 18) {
    System.out.println("18歳未満です");
}

while (count > 0) {
    count--;
}
\end{codescreen}

プログラマがこのコードを見たとき、最初のトークンだけで文の種類を判断できます。

\begin{itemize}
  \item ifで始まる → if文だ！
  \item whileで始まる → while文だ！
\end{itemize}

LL(1)は、まさにこの「最初の1トークンで判断」というアイデアをアルゴリズム化したものです。

しかし、すべての構文がこのように単純ではありません。たとえば、以下の単純な算術式を考えてみましょう。\texttt{N}は数値をあらわす終端記号とします。

\begin{lstlisting}
  E  → N
  E  → ( E )
  E  → - E
\end{lstlisting}

この場合、\texttt{E}は以下のいずれかで始まる可能性があります。

\begin{itemize}
\item
  \texttt{N}（例：\texttt{42}）
\item
  左括弧（例：\texttt{(3\ +\ 5)}）
\item
  マイナス記号（例：\texttt{-10}）
\end{itemize}

\texttt{E}は複数のトークンで始まり得ますが、それ自体は問題ではありません。困るのは、\textgt{同じトークンで始まる規則が複数ある場合に、最初の1トークンだけではどの規則を選ぶか決められない}ことです。

たとえば、

\begin{lstlisting}
  E → id
  E → id + E
\end{lstlisting}

\noindent のように、\texttt{id}が式の先頭として再利用される文法では、先読み1トークンだけではどちらの規則を適用するのか判断できません。トークンが複数の構文で使い回されると、このような曖昧さが生じます。ここで触れた「1トークンで判断できない」状況は、後で扱う共通前置辞問題（4.9.7）の最も短い形だと考えてください。

\subsection{なぜFIRST集合とFOLLOW集合が必要か}

LL(1)を実現するには、以下の2つの問題を解決する必要があります。

\begin{enumerate}
\item
  複数のトークンで始まる構文

上の算術式の例のように、一つの非終端記号が複数のトークンで始まる場合があります。これを体系的に扱うために\textbf{\sffamily FIRST}\textgt{集合}（その非終端記号から始まる可能性のあるトークンの集合）が必要です。

\item
  省略可能な要素

Javaのif文には、else節があるものとないものがあります。

\begin{codescreen}{}
// else節なし
if (condition) { ... }

// else節あり
if (condition) { ... } else { ... }
\end{codescreen}

else節は「あってもなくてもいい」要素です。このような場合、else節の後に何が来る「可能性がある」かを知る必要があります。これが\textbf{\sffamily FOLLOW}\textgt{集合}（その非終端記号の後に来る可能性のあるトークンの集合）です。
\end{enumerate}

\subsection{FIRST集合とFOLLOW集合}

LL(1)構文解析を実現するには、これらの集合を計算して利用します。

\subsubsection{FIRST集合}

ある非終端記号から導出される文字列の「最初に現れうるトークンの集合」を\textbf{\sffamily FIRST}\textgt{集合}と呼びます。

正式には、非終端記号\texttt{A}に対して以下のように定義されます。

\begin{lstlisting}
FIRST(A) = { a | A ⇒* aβ であるような終端記号 a } ∪ { ε | A ⇒* ε }
\end{lstlisting}

\texttt{A\ ⇒*\ aβ}は、非終端記号\texttt{A}が何らかの導出を経て、終端記号\texttt{a}で始まる文字列に変換できることを意味します。

たとえば、先ほどでてきた以下の算術式の文法を考えます。

\begin{lstlisting}
  E  → N
  E  → ( E )
  E  → - E
\end{lstlisting}

この場合、FIRST集合は次のようになります。

\begin{lstlisting}
  FIRST(E) = { '(', '-', N }
\end{lstlisting}

\subsubsection{nullable --- 空文字列を生成できるか}

LL(1)構文解析で重要な概念として、\textbf{\sffamily nullable}（ヌラブル）があります。これは、ある非終端記号が空文字列εを生成できるかどうかを示すものです。

\textbf{\sffamily nullable}という名前はJavaの\texttt{null}とは関係ありません。「空にできる」つまり「何も生成しないことができる」という意味です。

先ほどと同じ以下の文法を考えます。

\begin{lstlisting}
  E  → N
  E  → ( E )
  E  → - E
\end{lstlisting}

この算術式文法でのnullable集合は次のようになります。

\begin{lstlisting}
  nullable(E) = false
\end{lstlisting}

Eは明らかに空文字列を生成しないので\texttt{nullable(E)\ =\ false}となります。

ところで、なぜ\texttt{nullable}が重要かというと、\textbf{\sffamily FIRST}\textgt{集合}及び\textbf{\sffamily FOLLOW}\textgt{集合}を計算するアルゴリズムで必要になるためです。\textbf{\sffamily FIRST}\textgt{集合}の計算では、非終端記号が空文字列を生成できるかどうかを確認する必要があります。\textbf{\sffamily FOLLOW}\textgt{集合}の計算でも、非終端記号が空文字列を生成できる場合、その後に来るトークンを正しく扱うために\texttt{nullable}が必要です。

\subsubsection{FOLLOW集合}

非終端記号の後に現れうるトークンの集合を\textbf{\sffamily FOLLOW}\textgt{集合}と呼びます。これは、\textbf{\sffamily nullable}な非終端記号（空文字列を生成できる非終端記号）を扱うために特に重要です。

先ほどと同様に以下の文法を考えてみます。

\begin{lstlisting}
  E  → N
  E  → ( E )
  E  → - E
\end{lstlisting}

この場合\texttt{FOLLOW}集合は次のようになります。

\begin{lstlisting}
  FOLLOW(E) = { '$', ')' }
\end{lstlisting}

Eに続く可能性のあるトークンは、入力の終端記号\texttt{\$}（入力の終わり）と右括弧\texttt{)}です。

\subsection{FIRST集合、FOLLOW集合、nullableの計算方法}

これまでは概要でしたが、これらの集合を「実際に」計算するアルゴリズムを見ていきましょう。

\subsubsection{nullableの計算アルゴリズム}

\begin{enumerate}
\item
  全ての非終端記号Eについて\texttt{nullable(E)\ =\ false}に初期化
\item
  \(A \to \varepsilon\) の形の規則があれば\texttt{nullable(A)\ =\ true}
\item
  \(A \to X_1\,X_2\,\ldots\,X_k\) について、すべての \(X_i\)
  が\texttt{nullable}なら\texttt{nullable(A)\ =\ true}
\item
        \texttt{nullable = false ->
        nullable = true}への変化がなくなるまで繰り返す
\end{enumerate}

ここで \(X_1, X_2, \ldots\)
は、規則の右辺に並ぶそれぞれの記号（位置付きで区別したもの）を表します。

このような形になります。算術式文法を使って計算してみましょう。

\begin{lstlisting}
  E → N
  E → ( E )
  E → - E
\end{lstlisting}

計算過程：

\begin{itemize}
\item
  初期化：すべて\texttt{false}
\item
  \texttt{E\ →\ N}：Nは終端記号なので、Eはnullableではない
\item
  \texttt{E\ →\ (\ E\ )}：
  3つの記号があり、どれもnullableではないので、Eはnullableではない
\item
  \texttt{E\ →\ -\ E}：
  2つの記号があり、どちらもnullableではないので、Eはnullableではない
\item
  どの規則もεを生成せず、右辺もnullableでないので\texttt{false}のまま
\end{itemize}

結果：

\begin{itemize}
\item
  \texttt{nullable(E)\ =\ false}
\end{itemize}

Eは常に何らかの記号を生成するため、空文字列になることはありません。

\subsubsection{FIRST集合の計算アルゴリズム}

FIRST集合全体を次の手順で計算します。

\begin{enumerate}
\item
  すべての非終端記号の\(\mathrm{FIRST}\)を空集合に初期化
\item
  各規則 \(A \to \alpha\) について、\(\alpha = X_1\,X_2\,\ldots\,X_n\) とする
\item
  \(X_1\) が終端記号なら、\(\mathrm{FIRST}(A)\) に \(X_1\) を追加
\item
  \(X_1\) が非終端記号なら

  \begin{itemize}
  \item
    \(\mathrm{FIRST}(X_1) \mathbin{\backslash} \{\varepsilon\}\) を
    \(\mathrm{FIRST}(A)\) に追加
  \item
    \(X_1\) が\textbf{\sffamily nullable}なら \(X_2\) も確認（以下同様）
  \end{itemize}
\item
  すべての \(X_i\) が\textbf{\sffamily nullable}なら、\(\varepsilon\) を
  \(\mathrm{FIRST}(A)\) に追加
\item
  すべての規則を処理しても集合に変化がなくなるまで繰り返す
\end{enumerate}

ここで \(X_1, X_2, \ldots, X_n\)
は、規則の右辺に並ぶそれぞれの記号を位置付きで区別するための書き方です（同じ非終端記号でも、出現位置ごとに別の
\(X_i\) として扱います）。

同様に算術式文法で計算してみましょう。

\begin{lstlisting}
  E → N
  E → ( E )
  E → - E
\end{lstlisting}

各規則のFIRST集合を計算します。

\begin{itemize}
\item
  \texttt{E\ →\ N}：
  第1記号がN（終端記号）なので、\texttt{N}をFIRST(E)に追加
\item
  \texttt{E\ →\ (\ E\ )}：
  第1記号が\texttt{(}（終端記号）なので、\texttt{(}をFIRST(E)に追加
\item
  \texttt{E\ →\ -\ E}：
  第1記号が\texttt{-}（終端記号）なので、\texttt{-}をFIRST(E)に追加
\end{itemize}

\textgt{結果：}
\texttt{FIRST(E)\ =\ \{N,\ \textquotesingle{}(\textquotesingle{},\ \textquotesingle{}-\textquotesingle{}\}}

\subsubsection{FOLLOW集合の計算アルゴリズム}

\begin{enumerate}
\item
  すべての非終端記号の\texttt{FOLLOW}を空集合に初期化
\item
  開始記号\texttt{S}に対して\texttt{FOLLOW(S)\ =\ \{\$\}}（入力終端）
\item
  規則\texttt{B\ →\ αAβ}に対して

  \begin{itemize}
  \item
    \texttt{FIRST(β)\ -\ \{ε\}}を\texttt{FOLLOW(A)}に追加
  \item
    \texttt{β}がnullableなら\texttt{FOLLOW(B)}を\texttt{FOLLOW(A)}に追加
  \end{itemize}
\item
  規則\texttt{B\ →\ αA}（末尾に\texttt{A}）に対して

  \begin{itemize}
  \item
    \texttt{FOLLOW(B)}を\texttt{FOLLOW(A)}に追加
  \end{itemize}
\item
  変化がなくなるまで繰り返す
\end{enumerate}

同様に算術式文法で計算してみましょう。

\begin{lstlisting}
  E → N
  E → ( E )
  E → - E
\end{lstlisting}

初期状態：

\begin{lstlisting}
- FOLLOW(E) = {$}（開始記号として）
\end{lstlisting}

規則 \texttt{E\ →\ (\ E\ )} から：

\begin{lstlisting}
- Eの後に)が来るので、)をFOLLOW(E)に追加
- FOLLOW(E) = {$, ')'}
\end{lstlisting}

規則 \texttt{E\ →\ -\ E} から：

\begin{lstlisting}
- Eが末尾にあるので、FOLLOW(E)をFOLLOW(E)に追加（変化なし）
\end{lstlisting}

規則 \texttt{E\ →\ (\ E\ )} から（2回目）：

\begin{lstlisting}
- すでに ) は追加済みなので、変化なし
\end{lstlisting}

規則 \texttt{E\ →\ -\ E} から（2回目）：

\begin{lstlisting}
- 変化なし。すべての規則を見て変化がなくなったので終了
\end{lstlisting}

\textbf{結果：}
\texttt{FOLLOW(E)\ =\ \{\$,\ \textquotesingle{})\textquotesingle{}\}}

Eが唯一の非終端記号なのでFOLLOW集合の計算も簡単です。

\subsection{LL(1)構文解析でのFIRST/FOLLOW活用}

LL(1)構文解析の核心は、FIRST集合とFOLLOW集合を使って「どの規則を適用すべきか」を効率的に決定することです。

\subsubsection{FIRST/FOLLOWを使った規則選択の基本原理}

LL(1)構文解析では、現在解析中の非終端記号と次の入力トークンの組み合わせから、適用する規則を一意に決定します。

\textbf{\sffamily 1. FIRST}\textgt{集合による判断} --- ある規則 \texttt{A\ →\ α}について

\begin{itemize}
\item \texttt{α} が特定のトークン集合で始まる可能性がある場合、先読みトークン\texttt{a}が\texttt{FIRST(α)}に含まれている場合、この規則を適用する
\end{itemize}

\textbf{\sffamily 2. FOLLOW}\textgt{集合による判断} --- 規則 \texttt{A\ →\ α}の右辺がnullableな場合

\begin{itemize}
\item 先読みトークンがFOLLOW(A) に含まれる場合、この規則を適用する
\end{itemize}

\textbf{\sffamily 3. }\textgt{決定性の保証} --- 各状況で適用可能な規則が

\begin{itemize}
\item 一意に決まる場合、その文法は LL(1) である
\item 複数の規則が適用可能な場合、LL(1)コンフリクトが発生する
\end{itemize}

以上がLL(1)構文解析の基本的な考え方です。これにより、構文解析器は先読みトークンを1つだけ見て、次に適用すべき規則を決定できます。

\subsubsection{実用的な実装での最適化}

実際のLL(1)パーサーを構文解析器生成系で生成する場合は、上記の原理を使って構文解析表を事前に作成しておきます。これにより、解析時は単純な表引きだけで規則を決定できるため、高速な解析が可能になります。

具体的には、「非終端記号と先読みトークンの組み合わせ → 適用する規則」という対応表を作成し、解析器はこの表を参照するだけで動作します。

このようにして構築された解析表により、LL(1)パーサーは効率的かつ決定的に動作します。

\subsection{LL(1)の問題点と限界}

LL(1)はシンプルで実用的ですが、いくつかの限界があります。

\begin{enumerate}
\item
  共通前置辞問題

LL(1)では、同じ非終端記号から始まる複数の規則がある場合、先読みトークンだけではどの規則を適用すべきか判断できません。

\begin{lstlisting}
  S -> a B
  S -> a C
\end{lstlisting}

両方の規則が\texttt{a}で始まるため、1文字先読みでは判断できません。この問題は\textgt{左因子化}（Left
Factoring）で解決できます。

\begin{lstlisting}
  S  -> a S'
  S' -> B
  S' -> C
\end{lstlisting}

左因子化は直観的には、共通の前置きを抽出して新しい非終端記号を導入することで、先読みトークンだけで規則を一意に決定できるようにする変換操作です。

とはいえ、左因子化は万能ではありません。たとえば偶数長の回文を生成する文法

\begin{lstlisting}
  A → a A a | b A b | ε
\end{lstlisting}

\noindent は、どれだけ変形しても先頭の文字だけでは分岐できず、LL(1)にはなりません。左因子化をほどこしても、LL(1)に収まらない文法が存在する、という具体例です。また、左因子化を重ねると非終端記号が増えて抽象構文木との対応が見えにくくなり、規則の直感的な形を失いやすいという副作用もあります。

\item
  左再帰の問題

LL(1)の別の欠点は、\textgt{左再帰}を扱えないことです。たとえば、以下のような左再帰文法を扱うことができません。

\begin{lstlisting}
  E → E + T
  E → T
  T → F
  F → ( E ) | id
\end{lstlisting}

この問題は\textgt{左再帰の除去}で解決できます。

\begin{lstlisting}
  E → T E'
  E' → + T E' | ε
  T → F
  F → ( E ) | id
\end{lstlisting}

ただし、この変換により文法（規則の形）の直感性が失われやすく、抽象構文木の構築も複雑になります。
\end{enumerate}

ここまでで下向き構文解析の基本的な考え方とLL(1)について学びました。次は、より効率的な上向き構文解析の代表格である\textbf{\sffamily LR}\textgt{構文解析}を見ていきましょう。

\section{LR構文解析 --- 素朴なシフト還元の効率化}
\index{LR構文解析}
\index{えるあーるほう@LR法}
\index{LR(0)}
\index{SLR(1)}
\index{LR(1)}
\index{LALR(1)}
\index{LR(0)オートマトン}
\index{ゆうげんおーとまとん@有限オートマトン}

\begin{quote}
要点：
素朴なシフト還元では「いつシフト／いつ還元／どの規則で還元するか」を毎回スタック全体を見て決めていました。LR構文解析は、この判断を事前に構築した有限オートマトン（LR(0)オートマトン）に置き換えることで効率化したものです。本節以降ではLR(0)・SLR(1)・LR(1)・LALR(1)の順に発展していきますが、現在の実用パーサージェネレータ（Yacc/Bison）が採用しているのはLALR(1)で、初読では「LR系はLALR(1)が事実上の標準」と押さえれば次に進めます。
\end{quote}

最初に見た素朴なシフト還元構文解析には、大きな問題がありました。

\begin{enumerate}
\item
  \textgt{いつシフトして、いつ還元するか？} ---
  毎回スタック全体を調べる必要がある
\item
  \textgt{どの規則で還元するか？} --- すべての規則を順番に試す必要がある
\end{enumerate}

これらの判断を効率的に行うのが\textbf{\sffamily LR}\textgt{構文解析}です。

\subsection{素朴な方法の問題点}

先ほど（「例で学ぶシフト還元構文解析」（4.6.3節）で示した）\texttt{DyckShiftReduce}の実装を思い出してください。

\begin{codescreen}{}
private boolean tryReduce() {
    for (Rule rule : rules) {  // すべての規則を試す！
        if (rule.matches(stack)) {  // スタック全体をチェック！
            // 還元処理
        }
    }
    return false;
}
\end{codescreen}

入力が長くなると、この方法は非効率的です。LR構文解析は、この問題を解決します。

\subsection{LR構文解析のアイデア}

素朴なシフト還元構文解析の問題点は、毎回スタック全体を見て、どの規則が適用できるかを総当たりで確認する必要があることでした。しかし、ここで重要な観察があります。

スタックの内容によって「次に何ができるか」は\textgt{決まっている}ということです。

たとえば、Dyck言語の解析で考えてみましょう。

\begin{lstlisting}
  スタック: [ '(' ]
  次の入力: )
\end{lstlisting}

この状態では、次にできることは以下です。

\begin{enumerate}
\item
  \texttt{)}をシフトしてスタックに積む
\item
  その後、\texttt{(\ )}をXに還元する
\end{enumerate}

別の例を考えてみましょう。

\begin{lstlisting}
  スタック: [(, (, X  ]
  次の入力: )
\end{lstlisting}

還元を優先する\texttt{DyckShiftReduce}では、次のように進みます。

\begin{enumerate}
\item
  スタック末尾の\texttt{X}を\texttt{P}に還元する
\item
  \texttt{)}をシフトしてスタックに積む
\item
  その後、\texttt{(\ P\ )}を\texttt{X}に還元する
\end{enumerate}

\textgt{今どの状態にいるか}が分かれば、\textgt{次の行動が決まる}のです。

\subsection{状態の発見}

では、「状態」とは何でしょうか？実は、「今どの規則をどこまで読んでいるか」を表す情報です。

たとえば、規則 \texttt{X\ →\ (\ P\ )} について考えてみます。

\begin{enumerate}
\item
  \textgt{「(を読んだ直後の状態」}

  \begin{itemize}
  \item
    次は\texttt{P}を読みたい
  \item
    \texttt{P}は\texttt{ε}か\texttt{X\ P}のどちらか
  \end{itemize}
\item
  \textgt{「(}\textbf{\sffamily P}\textgt{まで読んだ状態」}

  \begin{itemize}
  \item
    次は\texttt{)}を読みたい
  \item
    \texttt{)}が来たらシフト
  \end{itemize}
\item
  \textgt{「}\textbf{\sffamily ( P )}\textgt{全部読んだ状態」}

  \begin{itemize}
  \item
    \texttt{X\ →\ (\ P\ )} で還元できる！
  \end{itemize}
\end{enumerate}

このように、各状態で「次に何を期待しているか」「何ができるか」が明確です。

\subsection{状態遷移の規則性}

さらに重要なことは、この「状態」は規則的に遷移するということです。

\begin{lstlisting}
  状態「(を読んだ直後」 + 記号「(」 → 状態「(を読んだ直後」
  状態「(を読んだ直後」 + 記号「)」 → 状態「( ) を読んだ」
\end{lstlisting}

この状態遷移は決定的で、事前にすべて計算できます。3章で学んだことを思い出してください。このような状態遷移はオートマトンとして表現できるのです。

\subsection{オートマトンの必然性}

状態と状態遷移を定義することで、素朴なシフト還元構文解析でやっていた流れをオートマトンとして表現できます。たとえば、以下はDyck言語の文法

\begin{lstlisting}
  D → P
  P → ( P ) P
  P → ε
\end{lstlisting}

\noindent に対して構築された簡略化されたオートマトンです。

前節ではシフト還元構文解析のために\texttt{\$}や\texttt{X}を導入した拡張文法を使いましたが、ここでは状態の見通しを優先して元の簡潔な文法に戻しています（拡張文法でも同様にオートマトンを作れますが、状態数が増えて図が煩雑になるためです）。

\begin{center}
\begin{tikzpicture}[
  >=stealth',
  node distance=3cm,
  state/.style={circle, draw, minimum size=12mm},
  accept/.style={state, double},
  every edge/.style={draw, ->, >=stealth'}
]
% 状態の定義（LR(0)項集合に基づく）
\node[state] (I0) at (0, 0) {\small $I_0$};
\node[state] (I1) at (4, 0) {\small $I_1$};
\node[accept] (I2) at (8, 0) {\small $I_2$};
\node[state] (I3) at (2, -3) {\small $I_3$};
\node[state] (I4) at (6, -3) {\small $I_4$};
\node[state] (I5) at (4, -6) {\small $I_5$};

% 遷移の定義
% I0からの遷移
\draw (I0) edge node[above] {D} (I1);
\draw (I0) edge[bend right=20] node[left] {P} (I3);
\draw (I0) edge[bend left=30] node[above left] {(} (I4);

% I1からの遷移（受理状態へ）
\draw (I1) edge node[above] {\$} (I2);

% I3からの遷移
\draw (I3) edge[bend left=20] node[above] {(} (I4);

% I4からの遷移
\draw (I4) edge node[below] {P} (I5);
\draw (I4) edge[loop above] node[above] {(} ();

% I5からの遷移
\draw (I5) edge[bend left=20] node[right] {)} (I3);
\draw (I5) edge[bend right=40] node[below] {P} (I3);

% 開始矢印
\draw[->] ([xshift=-1.5cm]I0.west) -- (I0);
\end{tikzpicture}
\end{center}

状態名が\(I_0\)のような番号になっているのは、後で導入するLR(0)項集合に対応づけるためです。イメージしやすいように、ここでは次のように読み替えてください。

\begin{itemize}
\item
  \(I_0\)：まだ何も読んでいない初期状態
\item
  \(I_1\)：\texttt{D}を読み終えて終端\texttt{\$}を待っている状態
\item
  \(I_2\)：入力全体を受理した状態
\item
  \(I_3\)：括弧列\texttt{P}を読み進めている状態
\item
  \(I_4\)：
  \texttt{(}を読んだ直後で、その中身の\texttt{P}をこれから処理する状態
\item
  \(I_5\)：\texttt{(\ P}まで読み、対応する\texttt{)}を待っている状態
\end{itemize}

このオートマトン（後で説明しますが、LR(0)オートマトンと呼ばれます）は、単に

\begin{itemize}
\item
  \textgt{状態}：今どの規則をどこまで読んでいるか
\item
  \textgt{遷移}：記号を読んだら次の状態へ
\end{itemize}

\noindent を表現したものです。

このオートマトンを事前に構築しておけば、解析時には以下の単純な操作で構文解析が可能になります。

\begin{enumerate}
\item
  現在の状態を確認
\item
  次の入力記号を見る
\item
  オートマトンの遷移をたどる
\item
  指示された動作（シフト/還元）を実行（こちらは別途計算が必要です。後述）
\end{enumerate}

これがLR構文解析の本質です。ここから、算術式の文法を簡略化した以下の文法を使って、LRオートマトンの構築方法について説明していきます。

\begin{lstlisting}
  S → T
  S → T + S
  T → x
\end{lstlisting}

\subsection{拡張開始記号の必要性}
\index{かくちょうかいしきごう@拡張開始記号}
\index{にゅうりょくしゅうたん@入力終端}
\index{じゅり@受理}

その前に、LR構文解析でよく使われる\textgt{拡張開始記号}（augmented start
symbol）を追加するのが一般的です。たとえば、先ほどの文法に対してはまず最初に、新しい開始記号\texttt{S\textquotesingle{}}と規則
\texttt{S\textquotesingle{}\ →\ S\ \$} を追加します。

拡張開始記号を追加した文法は以下のようになります。

\begin{lstlisting}
  S' → S $
  S → T
  S → T + S
  T → x
\end{lstlisting}

\texttt{S\textquotesingle{}}が拡張開始記号であり、\texttt{\$}は入力の終端を示します。

元の文法では、構文解析が「いつ完了するか」が曖昧でした。拡張開始記号を追加することで「\texttt{S\textquotesingle{}\ →\ S\ \$・}まで読み終えれば、構文解析が終了する」という形で受理状態（構文解析の完了）を明確に示すことができます。

拡張開始記号自体は本質的ではありませんが、実用上は重要なので覚えておくと良いでしょう。

\subsection{LR(0)項 --- 解析の進行状況を表す}
\index{LR(0)項}
\index{どっと@ドット}
\index{かんげん@還元}

オートマトン構築のために必要な概念が\textbf{\sffamily LR(0)}\textgt{項}です。LR(0)項は、構文解析の進行状況を表現するための基本的な単位です。

文法規則の中にドット（・）を挿入して、「どこまで認識したか」を表し、各々をLR(0)項とします。

たとえば、規則 \texttt{S\ →\ T\ +\ S} に対するLR(0)項は

\begin{itemize}
\item
  \texttt{S\ →\ ・T\ +\ S}：まだ何も認識していない
\item
  \texttt{S\ →\ T・\ +\ S}：Tまで認識した
\item
  \texttt{S\ →\ T\ +・\ S}：T + まで認識した
\item
  \texttt{S\ →\ T\ +\ S・}：すべて認識した（還元可能）
\end{itemize}

\noindent の4つとなります。

\subsection{LR(0)項集合 --- 同じ「状況」をまとめる}
\index{LR(0)項集合}
\index{じょうたい@状態}

ここで重要な概念が\textbf{\sffamily LR(0)}\textgt{項集合}です。構文解析の過程で、複数のLR(0)項が同時に「有効」になることがあります。これらをまとめて1つの「状態」として扱います。

たとえば、以下のような状況を考えてみましょう。

\begin{lstlisting}
  現在の状態 = {
      S → T・+ S    （Tを認識済み、次は+を期待）
      S → T・       （Tを認識済み、ここで還元も可能）
  }
\end{lstlisting}

この状態では、2つの可能性が同時に存在しています。

\begin{enumerate}
\item
  次に\texttt{+}が来れば、最初の規則に従って解析を続ける
\item
  ここで\texttt{S\ →\ T}として還元することもできる
\end{enumerate}

LR(0)項集合は、このような「同じ状況で有効な項の集まり」を表現します。オートマトンの各状態は、実はLR(0)項の集合なのです。

なぜ集合として扱うのでしょうか？それは、構文解析の過程で複数の可能性を同時に追跡する必要があるからです。たとえば、\texttt{if}文の解析中に、\texttt{else}節があるかないかの両方の可能性を考慮する必要があるような場合です。

\subsection{閉包 --- 項目の展開}
\index{へいほう@閉包}
\index{closure@Closure}
\index{LR(0)項}

LR(0)項集合を構築する上で最も重要なのが、LR(0)項の\textgt{閉包}（Closure）です。これは、あるLR(0)項から「論理的に導かれる全ての項」を表す概念です。

たとえば、LR(0)項 \texttt{S\textquotesingle{}\ →\ ・S\ \$}
があるとします。これは「これからSを認識したい」という状況を表しています。しかし、Sを認識するためには、Sから始まる規則も考慮する必要があります。このために閉包を求めます。

\begin{enumerate}
\item
  \texttt{S\textquotesingle{}\ →\ ・S\ \$} （これからSを認識したい）
\item
  Sを認識するには、Sから始まる規則も必要

\begin{lstlisting}
  S → ・T + S
  S → ・T
\end{lstlisting}

\item
  Tを認識するためには以下の規則も必要

\begin{lstlisting}
  T → ・x
\end{lstlisting}

\item
  最終的に、LR(0)項\texttt{S\textquotesingle{}\ →\ ・S\ \$}の閉包は以下のようになります。

\begin{lstlisting}
  {
    S' → ・S $,
    S → ・T + S,
    S → ・T,
    T → ・x
  }
\end{lstlisting}
\end{enumerate}

このように、ドットの直後に非終端記号がある場合、その非終端記号が左辺にある全ての規則を追加していく必要があります。

\subsubsection{LR(0)項の閉包を求めるアルゴリズム}

以下はLR(0)項の閉包を計算するアルゴリズムです。

\begin{lstlisting}[escapeinside={/*}{*/}]
  closure(I) {
      J = I                                   // /*{\footnotesize 初期項目集合をコピー} */
      repeat {
          for (各項目 A → α・Bβ in J) {    // /*{\footnotesize ドットの後に非終端記号Bがある項目を探す}*/
              for (各規則 B → γ) {          // /*{\footnotesize Bから始まるすべての規則を取得}*/
                  J = J ∪ {B → ・γ}        // /*{\footnotesize ドットを先頭に置いた項目を追加}*/
              }
          }
      } until Jに変化がなくなる             // /*{\footnotesize 新しい項目が追加されなくなるまで繰り返す}*/
      return J                              // /*{\footnotesize 閉包された項目集合を返す}*/
  }
\end{lstlisting}

\textgt{アルゴリズムの動作解説}

\begin{enumerate}
\item
  \textgt{初期化}：結果集合Jに入力項目集合Iをコピー
\item
  \textgt{項目のスキャン}：現在のJの各項目を調べ、ドットの直後に非終端記号Bがあるかチェック
\item
  \textgt{規則の展開}：Bから始まるすべての文法規則を取得し、ドットを先頭に置いた項目を作成
\item
  \textgt{項目の追加}：新しい項目をJに追加（重複は自動で除去）
\item
  \textgt{収束判定}：新しい項目が追加されなくなったら終了
\end{enumerate}

このアルゴリズムにより、あるLR(0)項に対する閉包を求めることができます。

\subsection{LR(0)オートマトンの全体像}

閉包を理解したところで、これをどのようにオートマトン構築に活用するかを見ていきましょう。まず、完成したLR(0)オートマトンがどのような構造になるかを見てみます。

先ほどと同じ、以下の文法に対するLR(0)オートマトンを考えます。

\begin{lstlisting}
  S' → S $
  S → T
  S → T + S
  T → x
\end{lstlisting}

このオートマトンは以下のような構造になります。

\begin{center}
\begin{tikzpicture}[
  scale=0.9, transform shape,
  >=stealth',
  node distance=3cm,
  state/.style={circle, draw, minimum size=12mm},
  accept/.style={state, double},
  every edge/.style={draw, ->, >=stealth'}
]
% 状態の定義
\node[state] (I0) at (0, 0) {\small $I_0$};
\node[state] (I1) at (4, 1) {\small $I_1$};
\node[state] (I2) at (4, -1) {\small $I_2$};
\node[state] (I3) at (8, -2) {\small $I_3$};
\node[accept] (I4) at (8, 1) {\small $I_4$};
\node[state] (I5) at (8, 0) {\small $I_5$};
\node[state] (I6) at (12, 0) {\small $I_6$};

% 遷移の定義
% I0からの遷移
\draw (I0) edge node[above left] {S} (I1);
\draw (I0) edge node[below left] {T} (I2);
\draw (I0) edge[bend right=20] node[below] {x} (I3);

% I1からの遷移（受理状態へ）
\draw (I1) edge node[above] {\$} (I4);

% I2からの遷移
\draw (I2) edge node[above] {+} (I5);

% I5からの遷移
\draw (I5) edge node[above] {S} (I6);
\draw (I5) edge[bend left=15] node[below] {T} (I2);
\draw (I5) edge[bend left=25] node[below] {x} (I3);

% 開始矢印
\draw[->] ([xshift=-1.5cm]I0.west) -- (I0);
\end{tikzpicture}
\end{center}

\textgt{オートマトンの構成要素}

\begin{enumerate}
\item
  \textgt{状態（円）}：各状態は LR(0)項の集合（閉包）を表す
\item
  \textgt{遷移（矢印）}：文法記号を読むことで状態が変化する
\item
  \textgt{初期状態（\(I_0\)）}：構文解析の開始時の状態
\item
  \textgt{受理状態（\(I_4\)）}：構文解析が完了した状態
\end{enumerate}

参考までに、図中の各状態 \(I_0\)〜\(I_6\) が表す
LR(0)項集合を先取りで示しておきます。具体的な作り方は4.10.11以降で順を追って説明しますが、まずは「各状態がどの項集合に対応するか」のイメージをつかんでください。

\begin{itemize}
\item
  \(I_0\) = \{S' → ・S \$, S → ・T, S → ・T + S, T → ・x\}
\item
  \(I_1\) = \{S' → S・\$\}
\item
  \(I_2\) = \{S → T・, S → T・+ S\}
\item
  \(I_3\) = \{T → x・\}
\item
  \(I_4\) = \{S' → S \$・\}（受理状態）
\item
  \(I_5\) = \{S → T +・S, S → ・T, S → ・T + S, T → ・x\}
\item
  \(I_6\) = \{S → T + S・\}
\end{itemize}

\textgt{オートマトンの動作}

\begin{itemize}
\item
  入力文字列を1文字ずつ読みながら状態を遷移する
\item
  各状態で「どの規則をどこまで読んだか」が分かる
\item
  受理状態に到達すれば、入力は正しく解析された
\end{itemize}

ここで一点補足しておきます。図をよく見ると、\(I_3\)（\texttt{T\ →\ x・}）や
\(I_6\)（\texttt{S\ →\ T\ +\ S・}）のように、外向きの矢印が描かれていない状態があります。これらは\textgt{還元状態}で、ドットが規則の右辺の最後まで進み終わっていることを表します。一見「行き詰まり」のように見えますが、実際には還元状態に到達するとパーサーは「\textgt{還元}」を実行し、スタックから対応する記号を取り除いて代わりに左辺の非終端記号を積み、別の状態へ復帰します。

たとえば \(I_3\) に到達した場合は規則 \texttt{T\ →\ x}
で還元され、\texttt{x} をスタックから取り除いて \texttt{T}
を積み、ひとつ前の状態（\(I_0\) か \(I_5\)）から \texttt{T}
での遷移をたどって \(I_2\)
に戻ります。つまり、図中の矢印はあくまで「シフト」と「\texttt{T} や
\texttt{S}
のような非終端記号での遷移（GOTO）」だけを示しており、還元による状態の戻り方は次節以降で扱う構文解析表で具体的に与えられます。

ともあれ、オートマトンを構築することで、構文解析の進行状況を効率的に追跡できるのです。

ここからはステップバイステップで先ほどの文法に対して、このLR(0)オートマトンを構築する過程を見ていきましょう。

\subsection{ステップ1：初期状態（$I_0$）の作成}

初期状態は、拡張開始記号から始まる項目の閉包です。

\(I_0\) = closure(\{S' → ・S \$\})

項目の閉包を計算すると

\begin{enumerate}
\item
  \texttt{S\textquotesingle{}\ →\ ・S\ \$} がある
\item
  ドットの後にSがあるので、Sから始まる規則を追加

\begin{itemize}
\item
  \texttt{S\ →\ ・T}
\item
  \texttt{S\ →\ ・T\ +\ S}
\end{itemize}

\item
  \texttt{S\ →\ ・T} のドットの後にTがあるので、Tから始まる規則を追加

\begin{itemize}
\item
  \texttt{T\ →\ ・x}
\end{itemize}
\end{enumerate}

結果として、以下のような初期状態\(I_0\)が得られます。

\(I_0\) = \{S' → ・S \$, S → ・T, S → ・T + S, T → ・x\}

\subsection{ステップ2：GOTO関数による遷移の計算}
\index{GOTO関数}
\index{せんい@遷移}

GOTO関数は、ある状態から記号を読んだときの次の状態を計算します。

\begin{lstlisting}
  GOTO(I, X) = closure({A → αX・β | A → α・Xβ ∈ I})
\end{lstlisting}

\(I_0\)からの各遷移を計算してみましょう。

GOTO(\(I_0\), S) =
\(I_0\)でドットの後にSがある項目に対して状態\(I_1\)を求める。

\begin{enumerate}
\item
  \texttt{S\textquotesingle{}\ →\ ・S\ \$}
  のドットの後にSがあるので、\texttt{S\textquotesingle{}\ →\ S・\$}
  を得る
\item
  上記の項集合の閉包をとって状態\(I_1\)を得る

  \(I_1\) = closure(\{S' → S・\$\}) = \{S' → S・\$\}
\end{enumerate}

GOTO(\(I_0\), T) =
\(I_0\)でドットの後にTがある項目に対して状態\(I_2\)を求める。

\begin{enumerate}
\item
  \texttt{S\ →\ ・T} のドットの後にTがあるので、\texttt{S\ →\ T・}
  を得る
\item
  \texttt{S\ →\ ・T\ +\ S}
  のドットの後にTがあるので、\texttt{S\ →\ T・\ +\ S} を得る
\item
  上記の項集合の閉包をとって状態\(I_2\)を得る

  \(I_2\) = closure(\{S → T・, S → T・ + S\}) = \{S → T・, S → T・ + S\}
\end{enumerate}

GOTO(\(I_0\), x) =
\(I_0\)でドットの後にxがある項目に対して状態\(I_3\)を求める。

\begin{enumerate}
\item
  \texttt{T\ →\ ・x} のドットの後にxがあるので、\texttt{T\ →\ x・}
  を得る
\item
  上記の項目集合の閉包をとって状態\(I_3\)を得る

  \(I_3\) = closure(\{T → x・\}) = \{T → x・\}
\end{enumerate}

\subsection{ステップ3：全ての状態を探索}

\(I_1\)からの遷移：

GOTO(\(I_1\), \$) =
\(I_1\)でドットの後に\$がある項目に対して状態\(I_4\)を求める。

\begin{enumerate}
\item
  \texttt{S\textquotesingle{}\ →\ S・\$}
  のドットの後に\$があるので、\texttt{S\textquotesingle{}\ →\ S\ \$・}
  を得る
\item
  これは受理状態なので、\(I_4\)は受理状態として定義

  \(I_4\) = \{S' → S \$・\}（受理状態）
\end{enumerate}

\(I_2\)からの遷移：

GOTO(\(I_2\), +) =
\(I_2\)でドットの後に\(+\)がある項目に対して状態\(I_5\)を求める。

\begin{enumerate}
\item
  \texttt{S\ →\ T・+\ S}
  のドットの後に+があるので、\texttt{S\ →\ T\ +・S} を得る
\item
  \texttt{S}から始まる規則も追加するため、\texttt{S\ →\ ・T}、\texttt{S\ →\ ・T\ +\ S}
  を得る
\item
  さらに\texttt{T}から始まる規則も追加するため、\texttt{T\ →\ ・x}
  を得る
\item
  上記の項目集合の閉包をとって状態\(I_5\)を得る

  \(I_5\) = closure(\{S → T +・S, S → ・T, S → ・T + S, T → ・x\})\\
  = \{S → T +・S, S → ・T, S → ・T + S, T → ・x\}
\end{enumerate}

\(I_5\)からの遷移：

GOTO(\(I_5\), S) =
\(I_5\)でドットの後にSがある項目に対して状態\(I_6\)を求める。

\begin{enumerate}
\item
  \texttt{S\ →\ T\ +・S}
  のドットの後にSがあるので、\texttt{S\ →\ T\ +\ S・} を得る
\item
  上記の項目集合の閉包をとって状態\(I_6\)を得る

  \(I_6\) = \{S → T + S・\}
\end{enumerate}

GOTO(\(I_5\), T) =
\(I_5\)でドットの後にTがある項目は状態\(I_2\)と同じになります。

GOTO(\(I_5\), x) =
\(I_5\)でドットの後にxがある項目は状態\(I_3\)と同じになります。

\subsection{完成したLR(0)オートマトン}

状態：

\begin{itemize}
\item
  \(I_0\) = \{S' → ・S \$, S → ・T, S → ・T + S, T → ・x\}
\item
  \(I_1\) = \{S' → S・\$\}
\item
  \(I_2\) = \{S → T・, S → T・+ S\}
\item
  \(I_3\) = \{T → x・\}
\item
  \(I_4\) = \{S' → S \$・\}（受理状態）
\item
  \(I_5\) = \{S → T +・S, S → ・T, S → ・T + S, T → ・x\}
\item
  \(I_6\) = \{S → T + S・\}
\end{itemize}

遷移：

\begin{itemize}
\item
  \(I_0\) \texttt{-\/-S-\/-\textgreater{}} \(I_1\)
\item
  \(I_0\) \texttt{-\/-T-\/-\textgreater{}} \(I_2\)
\item
  \(I_0\) \texttt{-\/-x-\/-\textgreater{}} \(I_3\)
\item
  \(I_1\) \texttt{-\/-\$-\/-\textgreater{}} \(I_4\)
\item
  \(I_2\) \texttt{-\/-+-\/-\textgreater{}} \(I_5\)
\item
  \(I_5\) \texttt{-\/-S-\/-\textgreater{}} \(I_6\)
\item
  \(I_5\) \texttt{-\/-T-\/-\textgreater{}} \(I_2\)
\item
  \(I_5\) \texttt{-\/-x-\/-\textgreater{}} \(I_3\)
\end{itemize}

これを図にすると、以下のようになります（先ほど提示した図と同じです）。

\begin{center}
\begin{tikzpicture}[
  scale=0.9, transform shape,
  >=stealth',
  node distance=3cm,
  state/.style={circle, draw, minimum size=12mm},
  accept/.style={state, double},
  every edge/.style={draw, ->, >=stealth'}
]
% 状態の定義
\node[state] (I0) at (0, 0) {\small $I_0$};
\node[state] (I1) at (4, 1) {\small $I_1$};
\node[state] (I2) at (4, -1) {\small $I_2$};
\node[state] (I3) at (8, -2) {\small $I_3$};
\node[accept] (I4) at (8, 1) {\small $I_4$};
\node[state] (I5) at (8, 0) {\small $I_5$};
\node[state] (I6) at (12, 0) {\small $I_6$};

% 遷移の定義
% I0からの遷移
\draw (I0) edge node[above left] {S} (I1);
\draw (I0) edge node[below left] {T} (I2);
\draw (I0) edge[bend right=20] node[below] {x} (I3);

% I1からの遷移（受理状態へ）
\draw (I1) edge node[above] {\$} (I4);

% I2からの遷移
\draw (I2) edge node[above] {+} (I5);

% I5からの遷移
\draw (I5) edge node[above] {S} (I6);
\draw (I5) edge[bend left=15] node[below] {T} (I2);
\draw (I5) edge[bend left=25] node[below] {x} (I3);

% 開始矢印
\draw[->] ([xshift=-1.5cm]I0.west) -- (I0);
\end{tikzpicture}
\end{center}

\subsection{LR(0)構文解析表の構築}
\index{こうぶんかいせきひょう@構文解析表}
\index{ACTION表}
\index{GOTO表}

LR(0)オートマトンが完成したら、次は構文解析表を作成します。LR(0)構文解析表は、オートマトンの状態と入力記号に基づいて、どのアクションを取るべきかを決定するためのものです。表には2種類の情報を記録します。

\begin{enumerate}
\item
  ACTION表：終端記号に対するアクション

\begin{itemize}
\item
  \(s_n\)：シフトして状態nへ遷移
\item
  \(r_n\)：規則nで還元
\item
  acc：受理
\end{itemize}

\item
  GOTO表：非終端記号に対する遷移
\end{enumerate}

\subsection{構築規則}

\begin{enumerate}
\item
  項目 \texttt{A\ →\ α・aβ} が状態iにあり、GOTO(i, a) = j なら

\begin{itemize}
\item
  ACTION{[}i, a{]} = \(s_j\)
\end{itemize}

\item
  項目 \texttt{A\ →\ α・} が状態iにあり、\texttt{A}が拡張開始記号でない（還元可能）なら

\begin{itemize}
\item
  全ての終端記号に対してACTION{[}i, *{]} = r (A→α)
\end{itemize}

\item
  項目 \texttt{S\textquotesingle{}\ →\ S\ \$・} が状態iにあるなら

\begin{itemize}
\item
  ACTION{[}i, \${]} = acc
\end{itemize}

\item
  GOTO(i, A) = j なら

\begin{itemize}
\item
  GOTO{[}i, A{]} = j
\end{itemize}
\end{enumerate}

\subsubsection{完成した解析表}

\begin{table}[htb]
  \begin{flushleft}
  \begin{tabular}{cccccc}\hline
状態 & x & + & \$ & S & T\\ \hline
\(I_0\) & \(s_3\) & - & - & 1 & 2\\ \hline
\(I_1\) & - & - & \(s_4\) & - & -\\ \hline
\(I_2\) & - & \(s_5\) & \(r_1\) & - & -\\ \hline
\(I_3\) & - & \(r_3\) & \(r_3\) & - & -\\ \hline
\(I_4\) & - & - & acc & - & -\\ \hline
\(I_5\) & \(s_3\) & - & - & 6 & 2\\ \hline
\(I_6\) & - & - & \(r_2\) & - & -\\ \hline
  \end{tabular}
  \end{flushleft}
\end{table}

凡例：

\begin{itemize}
\item
  \(s_n\)：シフトして状態nへ
\item
  \(r_n\)：規則\(n\)で還元（\(r_1\): S→T, \(r_2\): S→T+S, \(r_3\): T→x）
\item
  \texttt{acc}：受理
\item
  数字のみ：GOTO遷移
\end{itemize}

なお、上記の表は説明の見通しを優先して整理したものです。厳密に「還元可能な項目があれば任意の終端に対して還元する」というLR(0)の規則をそのまま適用すると、状態\(I_2\)では
\texttt{S\ →\ T・} による還元と \texttt{S\ →\ T・+\ S} の \texttt{+}
シフトが競合するため\textgt{シフト/還元コンフリクト}が生じます。この競合を「\texttt{+}
ならシフト、\texttt{\$}
なら還元」と区別するには、Sの後ろにきうる記号（FOLLOW(S) =
\{\$\}）を見て判断する必要があり、それが次節で扱うSLR(1)の発想です。本節の表は「LR(0)の状態構築までは厳密に行い、アクションの埋め方では先取りしてSLR(1)的な絞り込みを行ったもの」と理解してください。

また、状態\(I_4\)の\texttt{acc}は、入力終端\texttt{\$}をシフトして
\texttt{S\textquotesingle{}\ →\ S\ \$・}
に到達した受理状態であることを示しています。実行時にはすでに\texttt{\$}を読み終えていますが、表では「この状態に来たら受理する」という印として便宜上\texttt{\$}列に\texttt{acc}を書いています。

\subsection{LR(0)オートマトン由来の表による解析の実行}

構築した構文解析表を使って、入力 \texttt{x\ +\ x\ \$}
を解析してみましょう。各ステップでは現在の状態と入力記号からアクションを決定します。

ステップ1：最初のxを処理

\begin{itemize}
\item
  スタック: \texttt{{[}0{]}}
\item
  入力: \texttt{x\ +\ x\ \$}
\item
  アクション: \texttt{ACTION{[}0,\ x{]}} = \(s_3\)
\item
  実行: xをシフトして状態3へ
\item
  結果: スタック \texttt{{[}0,\ 3{]}}、入力\texttt{+ x \$}
\end{itemize}

ステップ2：T→xで還元

\begin{itemize}
\item
  スタック: \texttt{{[}0,\ 3{]}}
\item
  入力: \texttt{+ x \$}
\item
  アクション: \texttt{ACTION{[}3,\ +{]}} = \(r_3\) (\texttt{T\ →\ x})
\item
  実行: \texttt{T\ →\ x}で還元

  \begin{itemize}
  \item
    スタックから1つ削除: \texttt{{[}0{]}}
  \item
    \texttt{GOTO{[}0,\ T{]}\ =\ 2}
  \item
    状態2をプッシュ: \texttt{{[}0,\ 2{]}}
  \end{itemize}
\end{itemize}

ステップ3：+をシフト

\begin{itemize}
\item
  スタック: \texttt{{[}0,\ 2{]}}
\item
  入力: \texttt{+ x \$}
\item
  アクション: \texttt{ACTION{[}2,\ +{]}} = \(s_5\)
\item
  実行: +をシフトして状態5へ
\item
  結果: スタック \texttt{{[}0,\ 2,\ 5{]}}、入力 \texttt{x\ \$}
\end{itemize}

ステップ4：2番目のxをシフト

\begin{itemize}
\item
  スタック: \texttt{{[}0,\ 2,\ 5{]}}
\item
  入力: \texttt{x\ \$}
\item
  アクション: \texttt{ACTION{[}5,\ x{]}} = \(s_3\)
\item
  実行: xをシフトして状態3へ
\item
  結果: スタック \texttt{{[}0,\ 2,\ 5,\ 3{]}}、入力 \texttt{\$}
\end{itemize}

ステップ5：T→xで還元

\begin{itemize}
\item
  スタック: \texttt{{[}0,\ 2,\ 5,\ 3{]}}
\item
  入力: \texttt{\$}
\item
  アクション: \texttt{ACTION{[}3,\ \${]}} = \(r_3\) (\texttt{T\ →\ x})
\item
  実行: \texttt{T\ →\ x}で還元

  \begin{itemize}
  \item
    スタックから1つ削除: \texttt{{[}0,\ 2,\ 5{]}}
  \item
    \texttt{GOTO{[}5,\ T{]}\ =\ 2}
  \item
    状態2をプッシュ: \texttt{{[}0,\ 2,\ 5,\ 2{]}}
  \end{itemize}
\end{itemize}

ステップ6：S→Tで還元

\begin{itemize}
\item
  スタック: \texttt{{[}0,\ 2,\ 5,\ 2{]}}
\item
  入力: \texttt{\$}
\item
  アクション: \texttt{ACTION{[}2,\ \${]}} = \(r_1\) (\texttt{S\ →\ T})
\item
  実行: \texttt{S\ →\ T}で還元

  \begin{itemize}
  \item
    スタックから1つ削除: \texttt{{[}0,\ 2,\ 5{]}}
  \item
    \texttt{GOTO{[}5,\ S{]}\ =\ 6}
  \item
    状態6をプッシュ: \texttt{{[}0,\ 2,\ 5,\ 6{]}}
  \end{itemize}
\end{itemize}

ステップ7：S→T+Sで還元

\begin{itemize}
\item
  スタック: \texttt{{[}0,\ 2,\ 5,\ 6{]}}
\item
  入力: \texttt{\$}
\item
  アクション: \texttt{ACTION{[}6,\ \${]}} = \(r_2\)
  (\texttt{S\ →\ T\ +\ S})
\item
  実行: \texttt{S\ →\ T\ +\ S}で還元

  \begin{itemize}
  \item
    スタックから3つ削除: \texttt{{[}0{]}}
  \item
    \texttt{GOTO{[}0,\ S{]}\ =\ 1}
  \item
    状態1をプッシュ: \texttt{{[}0,\ 1{]}}
  \end{itemize}
\end{itemize}

ステップ8：入力終端\texttt{\$}をシフト

\begin{itemize}
\item
  スタック: \texttt{{[}0,\ 1{]}}
\item
  入力: \texttt{\$}
\item
  アクション: \texttt{ACTION{[}1,\ \${]}\ =\ \(s_4\)}
\item
  実行: \texttt{\$} をシフトして状態4へ
\item
  結果: スタック \texttt{{[}0,\ 1,\ 4{]}}、入力は空
\end{itemize}

ステップ9：受理

\begin{itemize}
\item
  スタック: \texttt{{[}0,\ 1,\ 4{]}}
\item
  入力: 空
\item
  状態\(I_4\)は \texttt{S\textquotesingle{}\ →\ S\ \$・} を含む受理状態
\item
  実行: 解析成功！
\end{itemize}

このように、この構文解析表を使った構文解析は

\begin{enumerate}
\item
  現在の状態と入力記号から表を引く
\item
  シフトか還元かを決定
\item
  還元の場合は、どの規則を使うかも一意に決まる
\end{enumerate}

\noindent という単純な操作の繰り返しで、効率的に構文解析を実現します。

\subsection{LR(0)の限界：コンフリクト}
\index{こんふりくと@コンフリクト}
\index{しふとかんげんこんふりくと@シフト/還元コンフリクト}

LR(0)は効率的ですが、重要な限界があります。ここでも先ほどと同じ文法を考えてみましょう。

\begin{lstlisting}
  S → T
  S → T + S
  T → x
\end{lstlisting}

この文法でLR(0)オートマトンを構築すると、ある状態\texttt{X}で以下の項を含むことになります。

\begin{lstlisting}
  X = {
      S → T・       （還元準備完了）
      S → T・+ S    （+を期待）
  }
\end{lstlisting}

この状態で問題が発生します。なぜなら

\begin{itemize}
\item
  入力が\texttt{+}なら → シフトして\texttt{+}を読むべき
\item
  入力が\texttt{\$}（終端）なら → \texttt{S\ →\ T}で還元すべき
\end{itemize}

しかし、LR(0)は\textgt{次の入力を見ない}ので、どちらを選ぶか決められません。これを\textgt{シフト/還元コンフリクト}と呼びます。

\subsection{なぜコンフリクトが起きるのか}

素朴なシフト還元では、規則の優先順位や「最長一致」のような暗黙のルールでなんとか動いていた選択を、LR(0)では「状態だけ」を見て判断しようとします。そのため、状態の中に「還元可能な項目」と「シフト可能な項目」が併存していると、どちらを選ぶべきか機械的には決められず、コンフリクトとして表面化するのです。

このコンフリクトの多さから、LR(0)は単独では実用パーサージェネレータにほとんど採用されていません。ただし、LR(0)で構築した状態機械は、続くSLR(1)・LALR(1)でもそのまま土台として再利用されます。LR(0)を学ぶ価値は、ここから先のすべてのLR系手法の出発点だという点にあります。

\section{SLR(1)、LR(1)、LALR(1) --- 先読みによる改良}
\index{SLR(1)}
\index{LR(1)}
\index{LALR(1)}
\index{さきよみ@先読み}
\index{FOLLOW集合}

LR(0)のコンフリクトを解決するため、\textgt{先読み}（次の入力を見る）を導入した手法が開発されました。

\begin{tcolorbox}[title=LR系構文解析の歴史]
本書ではLR(0) → SLR(1) → LR(1) → LALR(1)の順に解説していますが、これは「先読み機構が単純な順」に並べた教育上の順序であり、\textgt{歴史的な発見順とは逆向き}です。実際の年表は次のようになっています。

\begin{itemize}
\item
  \textbf{\sffamily 1965}\textgt{年}：Donald E. Knuth, ``On the Translation of Languages from Left to Right'' でLR(k)を導入。LR(1)を含む一般的なLR(k)が決定的文脈自由言語をきれいに扱えることを示しました。ただしLR(1)は状態数が膨大で、当時の計算機資源では実装が困難でした。
\item
  \textbf{\sffamily 1969}\textgt{年}：Frank DeRemerが博士論文 ``Practical Translators for LR(k) Languages''（MIT）で、LR(1)を実装可能なサイズに圧縮する手法としてLALR(1)を提案。
\item
  \textbf{\sffamily 1971}\textgt{年}：Frank DeRemer, ``Simple LR(k) grammars''（CACM）で、LR(0)オートマトンをそのまま使いFOLLOW集合で先読みを賄う、さらに簡素なSLR(k)を提案。SLRの``S''は``Simple''、つまり「LRをいかに単純化するか」という当時の問題意識を表しています。
\item
  \textbf{\sffamily 1982}\textgt{年}：DeRemer \& Pennello, ``Efficient Computation of LALR(1) Look-Ahead Sets''（TOPLAS）でLALR(1)の効率的構築アルゴリズムが確立され、Bisonなど後発の実用システムに広く採用されました。
\end{itemize}

つまり \textbf{\sffamily LR(1)}\textgt{が先にあり、それを実用化するために}\textbf{\sffamily LALR・SLR}\textgt{という「刈り込み」が後から作られた}というのが実情です。本書が逆順を採る理由は、LR(0)の素朴な状態構築から先読み機構を段階的に積み上げる方が、各手法の存在意義（何の問題を解いたのか）を把握しやすいからです。実装上の発見の順序と概念理解の順序は別物として整理しておくと、第5章で構文解析器生成系を選ぶ際にも感覚がつかみやすくなります。
\end{tcolorbox}

\subsection{SLR(1) --- FOLLOW集合による改良}

\begin{quote}
要点:
LR(0)で起こる「還元すべきか判断できない」問題を、還元する非終端記号Aの後ろにきうる記号（FOLLOW(A)）に先読み記号が含まれているときだけ還元する、というルールで解決した手法。多くのコンフリクトが消えるが、FOLLOW集合は文法全体から大雑把に計算されるため、より精密な文脈判断が必要なケースには対応しきれない（→
LR(1) へ）。
\end{quote}

\textbf{\sffamily SLR(1)}（Simple
LR(1)）\footnote{Frank DeRemer, ``Simple LR(k) grammars'',
  Communications of the ACM, 14(7), 1971.
  SLR法を提案した論文です。なお、LALR法そのものはDeRemerの博士論文
  ``Practical Translators for LR(k) Languages'' (MIT, 1969)
  で提案され、効率的な構築アルゴリズムは DeRemer \& Pennello,
  ``Efficient Computation of LALR(1) Look-Ahead Sets'', ACM TOPLAS,
  4(4), 1982 で確立されました。}は、LL(1)で学んだFOLLOW集合の考え方を使います。

\noindent\textgt{基本アイデア：}
\begin{quotation}
  「ある非終端記号Aの後にきうる記号」が分かれば、還元すべきタイミングが分かる。
\end{quotation}

先ほどのコンフリクトの例を考えてみましょう。

\begin{lstlisting}
  X = {
      S → T・       （還元準備完了）
      S → T・+ S    （+を期待）
  }
\end{lstlisting}

この状態で、次の入力が\texttt{+}か\texttt{\$}かによって、どのアクションを取るべきかが変わります。SLR(1)では、以下のように判断します。

\begin{itemize}
\item
  \texttt{FOLLOW(S)\ =\ \{\$\}}（この文法では、Sの後ろには入力終端だけがくる）
\item
  次の入力が\texttt{+}なら
  \begin{itemize}
  \item
    シフト（\texttt{S\ →\ T・+\ S}のため）
  \item
    \texttt{+}は\texttt{FOLLOW(S)}に含まれないので、\texttt{S\ →\ T}では還元しない
  \end{itemize}
\item
  次の入力が\texttt{\$}なら
  \begin{itemize}
  \item
    還元のみ（\texttt{\$} ∈ FOLLOW(S)）
  \end{itemize}
\end{itemize}

これで多くのコンフリクトが解消されます。

\noindent\textbf{\sffamily SLR(1)}\textgt{の改善点：}

LR(0)では「還元可能な状態では全ての入力に対して還元」していましたが、SLR(1)では「その非終端記号の後に実際にくる可能性がある記号（FOLLOW集合）の時だけ還元」します。これにより、多くのコンフリクトが解決されます。

\noindent\textbf{\sffamily SLR(1)}\textgt{の限界：}

ただし、FOLLOW集合は文法全体から計算される「大まかな」情報なので、より複雑な文法では依然としてコンフリクトが残る場合があります。

\subsection{LR(1) --- より精密な先読み}

\begin{quote}
要点:
SLR(1)が「文法全体から計算した大雑把なFOLLOW」を使うのに対し、LR(1)は項目自体に「この文脈で還元した直後にきうる記号」を埋め込んで持ち回ることで、文脈ごとに精密な判断を行う手法。理論上、あらゆる決定的文脈自由言語を解析できる最強クラスだが、状態数が爆発しやすい（→
実用化のためにLALR(1)へ）。
\end{quote}

\noindent\textbf{\sffamily LR(1)}\textgt{の基本アイデア：}

SLR(1)では文法全体から計算したFOLLOW集合を使っていましたが、LR(1)では「その具体的な解析文脈での先読み記号」を使います。

\begin{lstlisting}
  LR(1)項の形式： [A → α・β, a]
\end{lstlisting}

ここで\texttt{a}は「この項目で還元した後に期待される記号」を表します。同じ規則でも、到達した文脈によって異なる先読み記号を持つため、SLR(1)では区別できなかった状況も正確に判断できます。

\noindent\textbf{\sffamily LR(1)}\textgt{の特徴：}

\begin{itemize}
\item
  最強の能力：理論上、あらゆる決定的文脈自由言語を解析可能
\item
  メモリコスト：先読み記号の組み合わせにより状態数が大幅に増加する場合がある
\end{itemize}

同じ\texttt{E\ →\ T・}でも、文脈によって異なる先読み記号を持つため、より精密な判断が可能です。

\subsection{LALR(1) --- 実用的な妥協点}

\begin{quote}
要点:
LR(1)の状態を「コア部分（ドットの位置と規則）が同じで先読み記号だけが異なる」もの同士で統合し、状態数をLR(0)と同程度に抑えた手法。LR(1)の解析能力をほぼ保ったまま実用的なメモリ量に収まるため、Yacc/Bisonなど主要なパーサージェネレータが採用しています。LR(1)で解けるが
LALR(1)
では解けない（マージで再発するコンフリクトがある）特殊な文法は存在しますが、現実のプログラミング言語ではほとんど問題になりません。
\end{quote}

LR(1)は強力ですが、状態数の爆発という問題があります。LR(0)が文法によって十数〜数十状態程度なのに対し、LR(1)は同じ規模の文法でも数百〜数千状態以上に膨らむことがあります（具体的な数は文法のサイズと構造に依存します）。

LALR(1)（Look-Ahead
LR(1)）は、この問題を解決する「実用的な妥協点」です。

LALR(1)の基本アイデアは、LR(1)で生成される多くの状態が「コア部分（ドットの位置と規則）は同じで、先読み記号だけが異なる」という特徴に着目することです。LALR(1)では、これらの状態を統合します。

LR(1)の2つの状態

\begin{itemize}
\item
  状態A: \texttt{{[}E → T・, +{]}}
\item
  状態B: \texttt{{[}E → T・, \${]}}
\end{itemize}

LALR(1)では統合

\begin{itemize}
\item
  新状態: \texttt{{[}E → T・, \{+, \$\}{]}}
\end{itemize}

LALR(1)の特徴は、状態数がLR(0)と同程度に抑えられること、SLR(1)より強力でほとんどの実用言語を扱えること、そしてYacc/Bisonなど主要なパーサージェネレータが標準採用していることです。LR(1)の解析能力をほぼ保ちながら、状態数をSLR(1)と同程度に抑えられる「実用的な妥協点」として広く使われています。

\begin{lstlisting}
  LR(1) 状態A = {
      [A → α・β, a]
      [B → γ・δ, b]
  }

  LR(1) 状態B = {
      [A → α・β, c]
      [B → γ・δ, d]
  }
\end{lstlisting}

これをマージすると、以下のようになります。

\begin{lstlisting}
  LALR(1) 状態 = {
      [A → α・β, {a, c}]
      [B → γ・δ, {b, d}]
  }
\end{lstlisting}

先読み記号が和集合になることに注目してください。

\subsection{LALR(1)の構築方法}

LALR(1)オートマトンの構築には主に2つの方法があります。

\subsubsection{方法1：LR(1)からの変換}

\begin{enumerate}
\item
  完全なLR(1)オートマトンを構築
\item
  同じコアを持つ状態を識別
\item
  それらの状態をマージ（先読み記号は和集合）
\item
  遷移関係を調整
\end{enumerate}

この方法は概念的に分かりやすいですが、一時的に大きなLR(1)オートマトンを構築する必要があります。

\subsubsection{方法2：直接構築}

\begin{enumerate}
\item
  LR(0)オートマトンを構築
\item
  各状態に対して、効率的に先読み記号を計算
\item
  先読み伝播グラフを使って先読み記号を伝播
\end{enumerate}

実用的な構文解析器生成系（Yacc）は、効率的な方法2を採用しています。

\subsection{LALR(1)で新たに生じるコンフリクト}
\index{LALR(1)}
\index{こんふりくと@コンフリクト}
\index{かんげんかんげんこんふりくと@還元/還元コンフリクト}

状態のマージにより、LR(1)では解決できていたコンフリクトが再発することがあります。以下の文法を考えてみましょう。

\begin{lstlisting}
  S' → S $
  S → a A d | b B d | a B e | b A e
  A → c
  B → c
\end{lstlisting}

この文法は、前半の記号（aかb）と後半の記号（dかe）の組み合わせで、中間のAかBを決定する必要があります。

LR(1)では、以下のような、同じLR(0)コアを持つ状態が別々に存在します。

\begin{lstlisting}
  状態X: {
      [A → c・, d]
      [B → c・, e]
  }

  状態Y: {
      [A → c・, e]
      [B → c・, d]
  }
\end{lstlisting}

LR(1)では文脈ごとに先読み記号が分かれているため、それぞれの状態では還元先を一意に決められます。しかしLALR(1)では、同じLR(0)コアを持つこれらの状態がマージされます。

\begin{lstlisting}
  マージ後の状態: {
      [A → c・, d/e]
      [B → c・, d/e]
  }
\end{lstlisting}

この状態で、先読みがdまたはeの時、AとBのどちらで還元すべきか決定できません（reduce/reduceコンフリクト）。

\subsection{LALR(1)とLR(1)の実用的な比較}

\noindent 実際のプログラミング言語の文法での比較：

\begin{table}[htb]
  \begin{flushleft}
    \begin{tabular}{lll}\hline
特性 & LR(1) & LALR(1)\\ \hline
状態数 & 数千～数万 & 数百（LR(0)と同程度）\\ \hline
構文解析表サイズ & 非常に大きい & 実用的なサイズ\\ \hline
構築時間 & 遅い & 高速\\ \hline
解析能力 & 最も強力 & わずかに劣る\\ \hline
実用性 & メモリ制約で困難 & 十分実用的\\ \hline
    \end{tabular}
  \end{flushleft}
\end{table}

\subsection{なぜYaccはLALR(1)を選んだのか}

1975年にStephen C.
JohnsonがYaccを開発した際、LALR(1)を採用した理由は以下のようなものです。

\begin{enumerate}
\item
  メモリ制約：当時のコンピュータでは、LR(1)の巨大な解析表は非現実的
\item
  十分な表現力：ほとんどのプログラミング言語の文法はLALR(1)で記述可能
\item
  効率的な実装：DeRemerの効率的なアルゴリズムが利用可能
\end{enumerate}

\subsection{LALR(1)の実用例}

実際のプログラミング言語の構文解析ではLALR(1)が使われていることが多いです。

\begin{itemize}
\item
  C言語：元々Yaccで記述されたLALR(1)文法
\item
  Ruby：Yaccを使用（LALR(1)）
\item
  PostgreSQL：SQL文法をYaccで解析
\end{itemize}

\subsection{LALR(1)の限界を超えて}

LALR(1)でもコンフリクトが解決できない場合には以下のような対処法があります。

\begin{enumerate}
\item
  文法の書き換え：多くの場合、等価な文法に変換可能
\item
  GLR（Generalized LR）：非決定的な解析を許容
\item
  優先順位と結合性：Yaccの\texttt{\%left}、\texttt{\%right}、\texttt{\%prec}
\item
  意味的な処理：構文解析後に意味解析で解決
\end{enumerate}

GLR\footnote{GLR (Generalized LR) は Masaru Tomita, \emph{Efficient
  Parsing for Natural Language}, Kluwer, 1985
  で自然言語処理向けに提案された手法で、コンフリクトを「すべての可能性を並行して試す」ことで処理し、結果として曖昧な文法も解析可能にします。同様の発想を下向き型のLLに対して適用したものとして
  GLL（Scott \& Johnstone, ENTCS, 2010）もあります。}はオリジナルのYaccでなく互換実装のBisonでサポートされていますが、あまり使われていません。ほとんどのケースではLALR(1)で十分です。

\subsection{まとめ：素朴な方法からの進化}

\begin{enumerate}
\item
  素朴なシフト還元：毎回すべての規則をチェック
\item
  LR(0)：オートマトンで効率化、ただし先読みなし
\item
  SLR(1)：FOLLOW集合で簡単な先読み
\item
  LR(1)：文脈ごとの先読みで精密化
\item
  LALR(1)：実用的なバランス
\end{enumerate}

この流れを理解することで、各手法の存在意義が明確になります。

\subsection{LR系の早見表}

LR系の各手法の特徴を、後で参照しやすいように一覧にまとめておきます。
\clearpage
\begin{table}[htb]
  \small
  \begin{center}
    \begin{tabular}{lp{2cm}p{2cm}p{3.5cm}p{3cm}}\hline
      手法 & 解析能力 & 状態数の傾向 & 主な用途・採用例 & 補足\\ \hline
LR(0) & 弱い（コンフリクト多発） & 少ない & 教育目的が中心。実用パーサージェネレータでは単独採用はまれ & 還元時に先読みを見ない\\ \hline
SLR(1) & 中程度 & LR(0)と同じ & 比較的単純な文法 & FOLLOW集合で還元タイミングを絞る\\ \hline
LR(1) & 最強（あらゆる決定的文脈自由言語） & 数百〜数千以上に膨らむ & 主に理論研究や、状態数を許容できる小規模文法 & 項目自体に先読み記号を持たせる\\ \hline
LALR(1) & LR(1)にほぼ匹敵 & LR(0)と同程度 & Yacc/Bison、PostgreSQL、Rubyの構文解析など実用 & LR(1)の状態をコアごとに統合した妥協点\\ \hline
    \end{tabular}
  \end{center}
\end{table}

Yaccで shift/reduce conflictやreduce/reduce
conflictに出くわした場合、本書のSLR(1)〜LALR(1)節で扱った「コンフリクトの起こる仕組み」を念頭に置くと、警告の意味とその対処（先読みで解消できるのか、文法自体を書き換える必要があるのか）が判断しやすくなります。

\section{Parsing Expression Grammar(PEG) --- 新しいアプローチ}
\index{PEG}
\index{Parsing Expression Grammar}
\index{かいせきひょうげんぶんぽう@解析表現文法}
\index{じゅんじょつきせんたく@順序付き選択}
\index{ばっくとらっく@バックトラック}
\index{ぶんみゃくいぞんげんご@文脈依存言語}

2004年にBryan Fordが提案した\textbf{\sffamily Parsing Expression
Grammar（PEG）}は、全く違うアプローチを取ります。これまでの構文解析手法は大前提として、文脈自由文法を元にしていましたが、PEGは文脈自由文法に一見よく似た、それでいて異なる新たな形式文法です。異なる形式文法であるため、CFGと表現能力はことなります。

正確にいうと、PEG自体を構文解析アルゴリズムというのは多少不適切なのですが、PEG自体に標準的な操作的意味論（実行の方法くらいの意味）が定義されているため、ここではPEGを構文解析アルゴリズムとして扱います。

\subsection{PEGの基本アイデア}

PEGの最大の特徴は以下のとおりです。

\begin{enumerate}
\item
  順序付き選択とバックトラック：\texttt{A\ /\ B}は「まずAを試し、失敗したらBを試す」
\item
  字句解析が不要：文字列レベルで直接解析
\item
  いくつかの文脈依存言語を扱える：
  たとえば、\(a^n\ b^n\ c^n\)も扱える
\item
  全ての決定的文脈自由言語を表現可能
\item
  非常に単純：アルゴリズムが非常に単純で、実装が容易
\end{enumerate}

\subsection{第2章で実装したPEG}

第2章で実装した\texttt{PegJsonParser}がまさにPEGの実装でした。

\begin{codescreen}{}
public Ast.JsonArray parseArray() {
    int backup = cursor;
    try {
        // LBRACKET RBRACKET
        parseLBracket();
        parseRBracket();
        return new Ast.JsonArray(new ArrayList<>());
    } catch (ParseException e) {
        cursor = backup;  // バックトラック！
        // LBRACKET value {COMMA value} RBRACKET
        parseLBracket();
        // ... 省略 ...
    }
}
\end{codescreen}

この「バックトラック」がPEGの核心です。決定的なLL/LRでは、構文解析は「一度決めたら戻れない」ので、原則としてバックトラックができません（GLLやGLRなどは除きます）。しかし、PEGでは「失敗したら前の状態に戻って別の選択肢を試す」ことができます。決定的なLL/LRが非常に複雑なのは、ひとえに「バックトラックができない」がゆえに、先読みやFIRST/FOLLOW集合といった仕掛けを駆使して規則を一意に選ぶ必要があるからでもあります。PEGはこのバックトラックを行うことで、非常にシンプルな構文解析を実現しています。

\subsection{PEGの形式的定義}
\index{PEG!形式的定義}
\index{こうていじゅつご@肯定述語}
\index{ひていじゅつご@否定述語}
\index{れんせつ@連接}
\index{くりかえし@繰り返し}

PEGは以下の8つの構成要素から成り立ちます。

\begin{enumerate}
\item
  空文字列： ε
\item
  終端記号： t
\item
  非終端記号： N
\item
  連接： e1 e2（e1の後にe2）
\item
  選択： e1 / e2（e1またはe2）
\item
  0回以上の繰り返し： e*
\item
  肯定述語： \&e（eにマッチするが消費しない）
\item
  否定述語： !e（eにマッチしないことを確認）
\end{enumerate}

特に7番と8番は文脈自由文法にはない、PEG特有の機能です。
\clearpage
\subsection{PEGとCFGの違い}
\index{PEG!CFGとの違い}
\index{CFG}
\index{あいまいせい@曖昧性}
\index{ひだりさいき@左再帰}
\index{かいぶん@回文}

\begin{table}[h]
  \begin{center}
    \begin{tabular}{lll}\hline
      特徴 & CFG & PEG\\ \hline
      選択演算子 & \texttt{｜}（非決定的） & \texttt{/}（順序付き）\\ \hline
      曖昧性 & あり得る & 常に一意\\ \hline
      字句解析 & 必須 & 不要\\ \hline
      左再帰 & 問題なし（LR系） & 無限ループ\\ \hline
      表現力 & 文脈自由言語 & LR(k) + 一部の文脈依存言語\\ \hline
    \end{tabular}
  \end{center}
\end{table}

言語クラスとしてみたとき、PEGが認識する言語のクラスはCFGに含まれません。CFGでは表現できない\(a^n\ b^n\ c^n\)のような言語をPEGは表現できるからです。逆向き、つまり「CFGの認識する言語クラスがPEGに含まれるか」は未解決問題（open
problem）ですが、以下のような状況証拠から成り立たないだろうと予想されています。

\begin{itemize}
\item
  回文を表現できない可能性が高い：
  \texttt{A\ →\ aAa\ \textbar{}\ bAb\ \textbar{}\ ε}
  のような回文を生成する文法と等価な文法は、PEGではまだ知られていません。PEGで回文に近いものを書こうとすると、2の累乗 − 2の長さの回文という非常に不自然な表現になることが知られています。
\item
  線形時間で解析が可能： PEGは後述するPackrat
  Parsingにより、線形時間で解析が可能です。CFG全体に対する線形時間の解析アルゴリズムは知られていません。
\end{itemize}

\section{Packrat Parsing --- PEGの線形時間化}
\index{Packrat Parsing}
\index{ぱっくらっとぱーしんぐ@パックラット構文解析}
\index{せんけいじかん@線形時間}
\index{めもか@メモ化}
\index{PEG}

PEGの最大の弱点は、バックトラックにより最悪の場合に指数関数時間がかかることです。\textbf{\sffamily Packrat
Parsing}\footnote{Bryan Ford, ``Packrat parsing: simple, powerful, lazy,
  linear time'', ICFP '02, 2002. PEGのオリジナル論文 (POPL '04)
  より前に発表されており、線形時間化のアイデアが先に提示されました。なお、PEGとPackrat
  Parsingで素朴には扱えなかった\textgt{左再帰}を扱うためのアルゴリズムは、Warth,
  Douglass, Millstein, ``Packrat Parsers Can Support Left Recursion'',
  PEPM '08, 2008 で提案されています。}は、メモ化（memoization）を使ってこの問題を解決します。

\subsection{メモ化とは？}
\index{めもか@メモ化}
\index{memoization}
\index{きゃっしゅ@キャッシュ}

メモ化（memoization）は、「同じ引数で関数を呼んだら、同じ結果が返る」ことを利用して、一度計算した結果をキャッシュする技法です。

「メモ化」という名前は「memoize」（記憶する）からきています。「メモリ化」ではなく「メモ化」であることに注意してください。

\subsection{フィボナッチ数で学ぶメモ化}

フィボナッチ数の素朴な実装を考えてみます。

\begin{codescreen}{}
public static long fib(long n) {
    if(n == 0) return 0L;
    if(n == 1) return 1L;
    else return fib(n - 1) + fib(n - 2); 
}
\end{codescreen}

この実装の問題点は、同じ値を何度も計算してしまうことです。メモ化を適用すると次のようになります。

\begin{codescreen}{}
private static Map<Long, Long> cache = new HashMap<>();

public static long fib(long n) {
    Long value = cache.get(n);
    if(value != null) return value;

    long result;
    if(n == 0) {
        result = 0L;
    } else if (n == 1) {
        result = 1L;
    } else {
        result = fib(n - 1) + fib(n - 2);
    }
    cache.put(n, result);
    return result;
}
\end{codescreen}

メモ化の効果はこのようなケースで絶大です。具体的には

\begin{itemize}
\item
  時間計算量：\(O(2^n)\) → \(O(n)\)
\item
  空間計算量：\(O(1)\) → \(O(n)\)
\end{itemize}

メモ化によって結果をキャッシュすることで大幅に計算量を削減できます。

\subsection{PEGパーサーのメモ化}

PEGパーサーにもメモ化を適用できます。単に全ての規則に対してメモ化を行ったものを\textbf{\sffamily Packrat
Parsing}と呼びます。Packrat
Parsingでは、各規則のパース結果をキャッシュし、同じ位置で同じ規則を再度パースする際にキャッシュを利用します。このようにすることで、PEGのバックトラックによる指数関数的な計算量を線形時間に抑えることができます。

Packrat Parsingの基本的な実装は以下のようになります。

\begin{codescreen}{}
public class PackratParser {
    private Map<Integer, Map<String, MemoEntry>> memo;
    
    static class MemoEntry {
        boolean success;
        int endPos;
    }
    
    private boolean memoized(String ruleName, Supplier<Boolean> parser) {
        // メモ化テーブルをチェック
        Map<String, MemoEntry> posCache = memo.get(pos);
        if (posCache != null) {
            MemoEntry entry = posCache.get(ruleName);
            if (entry != null) {
                // キャッシュヒット！
                pos = entry.endPos;
                return entry.success;
            }
        }
        
        // キャッシュミス：実際にパース
        int startPos = pos;
        boolean success = parser.get();
        int endPos = pos;
        
        // 結果をキャッシュに保存
        // ... 省略 ...
        
        return success;
    }
}
\end{codescreen}

ここでのポイントは\texttt{memo}というマップを使って、各位置（\texttt{pos}）と規則名（\texttt{ruleName}）に対するパース結果をキャッシュすることです。これにより、同じ位置で同じ規則を再度パースする際に、キャッシュを利用して高速化できます。

\subsection{Packrat Parsingの特徴}
\index{Packrat Parsing!特徴}
\index{せんけいじかん@線形時間}
\index{くうかんけいさんりょう@空間計算量}
\index{ひだりさいき@左再帰}

利点：

\begin{itemize}
\item
  線形時間：\(O(n)\)（\(n\)は入力長）
\item
  左再帰対応：工夫により可能
\item
  実装が単純（単なるメモ化）
\end{itemize}

欠点：

\begin{itemize}
\item
  メモリ使用量：\(O(nm)\)（\(m\)は規則数）
\item
  キャッシュミスのオーバーヘッド
\item
  並列化が困難
\end{itemize}

Packrat
Parsingは、PEGのバックトラックによる指数関数的な計算量を線形時間に抑えるための強力な手法です。しかし、理論上は線形時間であるものの、その代償としてメモリ使用量が大きくなるため、注意が必要です。

\section{構文解析アルゴリズムの計算量と表現力}
\index{けいさんりょう@計算量}
\index{じかんけいさんりょう@時間計算量}
\index{くうかんけいさんりょう@空間計算量}
\index{ひょうげんりょく@表現力}

各アルゴリズムの計算量をまとめると：

\begin{table}[htb]
  \begin{center}
    \begin{tabular}{llll}\hline
      アルゴリズム & 時間計算量 & 空間計算量 & 備考\\ \hline
LL(1) & \(O(n)\) & \(O(|N| × |T|)\) & \(|N|\):非終端記号数,
\(|T|\):終端記号数\\ \hline
LL(k) & \(O(n)\) & \(O(|N| × |T|^k)\) &
\(k\)が大きくなると表サイズが指数的に増大\\ \hline
SLR(1) & \(O(n)\) & O(状態数 × (\textbar T\textbar+\textbar N\textbar)) & 状態数はLR(0)と同じ\\ \hline
LR(1) & \(O(n)\) & O(状態数 × (\textbar T\textbar+\textbar N\textbar)) & 状態数が多い\\ \hline
LALR(1) & \(O(n)\) & O(状態数 × (\textbar T\textbar+\textbar N\textbar)) & 状態数はLR(0)と同程度\\ \hline
PEG & \(O(2^n)\) & \(O(n)\) & バックトラック用スタック\\ \hline
Packrat & \(O(n)\) & \(O(n × |P|)\) & \(|P|\): 規則数\\ \hline
    \end{tabular}
  \end{center}
\end{table}

表の「O表記」は、アルゴリズムの計算量を表す記法で、\(O(n)\)は「入力の長さに比例した時間」、\(O(n^2)\)は「入力の長さの2乗に比例した時間」を意味します。

注目すべきは、PEGを除くすべての手法が線形時間で解析できることです。Packrat
Parsingを使えばPEGも線形時間になります。

\section{まとめ}

この章では、文脈自由文法を実際に解析するためのさまざまなアルゴリズムを学びました。

\subsection{学んだアルゴリズムの整理}

\textgt{下向き構文解析（トップダウン）}

\begin{itemize}
\item
  予測的構文解析：Dyck言語を例に、スタックを使った実装
\item
  LL(1)構文解析：FIRST集合とFOLLOW集合による効率化
\end{itemize}

\textgt{上向き構文解析（ボトムアップ）}

\begin{itemize}
\item
  シフト還元構文解析：葉から根に向かって構文木を構築
\item
  LR(0)、SLR(1)、LALR(1)、LR(1)：段階的に強力になる手法
\end{itemize}

\textgt{その他の手法}

\begin{itemize}
\item
  PEG：順序付き選択による決定的な構文解析
\item
  Packrat Parsing：メモ化によるPEGの線形時間化
\end{itemize}

これらについて知っていれば、次のような判断ができるようになります。

\begin{enumerate}
\item
  構文解析器生成系の選択：各ツールの特性を理解して選択できる

\begin{itemize}
\item
  ANTLRはLL系でエラーリカバリーが優れている
\item
  Yaccは高速だがエラーメッセージが分かりにくい
\end{itemize}

\item
  エラーメッセージの理解：コンフリクトの意味がわかる

\begin{itemize}
\item
  「shift/reduce conflict」→ シフトと還元のどちらを選ぶか決まらない
\item
  「reduce/reduce conflict」→ 複数の還元規則から選べない
\end{itemize}

\item
  DSLの設計：パースしやすい文法設計ができる

\begin{itemize}
\item
  シンプルな設定ファイルならLL(1)で十分
\item
  複雑な表現が必要ならPEGやLALR(1)を検討
\end{itemize}

\item
  パフォーマンスの理解：適切な最適化を選択できる

\begin{itemize}
\item
  大量のデータ処理なら線形時間が保証される手法を選ぶ
\item
  IDEのリアルタイム解析ならインクリメンタル対応を考慮
\end{itemize}
\end{enumerate}

構文解析は難しく感じるかもしれませんが、基本的なアイデアは「文字列を解析して木構造を抽出する」というシンプルなものです。この章で学んだ各手法の特徴を理解しておけば、実際の開発で構文解析が必要になった時に、適切な選択ができるはずです。

次章では、ここで学んだアルゴリズムを実装したパーサージェネレータ（構文解析器生成系）について具体的に見ていきます。

\thispagestyle{empty}
\chapter{構文解析器生成系の世界}\label{ch05}
\index{こうぶんかいせききせいせいけい@構文解析器生成系}
\index{ぱーさーじぇねれーた@パーサージェネレータ}
\index{ぱーさーこんびねーた@パーサーコンビネータ}
\index{CFG}
\index{PEG}

第4章では現在知られている構文解析手法について、アイデアと概要を説明しました。しかし、構文解析の世界ではよく知られていることなのですが、第4章で説明した手法は毎回プログラマが手で実装する必要はありません。

というのは、CFG（文脈自由文法）やPEG（その類似表記も含む）によって記述された文法定義から特定の構文解析アルゴリズムを用いた構文解析器を生成する\textgt{構文解析器生成系}（Parser
Generator、パーサージェネレータ）というソフトウェアがあるからです。

簡単に言えば、「文法の定義を書いたら、対応したパーサーのコードを自動生成してくれるツール」です。これにより、複雑な構文解析アルゴリズムを知らなくても、文法定義さえ書ければ高性能なパーサーを作れます。

構文解析器生成系でもっとも代表的なものはYacc（ヤック）あるいはその互換実装であるGNU
Bison（バイソン）でしょう。YaccはLALR(1)法を利用したC言語の構文解析器を生成してくれるソフトウェアです。

ちなみに「Yacc」は「Yet Another Compiler
Compiler」（また別のコンパイラコンパイラ）の略で、「コンパイラを作るためのコンパイラ」という意味です。

この章では構文解析器生成系という種類のソフトウェアの背後にあるアイデアからはじまり、PEGから実際に対応する構文解析器を作る方法、多種多様な構文解析器生成系についての紹介を行います。

また、この章の最後ではある意味構文解析器生成系の一種とも言える\textgt{パーサーコンビネータ}の実装方法について踏み込んで説明します。

パーサーコンビネータとは、簡単に言えば「プログラミング言語に備わったファーストクラス関数やオブジェクトを組み合わせてパーサーを構築する方法」です。通常の構文解析器生成系が「文法定義ファイル」から「Javaコード」を生成するのに対し、パーサーコンビネータではJavaのメソッドを組み合わせて直接パーサーを作ります。単なる手書きパーサーと違い、パーサーコンビネータは規則や文法に対応した関数やオブジェクトがライブラリとして備わっているため、再利用性が高く、柔軟な構文解析が可能です。

構文解析器生成系を作るのは骨が折れますが、パーサーコンビネータであれば、いわゆる「ラムダ式」を持つほとんどのプログラミング言語で比較的簡単に実装できます。本書で利用しているJavaでも同様です。というわけで、本章を読めば皆さんもパーサーコンビネータを明日から自前で実装できるようになります。

\section{Dyck言語の文法とPEGによる構文解析器生成}
\index{Dyck言語}
\index{PEG}
\index{こうぶんかいせききせいせい@構文解析器生成}
\index{こーどせいせい@コード生成}
\index{ばっくとらっく@バックトラック}

これまで何度も登場したDyck言語（括弧の対応が取れた文字列を表す言語）は明らかにLL(1)法でもLR(1)法でもPEGによっても解析可能な言語です。しかし、手書きでパーサーを実装しようとすると、退屈な繰り返しコードが多くなりがちです。

実際のところ、Dyck言語を表現する文法があって、構文解析アルゴリズムがPEGということまでわかれば対応するJavaコードを\textgt{機械的に生成する}ことも可能そうに見えます。特に、構文解析はコード量が多いわりには退屈な繰り返しコードが多いものですから、文法からJavaコードを生成できれば劇的に工数を削減できそうです。

このように「文法と構文解析手法が決まれば、後のコードは自動的に決定可能なはずだから、機械に任せてしまおう」という考え方が構文解析器生成系というソフトウェアの背後にあるアイデアです。

早速ですが、以下のようにDyck言語を表す文法（第3章で示した
\texttt{P\ →\ (\ P\ )\ P\ \textbar{}\ ε}
に近いもの）がPEGで与えられたとして、構文解析器を生成する方法を考えてみましょう。

\subsection{Dyck言語のPEG}

PEGの記法では、\texttt{/}は「順序付き選択」を表し、BNFの\texttt{\textbar{}}とは異なり「最初の選択肢がマッチしたらそれを採用し、2番目以降は試さない」という意味になります。

\begin{codescreen}{}
D <- P;
P <- "(" P ")" P / ""; // "" は空文字列ε（イプシロン）を表す
\end{codescreen}

このPEGは\texttt{P} が \texttt{(P)P}
の形であるか、または空文字列であることを示します。PEGでは非終端記号の呼び出しは関数呼び出しとみなすことができますから、この定義に対応するパーサー関数（メソッド）のスケルトンを考えると、次のようなコードのイメージになります。

\begin{codescreen}{}
// 概念的なスケルトン
public boolean parseD() {
    return parseP(); // DはPに委譲
}

public boolean parseP() {
    // P <- "(" P ")" P / "";
    // 最初に "(" P ")" P を試す：
    //   "(" にマッチするか？
    //   マッチしたら、P を再帰的に呼び出す (1回目)
    //   Pの呼び出しが成功したら： 
    //     ")" にマッチするか？
    //     マッチしたら、P を再帰的に呼び出す (2回目)
    //     Pの呼び出しが成功したら：
    //       全体として成功
    // もし途中で失敗したら：
    //   バックトラックして "" を試す
    //   "" は常に成功（空文字列を消費）
    // どちらかが成功すれば parseP は成功
}
\end{codescreen}

次の項ではこのスケルトンを実際のJavaコードに変換して、Dyck言語の構文解析器を生成する方法を見ていきます。

\subsection{Dyck言語の構文解析器を生成する}

PEGからJavaコードへの機械的な変換は、以下の規則に従います。

\begin{itemize}
\item
  \textgt{非終端記号}\texttt{N}： \texttt{parseN()}
  という名前のメソッドを生成
\item
  \textgt{文字列リテラル}\texttt{str}：\texttt{match(str)}
  の呼び出しに変換
\item
  \textgt{連接} \texttt{e1\ e2}：
  \texttt{e1}の中身を再帰的に変換した結果と\texttt{e2}の中身を再帰的に変換した結果を連結
\item
  \textgt{順序付き選択} \texttt{e1\ /\ e2}：
  \texttt{e1}を試すコードと、失敗したらバックトラックするコード、\texttt{e2}を試すコードを連結する

  \begin{itemize}
  \item
    ここではバックトラックのために例外を使う
  \end{itemize}
\item
  \textgt{繰り返し}：ループ構造に変換
\end{itemize}

\begin{codescreen}{}
public class DyckParser {
    private String input;
    private int pos;

    public static class Failure extends RuntimeException {
        public Failure(String message) {
            super(message);
        }
    }
    
    public DyckParser(String input) {
        this.input = input;
        this.pos = 0;
    }
    
    // D <- P
    public void parseD() {
        parseP();
    }
    
    // P <- "(" P ")" P / ""
    public void parseP() {
        // 現在位置を保存（バックトラック用）
        int saved = pos;

        try {
            match("(");  // "(" にマッチするか？
            parseP();  // Pを再帰的に呼び出す（1回目）
            match(")");  // ")" にマッチするか？
            parseP();  // Pを再帰的に呼び出す（2回目）
            return;  // 成功
        } catch (Failure e) {
            // 失敗したらバックトラック
            pos = saved;
            // 空文字列 "" を試す（常に成功）
            return;
        }
    }
    
    // 文字列のマッチングを行うヘルパーメソッド
    private void match(String str) {
        if (pos + str.length() <= input.length() &&
            input.startsWith(str, pos)) {
            pos += str.length();
            return;
        }
        throw new Failure("Expected '" + str + "' at position " + pos);
    }
    
    // パース実行メソッド
    public boolean parse() {
        parseD();
        // 例外が起きず、全ての文字が消費されていれば成功
        if (pos == input.length()) {
            return true;  
        } else {
            return false;
        }
    }
}
\end{codescreen}

このようにして、比較的シンプルにPEGからDyck言語の構文解析器を生成することができました。例外を使ってバックトラックを実装していますが、これは単純化のためです。Javaで例外を投げるのはそこそこ重い処理なので、単純にif文で分岐する方法を使った方が効率的です。

\subsection{構文解析器生成系の実装}

実際の構文解析器生成系では、以下のようなステップで自動生成を行います。

\begin{enumerate}
\item
  \textbf{\sffamily PEG}\textgt{文法の解析}：PEG記法で書かれた文法定義自体を解析
\item
  \textgt{中間表現の生成}：文法規則を内部的なデータ構造に変換
\item
  \textgt{コード生成}：テンプレートを使ってJavaコードを出力
\end{enumerate}

以下は中間表現が生成された後のコード生成の例です。ここでは、PEGの規則を表すクラス\texttt{Rule}と、各種PEG式を表す\texttt{Expr}クラスを定義し、Javaコードを生成しています。

\begin{codescreen}{}
import java.util.*;
import java.util.stream.Collectors;
// 簡単な構文解析器生成系の例
public class PEG2Java {
    // 文法規則を表すレコード
    public record Rule(String name, Expr body) {}
    
    // PEG式を表す抽象クラス
    public abstract static class Expr {
        abstract String generate(int indentLevel);
    }
    
    // 文字列リテラル
    public static class Lit extends Expr {
        public final String value;
        public Lit(String value) {
            this.value = value;
        }
        String generate(int indentLevel) {
            return ("match(\"" + value + "\");").indent(indentLevel);
        }
    }

    // 非終端記号
    public static class NT extends Expr {
        public final String name;
        public NT(String name) {
            this.name = name;
        }
        // 非終端記号のコード生成
        String generate(int indentLevel) {
            return ("parse" + name + "();").indent(indentLevel);
        }
    }
    
    // 連接
    static class Seq extends Expr {
        public final List<Expr> exprs;
        public Seq(Expr... exprs) {
            this.exprs = List.of(exprs);
        }

        // 連接のコード生成。単純にN個の式を順に生成
        String generate(int indentLevel) {
            return exprs.stream()
                .map(e -> e.generate(indentLevel))
                .collect(Collectors.joining("\n")) + "\n";
        }
    }

    // 順序付き選択
    public static class Choice extends Expr {
        public final Expr alt1, alt2;
        public Choice(Expr alt1, Expr alt2) {
            this.alt1 = alt1;
            this.alt2 = alt2;
        }

        // 順序付き選択のコード生成
        // try-catchを使ってバックトラックを実装
        String generate(int indentLevel) {
            var code = """
                int saved = pos;
                try {
                    %s
                } catch (Failure e) {
                    pos = saved;
                    %s
                }
             """.formatted(
                alt1.generate(indentLevel + 4).stripTrailing(), 
                alt2.generate(indentLevel + 4));
            return code.indent(indentLevel);
        }
    }
    
    // コード生成のメインメソッド
    public static String generateParser(List<Rule> rules) {
        StringBuilder code = new StringBuilder();
        code.append("""
            public class Parser {
                public static class Failure extends RuntimeException {
                    public Failure(String message) {
                        super(message);
                    }
                }
                private String input;
                private int pos;
                public Parser(String input) {
                    this.input = input;
                    this.pos = 0;
                }
                public void match(String str) {
                    if (pos + str.length() <= input.length() &&
                        input.startsWith(str, pos)) {
                        pos += str.length();
                        return;
                    }
                    throw new Failure("Expected '" + str + "' at pos " + pos);
                }
        """.stripLeading());
        // 各規則に対してメソッドを生成
        for (Rule rule : rules) {
            code.append("    public void parse" + rule.name + "() {\n");
            code.append(rule.body.generate(8));
            code.append("\n");
            code.append("        return;\n");
            code.append("    }\n\n");
        }
        code.append("}\n");
        return code.toString();
    }
}
\end{codescreen}

先ほど作ったDyck言語のPEGをちょっと変形したものをこの構文解析器生成系に与えます。利用コードは次のようになります。

\begin{codescreen}{}
// D <- P;
// P <- "(" P ")" P / "";
List<Rule> rules = Arrays.asList(
    new Rule("D", new NT("P")),
    new Rule("P", 
        new Choice(
            new Seq(
                new Lit("("),
                new NT("P"),
                new Lit(")"),
                new NT("P")
            ),
            // 空文字列 ε（イプシロン）
            new Lit("")
         )
    )
);
// 構文解析器のコードを生成
String parserCode = PEG2Java.generateParser(rules);
System.out.println(parserCode);
\end{codescreen}

次のようなJavaコードが生成されます。

\begin{codescreen}{}
public class Parser {
    public static class Failure extends RuntimeException {
        public Failure(String message) {
            super(message);
        }
    }
    private String input;
    private int pos;
    public Parser(String input) {
        this.input = input;
        this.pos = 0;
    }
    public void match(String str) {
        if (pos + str.length() <= input.length() &&
            input.startsWith(str, pos)) {
            pos += str.length();
            return;
        }
        throw new Failure("Expected '" + str + "' at pos " + pos);
    }

    public void parseD() {
        parseP();
        return;
    }

    public void parseP() {
        int saved = pos;
        try {
            match("(");
            parseP();
            match(")");
            parseP();
        } catch (Failure e) {
            pos = saved;
            match("");
        }

        return;
    }
}
\end{codescreen}

空文字列への分岐（\texttt{match("")}）を残しているので、\texttt{()}のように\texttt{P}が空になるケースも、\texttt{(())}のようにネストするケースも成功します。\texttt{try}ブロックで失敗したときに位置を巻き戻してから空文字列を選択する、という5.1.2節と同じ挙動を保っています。

このコードは実際にDyck言語の変形版を解析できます。\texttt{parseD()}メソッドを呼び出すことで、Dyck言語の文字列が正しいかどうかをチェックできます。

このように、PEGの文法定義から機械的にJavaコードを生成できます。本体のコードだけでいえばわずか140行程度です。これでPEGの構文解析器生成系が実装できるのです。もちろん、実用的には抽象構文木（AST）の生成や、エラーメッセージの改善などが必要になりますが、コアの部分はこのようにシンプルに実装できます。

\section{構文解析器生成系の分類}
\index{こうぶんかいせききせいせいけい@構文解析器生成系!分類}
\index{LL(1)}
\index{LALR(1)}
\index{GLR}
\index{じくかいせききせいせいけい@字句解析器生成系}
\index{Lex}
\index{Yacc}

構文解析器生成系は1970年代頃から研究の蓄積があり、数多くの構文解析器生成系がこれまで開発されています。

基本的には、構文解析器生成系は採用している構文解析アルゴリズムによって分類されます。たとえば

\begin{itemize}
\item
  JavaCCはLL(k)構文解析器を出力するため、\textbf{\sffamily LL(k)}\textgt{構文解析器生成系}
\item
  YaccはLALR(1)構文解析器を出力するため、\textbf{\sffamily LALR(1)}\textgt{構文解析器生成系}
\end{itemize}

一般的な構文解析器生成系の処理フローは、おおむね以下のようになります。ただし、PEGを採用している構文解析器生成系では、字句解析器の部分が丸ごと不要になります。

\begin{figure}[htb]
\centering
\begin{tikzpicture}[
    scale=0.7, transform shape,
    node distance=2.2cm,
    box/.style={draw, rectangle, minimum width=3cm, minimum height=0.8cm, align=center, font=\small},
    diamond/.style={draw, ellipse, minimum width=2.8cm, minimum height=1cm, align=center, font=\small},
    arrow/.style={->, thick},
    label/.style={midway, above, font=\footnotesize}
]
    % ノードの定義
    \node[box] (A) {文法定義ファイル\\(.g, .y, .jjなど)};
    \node[diamond, right=1.8cm of A] (B) {構文解析器\\生成系};
    \node[box, above right=1cm and 1.5cm of B] (C) {字句解析器\\コード生成部};
    \node[box, right=1.8cm of C] (D) {生成された\\字句解析器コード\\(.java, .cなど)};
    \node[box, below right=1cm and 1.5cm of B] (E) {構文解析器\\コード生成部};
    \node[box, right=1.8cm of E] (F) {生成された\\構文解析器コード\\(.java, .cなど)};
    \node[diamond, right=3.5cm of B] (H) {コンパイラ};
    \node[box, right=1.8cm of H] (I) {実行可能な\\パーサー};

    % 矢印の定義
    \draw[arrow] (A) -- (B);
    \draw[arrow] (B) -- node[label,sloped] {字句解析ルール} (C);
    \draw[arrow] (C) -- (D);
    \draw[arrow] (B) -- node[label,sloped] {構文解析ルール} (E);
    \draw[arrow] (E) -- (F);
    \draw[arrow] (D) -- (H);
    \draw[arrow] (F) -- (H);
    \draw[arrow] (H) -- (I);
\end{tikzpicture}
\caption{一般的な構文解析器生成系の処理フロー}
\label{fig:parser-generator-flow}
\end{figure}

Yacc/Lexのように、字句解析器生成系（Lex）と構文解析器生成系（Yacc）が別々のツールとして提供され、連携して動作するケースもあります。

ただし、GNU BisonはYaccと違って、LALR(1)より広いGLR（Generalized
LR：一般化LR）構文解析器も生成できるので、GLR構文解析器生成系であるとも言えます。GLRは、通常のLR構文解析ではコンフリクト（シフト/還元コンフリクトなど）が発生するような曖昧な文法も扱えるように拡張された手法です。Bisonを使う場合、ほとんどはLALR(1)構文解析器を出力するので、GLRについて言及されることは少ないですが、知っておいても損はないでしょう。

より大きなくくりでみると、下向き構文解析（LL法やPEG）と上向き構文解析（LR法など）という観点から分類することもできますし、ともに文脈自由文法ベースであるLL法やLR法と、PEGなど他の形式言語を用いた構文解析法を対比してみせることもできます。

以下に、本書で紹介する代表的な構文解析器生成系の比較をまとめます。

\begin{table}[htb]
  \small
  \begin{center}
    \begin{tabular}{p{3cm}p{3cm}p{3cm}p{3cm}}\hline
      特徴項目 & JavaCC & Yacc & ANTLR\\ \hline
      採用アルゴリズム & LL(k)(デフォルトはLL(1)) & LALR(1) & ALL(*)\\ \hline
      生成コードの言語 & Java & C, C++ & Java, C++, Python, JavaScript, Go, C\#, Swift, Dart, PHP (多言語対応)\\ \hline
      左再帰の扱い & 不可(文法書き換えが必要) & 左再帰を扱える(直接・間接とも) & 直接左再帰を扱える(v4以降。間接左再帰は不可)\\ \hline
      曖昧性解決 & 先読みトークン数(k)の調整、意味アクション & 演算子の優先順位・結合規則指定、\%precなどで対応	 & 意味アクション、構文述語、ALL(*)による自動解決\\ \hline
      エラー報告/リカバリ機能 & 基本的 & errorトークンによる限定的なリカバリ & 高度なエラー報告、柔軟なエラーリカバリ戦略\\ \hline
      学習コスト & Javaユーザーには比較的容易 & やや高め(C言語と連携の知識も必要) & 機能が豊富で強力な分、やや高め\\ \hline
      その他特徴 & Javaに特化、構文がJavaライク & C言語との親和性が高い、歴史と実績がある & 強力な解析能力、豊富なターゲット言語、優れたツールサポート(GUIなど)\\ \hline
    \end{tabular}
  \end{center}
\end{table}

\section{JavaCC：Javaの構文解析器生成系の定番}
\index{JavaCC}
\index{LL(1)}
\index{LL(k)}
\index{LOOKAHEAD}
\index{TOKEN}
\index{SKIP}
\index{EOF}
\index{せまんてぃっくあくしょん@セマンティックアクション}

1996年、Sun
Microsystems（当時）は、Jackという構文解析器生成系をリリースしました。その後、Jackの作者が自らの会社を立ち上げ、JackはJavaCCに改名されて広く知られることとなりましたが、現在では紆余曲折の末、\href{https://javacc.github.io/javacc}{javacc.github.io}の元で開発およびメンテナンスが行われています。現在のライセンスは3条項BSDライセンスです。

JavaCCという名前は「Java Compiler
Compiler」の略で、Yacc同様「コンパイラを作るためのコンパイラ」という意味です。

JavaCCはLL(k)法を元に作られており、構文定義ファイルからLL(k)パーサーを生成します。以下は四則演算を含む数式を計算できる電卓をJavaCCで書いた場合の例です。

\begin{codescreen}{}
options {
  STATIC = false;
  JDK_VERSION = "21";
}

PARSER_BEGIN(Calculator)
package com.github.kmizu.calculator;
public class Calculator {
  public static void main(String[] args) throws ParseException {
    Calculator parser = new Calculator(System.in);
    System.out.println(parser.start());
  }
}
PARSER_END(Calculator)

SKIP : { " " | "\t"  | "\r" | "\n" }
TOKEN : {
  <ADD: "+">
| <SUBTRACT: "-">
| <MULTIPLY: "*">
| <DIVIDE: "/">
| <LPAREN: "(">
| <RPAREN: ")">
| <INTEGER: (["0"-"9"])+>
}

public int start() :
{int r = 0;}
{
  r=add() { return r; }
}

public int expression() :
{int r = 0;}
{
  r=add() <EOF> { return r; }
}

public int add() :
{int r = 0; int v = 0;}
{
    r=mult() ( <ADD> v=mult() { r += v; }| <SUBTRACT> v=mult() { r -= v; })* {
        return r;
    }
}


public int mult() :
{int r = 0; int v = 0;}
{
   r=primary() ( 
     <MULTIPLY> v=primary() { r *= v; }
   | <DIVIDE> v=primary() { r /= v; })* { return r; }
}

public int primary() :
{int r = 0; Token t = null;}
{
(
  <LPAREN> r=add() <RPAREN>
| t=<INTEGER> { r = Integer.parseInt(t.image); }
) { return r; }
}
\end{codescreen}

この\texttt{SKIP}から\texttt{TOKEN}までの部分はトークン定義になります。トークンとは、構文解析の基本単位となる要素で、「+」「123」「(」などのことです。第2章のJSONパーサーでも「文字列」「数値」「中括弧」などをトークンとして扱いましたね。ここでは、7つのトークンを定義しています。トークン定義の後が構文規則の定義になります。ここでは、

\begin{itemize}
\item
  \texttt{start()}
\item
  \texttt{expression()}
\item
  \texttt{add()}
\item
  \texttt{mult()}
\item
  \texttt{primary()}
\end{itemize}

\noindent の5つの構文規則が定義されています。各構文規則はJavaのメソッドに酷似した形で記述されますが、実際、ここから生成される.javaファイルには同じ名前のメソッドが定義されます。

\texttt{start()}と\texttt{expression()}は両方とも\texttt{add()}を呼び出すエントリポイントですが、違いは入力末尾の扱い方です。\texttt{expression()}は最後に\texttt{\textless{}EOF\textgreater{}}（入力終端トークン）を要求するため、入力全体が1つの式であることを厳密に検証したい場面（テストやファイル全体の構文チェックなど）に向きます。一方、\texttt{start()}は\texttt{\textless{}EOF\textgreater{}}を要求しないので、標準入力のような対話的な入力で「式を1つ読み終えたら抜ける」という使い方ができます。\texttt{main}メソッドでは標準入力を読むので\texttt{start()}を呼んでいます。

\texttt{add()}が\texttt{mult()}を呼び出して、\texttt{mult()}が\texttt{primary()}を呼び出すという構図は第1章ですでにみた形ですが、第1章と違って単純に宣言的に各構文規則の関係を書けばそれでOKなのが構文解析器生成系の強みです。

ちなみに、\texttt{\{int\ r\ =\ 0;\}}のような部分はJavaコードの埋め込みで、構文解析中に実行される処理を記述しています。これを「セマンティックアクション」（意味的動作）と呼びます。セマンティックアクションを使うことで、単に文法が正しいかチェックするだけでなく、解析しながら計算や抽象構文木の構築などの処理を行えます。

このようにして定義した電卓プログラムは次のようにして利用できます。

\begin{codescreen}{}
package com.github.kmizu.calculator;

import org.junit.jupiter.api.DisplayName;
import org.junit.jupiter.api.Test;
import com.github.kmizu.calculator.Calculator;
import java.io.*;

import static org.junit.jupiter.api.Assertions.assertEquals;

public class CalculatorTest {
    @Test
    @DisplayName("1 + 2 * 3 = 7")
    public void test1() throws Exception {
        Calculator calculator = new Calculator(
            new StringReader("1 + 2 * 3"))
        ;
        assertEquals(7, calculator.expression());
    }

    @Test
    @DisplayName("(1 + 2) * 4 = 12")
    public void test2() throws Exception {
        Calculator calculator = new Calculator(
            new StringReader("(1 + 2) * 4")
        );
        assertEquals(12, calculator.expression());
    }

    @Test
    @DisplayName("(5 * 6) - (3 + 4) = 23")
    public void test3() throws Exception {
        Calculator calculator = new Calculator(
            new StringReader("(5 * 6) - (3 + 4)")
        );
        assertEquals(23, calculator.expression());
    }
}
\end{codescreen}

この\texttt{CalculatorTest}クラスではJUnit5を使って、JavaCCで定義した\texttt{Calculator}クラスの挙動をテストしています。空白や括弧を含む数式を問題なく計算できているのがわかるでしょう。

このようなケースでは先読みトークン数が1のため、JavaCCのデフォルトで構いませんが、定義したい構文によっては先読み数を2以上に増やさなければいけないこともあります。以下のようにして先読み数を増やせます。

\begin{codescreen}{}
options {
  STATIC = false;
  JDK_VERSION = "21";
  LOOKAHEAD = 2;
}
\end{codescreen}

\texttt{LOOKAHEAD\ =\ 2}というオプションによって、先読みトークン数を2に増やしています。

先読みトークン数が2ということは、「次の2つのトークンを見て、どの規則を適用するか決定できる」という意味です。LOOKAHEADは固定されていれば任意の正の整数にできるので、JavaCCはデフォルトではLL(1)だが、オプションを設定することによってLL(k)になるともいえます。

JavaCCは構文定義ファイルの文法がかなりJavaに似ているため、生成されるコードの形を想像しやすいというメリットがあります。Javaの構文解析器生成系の中では最古の部類の割に今でも現役で使われているのは、Javaユーザーにとっての使いやすさが背景にあるように思います。

\section{Yacc：構文解析器生成系の老舗}
\index{Yacc}
\index{Bison}
\index{Lex}
\index{Flex}
\index{LALR(1)}
\index{error token@\texttt{error}トークン}

Yaccは1970年代にAT\&Tのベル研究所にいたStephen C.
Johnson（スティーブン・ジョンソン）によって作られたソフトウェアで、非常に歴史がある構文解析器生成系です。第4章で学んだLALR(1)法の構文解析器生成系を実用化した最初のツールでもあります。

現在広く使われているGNU
BisonはYaccのGNUによる再実装で、GNUプロジェクトの一部として開発されています。GNU
BisonはYaccと互換性があり、Yaccの機能を拡張したものです。

Yaccを使って、四則演算を行う電卓プログラムを作るにはまず字句解析器生成系であるLex（正確には、Vern
Paxsonによる互換実装であるFlex）用の定義ファイルを書く必要があります。Lexの定義ファイル\texttt{token.l}は次のようになります。

\begin{codescreen}{}
  %{
  #include "y.tab.h"
  extern int yylval;
  %}

  %%
  [0-9]+      { yylval = atoi(yytext); return NUM; }
  [-+*/()]    { return yytext[0]; }
  [ \t]       { /* ignore whitespace */ }
  "\r\n"      { return EOL; }
  "\r"        { return EOL; }
  "\n"        { return EOL; }
  .           { printf("Invalid character: %s\n", yytext); }
  %%
\end{codescreen}

\texttt{\%\%}から\texttt{\%\%}までがトークンの定義です。これはそれぞれ次のように読むことができます。

\begin{itemize}
\item
  \texttt{{[}0-9{]}+\ \ \ \ \ \ \{\ yylval\ =\ atoi(yytext);\ return\ NUM;\ \}}

  0-9までの数字が1個以上あった場合は数値として解釈し（\texttt{atoi(yytext)})、トークン\texttt{NUM}として返します。
\item
  \texttt{{[}-+*/(){]}\ \ \ \ \{\ return\ yytext{[}0{]};\ \}}

  \texttt{-},\texttt{+},\texttt{*},\texttt{/},\texttt{(},\texttt{)}についてはそれぞれの文字をそのままトークンとして返します。
\item
  \texttt{{[}\ \textbackslash{}t{]}\ \ \ \ \ \ \ \{\ /*\ ignore\ whitespace\ */\ \}}

  タブおよびスペースは単純に無視します
\item
  \texttt{"\textbackslash{}r\textbackslash{}n"\ \ \ \ \ \ \{\ return\ EOL;\ \}}

  改行はEOLというトークンとして返します
\item
  \texttt{"\textbackslash{}r"\ \ \ \ \ \ \ \ \{\ return\ EOL;\ \}}

  前に同じ
\item
  \texttt{"\textbackslash{}n"\ \ \ \ \ \ \ \ \{\ return\ EOL;\ \}}

  前に同じ
\item
  \texttt{.\ \ \ \ \ \ \ \ \ \ \ \{\ printf("Invalid\ character:\ \%s\textbackslash{}n",\ yytext);\ \}}

  それ以外の文字がきたらエラーを吐きます
\end{itemize}

このトークン定義ファイルを入力として与えると、flexは\texttt{lex.yy.c}という形で字句解析器を出力します。

次に、yaccの構文定義ファイルを書きます。

\begin{codescreen}{}
  %{
  #include <stdio.h>
  int yylex(void);
  void yyerror(char const *s);
  int yywrap(void) {return 1;}
  extern int yylval;
  %}

  %token NUM
  %token EOL
  %left '+' '-'
  %left '*' '/'

  %%

  input :
        | input line
        ;

  line : expr EOL
        { printf("Result: %d\n", $1); }
        | EOL
        | error EOL
        { yyerrok; }
        ;

  expr : NUM
        { $$ = $1; }
        | expr '+' expr
        { $$ = $1 + $3; }
        | expr '-' expr
        { $$ = $1 - $3; }
        | expr '*' expr
        { $$ = $1 * $3; }
        | expr '/' expr
        {
          if ($3 == 0) { yyerror("Cannot divide by zero."); YYERROR; }
          else { $$ = $1 / $3; }
        }
        | '(' expr ')'
        { $$ = $2; }
        ;

  %%

  void yyerror(char const *s)
  {
    fprintf(stderr, "Parse error: %s\n", s);
  }

  int main()
  {
    yyparse();
  }
\end{codescreen}

flexの場合と同じく、\texttt{\%\%}から\texttt{\%\%}までが構文規則の定義の本体です。実行されるコードが入っているので読みづらくなっていますが、それを除くと以下のようになります。

\begin{codescreen}{}
  %token NUM
  %token EOL
  %left '+' '-'
  %left '*' '/'

  %%
  input :
        | input line
        ;

  line  : expr EOL
        | EOL
        | error EOL
        ;

  expr  : NUM
        | expr '+' expr
        | expr '-' expr
        | expr '*' expr
        | expr '/' expr
        | '(' expr ')'
        ;
  %%
\end{codescreen}

だいぶ見やすくなりましたね。入力を表す\texttt{input}規則は、0個以上の\texttt{line}からなります。\texttt{line}は式を表す\texttt{expr}と改行からなるので、式を1行ずつ読み取って結果を出力できます。

次に、\texttt{expr}規則ですが、ここではyaccの優先順位を表現するための機能である\texttt{\%left}を使ったため、優先順位のためだけに規則を作る必要がなくなっており、定義が簡潔になっています。ともあれ、こうして定義された文法定義ファイルをyaccに与えると\texttt{y.tab.c}というファイルを出力します。

flexとyaccが出力したファイルを結合して実行ファイルを作ると、次のような入力に対して

\begin{codescreen}{}
1 + 2 * 3
2 + 3
6 / 2
3 / 0
\end{codescreen}

それぞれ、次のような出力を返します。

\begin{codescreen}{}
Result: 7
Result: 5
Result: 3
Parse error: Cannot divide by zero.
\end{codescreen}

Yaccはとても古いソフトウェアの1つですが、Rubyの文法定義ファイルparse.yはYacc用ですし、未だに各種言語処理系では現役で使われてもいます。

\section{ANTLR：多言語対応の強力な下向き構文解析器生成系}
\index{ANTLR}
\index{PCCTS}
\index{LL(k)}
\index{LL(*)}
\index{ALL(*)}
\index{Adaptive LL(*)}
\index{ひだりさいき@左再帰}
\index{じゅつご@述語}
\index{semantic predicate}

1989年にPurdue Compiler Construction Tool
Set（PCCTS）というものがありました。ANTLRはその後継というべきもので、これまでに、\texttt{LL(k) -> LL(*) -> ALL(*)}と構文解析アルゴリズムを拡張し、取り扱える文法の幅を広げつつアクティブに開発が続けられています。作者はTerence
Parrという方ですが、構文解析器一筋と言っていいくらいに、ANTLRにキャリアの多くを費やしてきている人です。

それだけに、ANTLRの完成度は非常に高いものになっています。また、一時期はLR法に比べてLL法の評価は低いものでしたが、Terence
ParrがLL(k)を改良していく過程で、LL(\texttt{*})やALL(\texttt{*})のようなLR法に比べてもなんら劣らない、実用的にも使いやすい構文解析法が発明されました。

ANTLRはJava、C++などいくつもの言語を扱うことができますが、特に安心して使えるのはJavaです。以下は先ほどと同様の、四則演算を解析できる数式パーサーをANTLRで書いた場合の例です。

ANTLRでは構文規則は、\texttt{規則名\ :\ 本体\ ;}
という形で記述しますが、LLパーサー向けの構文定義を素直に書き下すだけでOKです。

\begin{codescreen}{}
grammar Expression;

expression returns [int e]
    : v=additive {$e = $v.e;}
    ;

additive returns [int e = 0;]
    : l=multitive {$e = $l.e;} (
      '+' r=multitive {$e = $e + $r.e;}
    | '-' r=multitive {$e = $e - $r.e;}
    )*
    ;

multitive returns [int e = 0;]
    : l=primary {$e = $l.e;} (
      '*' r=primary {$e = $e * $r.e;}
    | '/' r=primary {$e = $e / $r.e;}
    )*
    ;

primary returns [int e]
    : n=NUMBER {$e = Integer.parseInt($n.getText());}
    | '(' x=expression ')' {$e = $x.e;}
    ;

LP : '(' ;
RP : ')' ;
NUMBER : INT ;
fragment INT :   '0' | [1-9] [0-9]* ; // no leading zeros
WS  :   [ \t\n\r]+ -> skip ;
\end{codescreen}

規則\texttt{expression}が数式を表す規則です。そのあとに続く\texttt{returns\ {[}int\ e{]}}はこの規則を使って解析を行った場合に\texttt{int}型の値を返すことを意味しています。これまで見てきたように構文解析をした後には抽象構文木をはじめとして何らかのデータ構造を返す必要があります。\texttt{returns\ ...}はそのために用意されている構文です。名前がすべて大文字の規則はトークンを表しています。

数式を表す各規則についてはこれまで書いてきた構文解析器と同じ構造なので読むのに苦労はしないでしょう。

規則\texttt{WS}は空白文字を表すトークンですが、これは数式を解析する上では読み飛ばす必要があります。\\
\texttt{{[}\ \textbackslash{}t\textbackslash{}n\textbackslash{}r{]}+\ -\textgreater{}\ skip}は

\begin{itemize}
\item
  スペース
\item
  タブ文字
\item
  改行文字
\end{itemize}

\noindent のいずれかが出現した場合は読み飛ばすということを表現しています。

ANTLRは下向き型の構文解析が苦手とする左再帰も、規則の中で直接的に自分自身を呼び出す形（直接左再帰）であれば扱えます（規則を経由した間接左再帰には対応していない点には注意してください）。先ほどの定義ファイルでは繰り返しを使っていましたが、これを直接左再帰に直した以下の定義ファイルも全く同じ挙動をします。

\begin{codescreen}{}
grammar LRExpression;

expression returns [int e]
    : v=additive {$e = $v.e;}
    ;

additive returns [int e]
    : l=additive op='+' r=multitive {$e = $l.e + $r.e;}
    | l=additive op='-' r=multitive {$e = $l.e - $r.e;}
    | v=multitive {$e = $v.e;}
    ;

multitive returns [int e]
    : l=multitive op='*' r=primary {$e = $l.e * $r.e;}
    | l=multitive op='/' r=primary {$e = $l.e / $r.e;}
    | v=primary {$e = $v.e;}
    ;

primary returns [int e]
    : n=NUMBER {$e = Integer.parseInt($n.getText());}
    | '(' x=expression ')' {$e = $x.e;}
    ;

LP : '(' ;
RP : ')' ;
NUMBER : INT ;
fragment INT :   '0' | [1-9] [0-9]* ; // no leading zeros
WS  :   [ \t\n\r]+ -> skip ;
\end{codescreen}

左再帰を使うことでより簡単に文法を定義できることもあるので、あるとうれしい機能だと言えます。

さらに、ANTLRは\texttt{ALL(*)}というアルゴリズムを採用しているため、通常のLLパーサーでは扱えないような文法定義も取り扱えます。以下の「最小XML」文法定義を見てみましょう。

\begin{codescreen}{}
grammar PetitXML;
@parser::header {
    import static com.github.asciidwango.parser_book.ch5.PetitXML.*;
    import java.util.*;
}

root returns [Element e]
    : v=element {$e = $v.e;}
    ;

element returns [Element e]
    : ('<' begin=NAME '>' es=elements '</' end=NAME '>' {
          $begin.text.equals($end.text)
      }?
      {$e = new Element($begin.text, $es.es);})
    | ('<' name=NAME '/>' {$e = new Element($name.text);})
    ;

elements returns [List<Element> es]
    : { $es = new ArrayList<>();} (element {$es.add($element.e);})*
    ;

LT:    '<';
GT:    '>';
SLASH: '/';
NAME:  [a-zA-Z_][a-zA-Z0-9]* ;

WS :   [ \t\n\r]+ -> skip ;
\end{codescreen}

\texttt{PetitXML}の名のとおり、属性やテキストなどは全く扱えず、\texttt{\textless{}a\textgreater{}}や\texttt{\textless{}a/\textgreater{}}、\texttt{\textless{}a\textgreater{}\textless{}b\textgreater{}\textless{}/b\textgreater{}\textless{}/a\textgreater{}}といった要素のみを扱えます。規則\texttt{element}が重要です。

\begin{codescreen}{}
element returns [Element e]
    : ('<' begin=NAME '>' es=elements '</' end=NAME '>' {
            $begin.text.equals($end.text)
       }?
      {$e = new Element($begin.text, $es.es);})
    | ('<' name=NAME '/>' {$e = new Element($name.text);})
    ;
\end{codescreen}

ここで空要素（\texttt{\textless{}e/\textgreater{}}など）に分岐するか、子要素を持つ要素（\texttt{\textless{}a\textgreater{}\textless{}b/\textgreater{}\textless{}/a\textgreater{}}など）に分岐するかを決定するには、\texttt{\textless{}}に加えて、任意の長さになり得るタグ名を先読みしなければいけません。通常のLLパーサーでは何文字（何トークン）先読みできるかは予め決定されているのでこのような文法定義を取り扱うことはできません。しかし、ANTLRの\texttt{ALL(*)}アルゴリズムやその前身となる\texttt{LL(*)}アルゴリズムでは任意個の文字数を先読みして分岐を決定できます。

ANTLRでは通常のLLパーサーと違い文法を記述する上での大きな制約がなく、非常に強力です。理論的な意味でも\texttt{ALL(*)}アルゴリズムは任意の決定的な文脈自由言語を取り扱えます。

また、\texttt{ALL(*)}アルゴリズム自体とは関係ありませんが、XMLのパーサーを書くときには開きタグと閉じタグの名前が一致している必要があります。この条件を記述するために\texttt{PetitXML}では次のように記述されています。

\begin{codescreen}{}
'<' begin=NAME '>' es=elements '</' end=NAME '>' {
  $begin.text.equals($end.text)
}?
\end{codescreen}

この中の\texttt{\{\$begin.text.equals(\$end.text)\}?}という部分はsemantic
predicate（意味的述語）と呼ばれ、プログラムとして書かれた条件式が真になるときにだけマッチします。この例では、開きタグの名前（\texttt{begin}）と閉じタグの名前（\texttt{end}）が一致しているかをJavaのコードで確認しています。semantic
predicateのような機能はプログラミング言語をそのまま埋め込むという意味で、正直「あまりきれいではない」と思わなくもないですが、実用上はsemantic
predicateを使いたくなる場面にしばしば遭遇します。

ANTLRはこういった実用上重要な痒いところにも手が届くように作られており、非常によくできた構文解析器生成系といえるでしょう。

\section{SComb：Scala製のパーサーコンビネータ}
\index{SComb}
\index{ぱーさーこんびねーた@パーサーコンビネータ}
\index{Scala}
\index{DSL}

手前味噌ですが、拙作のパーサーコンビネータである\href{https://github.com/kmizu/scomb}{SComb}も紹介しておきます。これまで紹介してきたものはすべて構文解析器生成系です。つまり、独自の言語を用いて作りたい言語の文法を記述し、そこから\textgt{対象言語}（CであったりJavaであったりさまざまですが）で書かれた構文解析器を生成するものだったわけですが、パーサーコンビネータは少々違います。

パーサーコンビネータでは対象言語のメソッドや関数、オブジェクトとして構文解析器を定義し、演算子やメソッドによって構文解析器を組み合わせることで構文解析器を組み立てていきます。「コンビネータ」という名前は「組み合わせる（combine）」からきており、小さなパーサーを組み合わせて大きなパーサーを作るという意味です。パーサーコンビネータではメソッドや関数、オブジェクトとして規則自体を記述するため、特別にプラグインを作らなくてもIDEによる支援が受けられることや、対象言語が静的型システムを持っていた場合、型チェックによる支援を受けられることなどがメリットとして挙げられます。

SCombで四則演算を解析できるプログラムを書くと以下のようになります。先ほど述べたようにSCombはパーサーコンビネータであり、これ自体がScalaのプログラム（\texttt{object}宣言）でもあります。

\begin{codescreen}{}
object Calculator extends SCombinator {
  // root <- E 
  def root: P[Int] = E

  // E <- A
  def E: P[Int] = rule(A)

  // A <- M ("+" M / "-" M)* 
  def A: P[Int] = rule(chainl(M) {
    $("+").map { op => (lhs: Int, rhs: Int) => lhs + rhs } |
    $("-").map { op => (lhs: Int, rhs: Int) => lhs - rhs }
  })

  // M <- P ("*" P / "/" P)* 
  def M: P[Int] = rule(chainl(P) {
    $("*").map { op => (lhs: Int, rhs: Int) => lhs * rhs } |
    $("/").map { op => (lhs: Int, rhs: Int) => lhs / rhs }
  })

  // P <- "(" E ")" | N
  def P: P[Int] = rule{
    (for {
      _ <- string("("); e <- E; _ <- string(")")} yield e) | N
  }
  
  // N <- [0-9]+ 
  def N: P[Int] = rule(set('0'to'9').+.map{ digits => digits.mkString.toInt})

  def parse(input: String): Result[Int] = parse(root, input)
}
\end{codescreen}

ScalaはJavaとは異なるプログラミング言語ですが、構文にJavaと似た部分が多く（JVM上で動く処理系がある）、Javaユーザーの方でも比較的読み進めやすいはずです。

各メソッドに対応するBNFによる規則をコメントとして付加してみましたが、BNFと比較しても簡潔に記述できています。Scalaは記号をそのままメソッドとして記述できるなど、元々DSL（ドメイン特化言語）に向いている特徴を持った言語なのですが、その特徴を活用しています。\texttt{chainl}というメソッドについてだけは見慣れない読者の方が多そうですが、これは

\begin{codescreen}{}
// A ::= M ("+" M | "-" M)*
\end{codescreen}

\noindent のような二項演算を簡潔に記述するためのコンビネータ（メソッド）です。パーサーコンビネータの別のメリットとして、BNF（あるいはPEG）に無い演算子を後付けで導入できることも挙げられます。構文規則からの値（意味値）の取り出しもScalaのfor式（Javaのfor文とは異なり、値を生成する式です）を用いて簡潔に記述できています。

筆者は自作言語Klassicの処理系作成のためにSCombを使っていましたが、かなり複雑な文法を記述できるにもかかわらず、SCombのコア部分はわずか600行ほどです。それでいて高い拡張性や簡潔な記述が可能なのは、Scalaという言語の能力と、SCombがベースとして利用しているPEGという手法のシンプルさがあってのものだと言えるでしょう。

\section{パーサーコンビネータJCombを自作しよう！}
\index{JComb}
\index{ぱーさーこんびねーた@パーサーコンビネータ!JComb}
\index{じぇねりくす@ジェネリクス}

コンパイラについて解説した本は数えきれないほどありますし、その中で構文解析アルゴリズムについて説明した本も少なからずあります。しかし、構文解析アルゴリズムについてのみフォーカスした本はParsing Techniquesほぼ一冊といえる現状です。その上でパーサーコンビネータの自作まで踏み込んだ書籍はほぼ皆無と言っていいでしょう。読者の方には「さすがにちょっとパーサーコンビネータの自作は無理があるのでは」と思われた方もいるのではないでしょうか。

しかし、現代的な言語であればパーサーコンビネータを自作するのは本当に簡単です。多くの読者の方々が拍子抜けしてしまうくらいに。この節ではJavaで書かれたパーサーコンビネータJCombを自作する過程を通じて皆さんにパーサーコンビネータとはどのようなものかを学んでいただきます。

パーサーコンビネータと構文解析器生成系の関係は、次のように考えられます。

\begin{itemize}
\item
  構文解析器生成系：文法定義ファイル → ツール → Javaコード
\item
  パーサーコンビネータ：Javaコードで直接文法を表現
\end{itemize}

つまり、パーサーコンビネータは「生成」のステップを省き、プログラミング言語の機能を活用して直接パーサーを組み立てる手法なのです。

まず復習になりますが、構文解析器というのは文字列を入力として受け取って、解析結果を返す関数（あるいはオブジェクト）とみなせるのでした。これはパーサーコンビネータ、特にPEGを使ったパーサーコンビネータを実装するときに有用な見方です。この「構文解析器はオブジェクトである」を文字通りとって、以下のようなジェネリックなインタフェース\texttt{JParser\textless{}R\textgreater{}}を定義します。

\begin{codescreen}{}
interface JParser<R> {
  Result<R> parse(String input);
}
\end{codescreen}

ここで構文解析器を表現するインタフェース\texttt{JParser\textless{}R\textgreater{}}は型パラメータ\texttt{R}を受け取ることに注意してください。型パラメータ\texttt{R}は、「このインタフェースは何かの型\texttt{R}に対して動作しますが、具体的な型は使用時に決まります」という意味です。

構文解析の結果は抽象構文木になりますが、インタフェースを定義する時点では抽象構文木がどのような形になるかはわかりようがないので、型パラメータにしておくのです。\texttt{JParser\textless{}R\textgreater{}}はたった1つのメソッド\texttt{parse()}を持ちます。\texttt{parse()}は入力文字列\texttt{input}を受け取り、解析結果を\texttt{Result\textless{}R\textgreater{}}として返します。

\texttt{JParser\textless{}R\textgreater{}}の実装はいったいぜんたいどのようなものになるの？という疑問を脇に置いておけば理解は難しくないでしょう。次に解析結果\texttt{Result\textless{}V\textgreater{}}をレコードとして定義します。

\begin{codescreen}{}
record Result<V>(V value, String rest){}
\end{codescreen}

レコード\texttt{Result\textless{}V\textgreater{}}は解析結果を保持するクラスです。\texttt{value}は解析結果の値を表現し、\texttt{rest}は解析した結果「残った」文字列を表します。たとえば、\texttt{"123abc"}という文字列から数値部分\texttt{"123"}を解析した場合、\texttt{value}は\texttt{123}（整数値）、\texttt{rest}は\texttt{"abc"}（残りの文字列）となります。

このインタフェース\texttt{JParser\textless{}R\textgreater{}}は次のように使えると理想的です。

\begin{codescreen}{}
JParser<Integer> calculator = ...;
Result<Integer> result = calculator.parse("1+2*3");
assert 7 == result.value();
\end{codescreen}

パーサーコンビネータは、このようなどこか都合のよい\texttt{JParser\textless{}R\textgreater{}}を、BNF（あるいはPEG）に近い文法規則を連ねていくのに近い使い勝手で構築するための技法です。前の節で紹介した\texttt{SComb}もパーサーコンビネータでしたが基本的には同じようなものです。

この節では最終的に上のような式を解析できるパーサーコンビネータを作るのが目標です。

\subsection{部品を考えよう}

これからパーサーコンビネータを作っていくわけですが、パーサーコンビネータの基本となる「部品」を作る必要があります。

まず最初に、文字列リテラルを受け取ってそれを解析できる次のような\texttt{string()}メソッドはぜひともほしいところです。

\begin{codescreen}{}
assert new Result<String>("123", "").equals(string("123").parse("123"));
\end{codescreen}

これはBNFで言えば文字列リテラルの表記に相当します。

次に、解析に成功したとしてその値を別の値に変換するための方法もほしいところです。

たとえば、\texttt{123}という文字列を解析したとして、Stringではなくintに変換したいところです。Javaのラムダ式（無名関数）を使えば、このような変換を簡潔に書けます。このようなメソッドは、ラムダ式で変換を定義できるように、次のような\texttt{map()}メソッドとして提供したいところです。

\begin{codescreen}{}
<T, U> JParser<U> map(JParser<T> parser, Function<T, U> function);
assert (new Result<Integer>(123, "")).equals(
    map(string("123"), v -> Integer.parseInt(v)).parse("123")
);
\end{codescreen}

これは構文解析器生成系でセマンティックアクションを書くのに相当すると言えます。

BNFで\texttt{a\ \textbar{}\ b}、つまり選択を書くのに相当するメソッドも必要です。これは次のような\texttt{alt()}メソッドとして提供します。

\begin{codescreen}{}
<T> JParser<T> alt(JParser<T> p1, JParser<T> p2);
assert (new Result<String>("bar", "")).equals(
    alt(string("foo"), string("bar")).parse("bar")
);
\end{codescreen}

同様に、BNFで\texttt{a\ b}、つまり連接を書くのに相当するメソッドも必要ですが、これは次のような\texttt{seq()}メソッドとして提供します。

\begin{codescreen}{}
record Pair<A, B>(A a, B b){}
<A, B> JParser<Pair<A, B>> seq(JParser<A> p1, JParser<B> p2);
assert (new Result<>(new Pair<>("foo", "bar"), "")).equals(
    seq(string("foo"), string("bar")).parse("foobar")
);
\end{codescreen}

最後に、BNFで\texttt{a*}、つまり0回以上の繰り返しに相当する\texttt{rep0()}メソッドや、

\begin{codescreen}{}
<T> JParser<List<T>> rep0(JParser<T> p);
\end{codescreen}

\texttt{a+}、つまり1回以上の繰り返しに相当する\texttt{rep1()}メソッドもほしいところです。

\begin{codescreen}{}
<T> JParser<List<T>> rep1(JParser<T> p);
assert (new Result<List<String>>(List.of("a", "a", "a"), "")).equals(
  rep1(string("a")).parse("aaa")
);
\end{codescreen}

この節ではこれらのプリミティブなメソッドの実装方法について説明していきます。

\subsection{string()メソッド}
\index{JComb}
\index{もじれつりてらる@文字列リテラル}

最初に\texttt{string(String\ literal)}メソッドで返す\texttt{JLiteralParser\textless{}String\textgreater{}}の中身を作ってみましょう。\texttt{JLiteralParser}クラスはただ1つのメソッド\texttt{parse()}を持ち、次のような実装になります。

\begin{codescreen}{}
class JLiteralParser implements JParser<String> {
  private String literal;
  public JLiteralParser(String literal) {
    this.literal = literal;

  }
  public Result<String> parse(String input) {
    if(input.startsWith(literal)) {
      return new Result<String>(literal, input.substring(literal.length()));
    } else {
      return null;
    }
  }
}
\end{codescreen}

このクラスは次のようにして使います。

\begin{codescreen}{}
assert new Result<String>("foo", "").equals(new JLiteralParser("foo").parse("foo"));
\end{codescreen}

リテラルを表すフィールド\texttt{literal}が\texttt{input}の先頭とマッチした場合、\texttt{literal}と残りの文字列からなる\texttt{Result\textless{}String\textgreater{}}を返します。そうでない場合は返すべき\texttt{Result}がないので\texttt{null}を返します。簡単ですね。

\texttt{startsWith}メソッドは、文字列がある文字列で始まるかを判定するJavaの標準メソッドです。\texttt{substring}メソッドは、文字列の一部を切り出すメソッドです。

ここでは簡潔さを優先して失敗時に\texttt{null}を返していますが、\texttt{Success}/\texttt{Failure}のような直和型にしておくと、戻り値に\texttt{null}が紛れ込む心配がなくなります。JCombでは説明を短くするため\texttt{null}にしていますが、実用では直和型を作るのがおすすめです。

あとはこのクラスのインスタンスを返す\texttt{string()}メソッドを作成するだけです。利便性のため、以降では各種メソッドはクラス\texttt{JComb}のstaticメソッドとして実装していきます。

\begin{codescreen}{}
public class JComb {
  public static JParser<String> string(String literal) {
    return new JLiteralParser(literal);
  }
}
\end{codescreen}

使うときは次のようになります。

\begin{codescreen}{}
JParser<String> foo = string("foo");
assert new Result<String>("foo", "_bar").equals(foo.parse("foo_bar"));
assert null == foo.parse("baz");
\end{codescreen}

\subsection{alt()メソッド}
\index{せんたく@選択}
\index{じゅんじょつきせんたく@順序付き選択}

2つのパーサーを取って「選択」パーサーを返すメソッド\texttt{alt()}を実装します。先ほどのようにクラスを実装してもいいですが、メソッドは1つだけなのでラムダ式（Java
8から導入された無名関数）にします。

\begin{codescreen}{}
public class JComb {
    // p1 / p2 (PEGの順序付き選択に対応)
    public static <A> JParser<A> alt(JParser<A> p1, JParser<A> p2) {
        return (input) -> {
            var result = p1.parse(input);//(1) p1を試す
            if(result != null) return result;
                                   //(2) p1が成功したらその結果を返す
            return p2.parse(input);//(3) p1が失敗したらp2を試す
        };
    }
}
\end{codescreen}

ラムダ式について復習しておきます。これは実質的に以下の匿名クラスを書いたのと同じになります。

\begin{codescreen}{}
public class JComb {
    public static <A> JParser<A> alt(JParser<A> p1, JParser<A> p2) {
        return new JAltParser<A>(p1, p2);
    }
}

class JAltParser<A> implements JParser<A> {
  private JParser<A> p1, p2;
  public JAltParser(JParser<A> p1, JParser<A> p2) {
    this.p1 = p1;
    this.p2 = p2;
  }
  public Result<A> parse(String input) {
    var result = p1.parse(input);
    if(result != null) return result;
    return p2.parse(input);
  }
}
\end{codescreen}

この定義では、

\begin{enumerate}
\item
  まずパーサー\texttt{p1}を試す
\item
  \texttt{p1}が成功した場合は\texttt{p2}を試すことなく値をそのまま返す
\item
  \texttt{p1}が失敗した場合、\texttt{p2}を試しその値を返す
\end{enumerate}

\noindent という挙動をします。これはBNFの選択 \texttt{\textbar{}}
とは異なり、PEGの「順序付き選択 \texttt{/}」に対応します。

実は今ここで作っているパーサーコンビネータである\texttt{JComb}は（\texttt{SComb}と同様に）PEGをベースとしたパーサーコンビネータだったのです。

PEGベースでないパーサーコンビネータを作ることもできるのですが実装がかなり複雑になってしまいます。PEGの挙動をそのままプログラミング言語に当てはめるのは非常に簡単であるため、今回はPEGを採用しましたが、もし興味があればBNFベース（文脈自由文法ベース）のパーサーコンビネータも作ってみてください。

\subsection{seq()メソッド}
\index{れんせつ@連接}

次に2つのパーサーを取って「連接」パーサーを返すメソッド\texttt{seq()}を実装します。これはPEGの連接
\texttt{e1\ e2} に対応します。先ほどと同じくラムダ式にしてみます。

\begin{codescreen}{}
record Pair<A, B>(A a, B b){}
// p1 p2 (PEGの連接に対応)
public class JComb {
    public static <A, B> JParser<Pair<A, B>> seq(
        JParser<A> p1, JParser<B> p2) {
        return (input) -> {
            var result1 = p1.parse(input); //(1-1) p1を試す
            if(result1 == null) return null; 
                                           //(1-2) p1が失敗したら全体も失敗
            var result2 = p2.parse(result1.rest()); 
                                           //(2-1) p1の残り入力でp2を試す
            if(result2 == null) return null; 
                                           //(2-2) p2が失敗したら全体も失敗
            // 両方成功したら、結果をペアにして返す
            return new Result<>(
                new Pair<A, B>(
                    result1.value(), result2.value()
                ), result2.rest()
            );//(2-3)
        };
    }
}
\end{codescreen}

先ほどの\texttt{alt()}メソッドと似通った実装ですが、\texttt{p1}が失敗したら全体が失敗するのがポイントです。\texttt{p1}と\texttt{p2}の両方が成功した場合は、2つの値のペアを返しています。

\subsubsection{rep0(), rep1()メソッド}
\index{くりかえし@繰り返し}

残りは0回以上の繰り返し（PEGの
\texttt{p*}）を表す\texttt{rep0()}と1回以上の繰り返し（PEGの
\texttt{p+}）を表す\texttt{rep1()}メソッドです。

まず、\texttt{rep0()}メソッドは次のようになります。

\begin{codescreen}{}
public class JComb {
    public static <A> JParser<List<A>> rep0(JParser<A> p) {
        return (input) -> {
            var result = p.parse(input); // (1)
            if(result == null) return new Result<>(List.of(), input); //(2)
            var value = result.value();
            var rest = result.rest();
            var result2 = rep0(p).parse(rest); //(3)
            if(result2 == null) return new Result<>(List.of(value), rest);
            List<A> values = new ArrayList<>();
            values.add(value);
            values.addAll(result2.value());
            return new Result<>(values, result2.rest());
        };
    }
}
\end{codescreen}

\texttt{(1)}でまずパーサー\texttt{p}を適用しています。ここで失敗した場合、0回の繰り返しにマッチしたことになるので、空リストからなる結果を返します（\texttt{(2)}）。そうでなければ、1回以上の繰り返しにマッチしたことになるので、繰り返し同じ処理をする必要がありますが、これは再帰呼出しによって簡単に実装できます（\texttt{(3)}）。シンプルな実装ですね。

\texttt{rep1(p)}は意味的には\texttt{seq(p,\ rep0(p))}なので、次のようにして実装を簡略化できます。

\begin{codescreen}{}
public class JComb {
    public static <A> JParser<List<A>> rep1(JParser<A> p) {
        JParser<Pair<A, List<A>>> rep1Sugar = seq(p, rep0(p));
        return (input) -> {
            var result = rep1Sugar.parse(input);//(1)
            if(result == null) return null;//(2)
            var pairValue = result.value(); // resultから値を取得
            var values = new ArrayList<A>();
            values.add(pairValue.a()); // Pairの最初の要素
            values.addAll(pairValue.b()); // Pairの2番目の要素 (List)
            return new Result<>(values, result.rest()); //(3)
        };
    }
}
\end{codescreen}

実質的な本体は\texttt{(1)}だけで、あとは結果の値が\texttt{Pair}なのを\texttt{List}に加工しているだけですね。

\subsection{map()メソッド}
\index{せまんてぃっくあくしょん@セマンティックアクション}
\index{らむだしき@ラムダ式}

パーサーを加工して別の値を生成するためのメソッド\texttt{map()}を実装してみましょう。\texttt{map()}は\texttt{JParser\textless{}R\textgreater{}}のメソッドとして実装するとメソッドチェインが使えて便利なので、インタフェースの\texttt{default}メソッドとして実装します。

\begin{codescreen}{}
interface JParser<R> {
    Result<R> parse(String input);

    default <T> JParser<T> map(Function<R, T> f) {
        return (input) -> {
           var result = this.parse(input);
           if (result == null) return null;
           return new Result<>(f.apply(result.value()),result.rest());//(1)
        };
    }
}
\end{codescreen}

\texttt{(1)}で\texttt{f.apply(result.value())}として値を加工しているのがポイントです。

\subsection{lazy()メソッド}
\index{ちえんひょうか@遅延評価}
\index{むげんさいき@無限再帰}
\index{そうごさいき@相互再帰}

パーサーを遅延評価するためのメソッド\texttt{lazy()}も導入します。

\begin{codescreen}{}
public class JComb {
    public static <A> JParser<A> lazy(Supplier<JParser<A>> supplier) {
        return (input) -> supplier.get().parse(input);
    }
}
\end{codescreen}

\texttt{lazy}メソッドの必要性について補足します。Javaは先行評価（関数が呼ばれたらすぐに評価）を行う言語であるため、再帰的な文法規則を直接メソッド呼び出しで表現しようとすると、パーサーオブジェクトの構築時に無限再帰が発生し
\texttt{StackOverflowError} となることがあります。たとえば、算術式の文法で
\texttt{expression} が \texttt{additive} を呼び出し、\texttt{additive}
が \texttt{primary} を呼び出し、\texttt{primary} が括弧表現の中で再び
\texttt{expression} を呼び出すような相互再帰構造を考えてみましょう。

\begin{codescreen}{}
// JComb を使った算術式のパーサー定義（簡略版・lazyなしのイメージ）
// 仮に直接代入しようとすると...
public static JParser<Integer> expression() {
    return additive(); 
}
public static JParser<Integer> additive() {
    return primary();   
}
public static JParser<Integer> primary() {
    return alt(
        number,
        seq(string("("), expression(), string(")"))
    );
}
\end{codescreen}

上記のように単純にメソッド呼び出しでパーサーを組み合わせようとすると、\texttt{expression}
の初期化時に \texttt{additive} が必要になり、その \texttt{additive}
の初期化に \texttt{primary} が、さらにその \texttt{primary} の初期化に
\texttt{expression}
が必要となり、循環参照によって初期化が終わらなくなります。

\texttt{lazy} は
\texttt{Supplier\textless{}JParser\textless{}A\textgreater{}\textgreater{}}
を引数に取ることで、実際の \texttt{JParser\textless{}A\textgreater{}}
オブジェクトの取得（\texttt{supplier.get()}）を、そのパーサーが実際に
\texttt{parse()}
メソッドで使われるときまで遅延させます。\texttt{Supplier}はJavaの関数型インタフェースで、「引数なしで値を返す関数」を表します。これにより、相互再帰するパーサー定義でも、オブジェクト構築時の無限再帰を避けられます。算術式の例では、\texttt{expression}
の定義内で \texttt{additive} を呼び出す部分を
\texttt{lazy(()\ -\textgreater{}\ additive())}
のように記述することで、この問題を解決します。

\subsection{regex()メソッド}
\index{せいきひょうげん@正規表現}

せっかくなので正規表現を扱うメソッド\texttt{regex()}も導入してみましょう。

\begin{codescreen}{}
public class JComb {
    public static JParser<String> regex(String regex) {
        return (input) -> {
            var matcher = Pattern.compile(regex).matcher(input);
            if(matcher.lookingAt()) {
                return new Result<>(
                    matcher.group(), input.substring(matcher.end())
                );
            } else {
                return null;
            }
        };
    }
}
\end{codescreen}

引数として与えられた文字列を\texttt{Pattern.compile()}で正規表現に変換して、マッチングを行うだけです。\texttt{matcher.lookingAt()}は文字列の先頭から正規表現にマッチする部分があるかを確認するメソッドです。これは次のようにして使うことができます。

\begin{codescreen}{}
var number = regex("[0-9]+").map(v -> Integer.parseInt(v));
assert (new Result<Integer>(10, "")).equals(number.parse("10"));
\end{codescreen}

\subsection{算術式のインタプリタを書いてみる}

ここまでで作ったクラス\texttt{JComb}と\texttt{JParser}などを使って簡単な算術式のインタプリタを書いてみましょう。仕様は次のとおりです。

\begin{itemize}
\item
  扱える数値は整数のみ
\item
  演算子は加減乗除（\texttt{+\textbar{}-\textbar{}*\textbar{}/}）のみ
\item
  \texttt{()}によるグルーピングができる
\end{itemize}

実装だけを提示すると次のようになります。

\begin{codescreen}{}
public class Calculator {
    // expression は加減算を担当 (左結合)
    // PEG: expression <- multitive ( ( "+" / "-" ) multitive )*
    public static JParser<Integer> expression() {
        return seq( // multitive と (( "+" / "-" ) multitive )* の連接
            lazy(() -> multitive()), // 左辺の multitive (乗除の項)
            rep0( // 0回以上の繰り返し
                seq( // ( "+" / "-" ) と multitive の連接
                    // "+" または "-"
                    alt(string("+"), string("-")), 
                    lazy(() -> multitive()) // 右辺の multitive
                )
            )
        ).map(p -> { // 解析結果を処理するラムダ式
            var left = p.a(); // 初期値 (最初の項)
            var rights = p.b(); // 残りの演算子と項のペアのリスト
            for (var rightPair : rights) { 
                                // 各 Pair<String, Integer> について
                var op = rightPair.a(); // 演算子 ( "+" または "-" )
                var rightValue = rightPair.b(); // 項の値
                if (op.equals("+")) {
                    left += rightValue;
                } else { // op.equals("-")
                    left -= rightValue;
                }
            }
            return left; // 計算結果
        });
    }

    // multitive は乗除算を担当 (左結合)
    // PEG: multitive <- primary ( ( "*" / "/" ) primary )*
    public static JParser<Integer> multitive() {
        return seq( // primary と ( ( "*" / "/" ) primary )* の連接
            lazy(() -> primary()), // 左辺の primary
            rep0( // 0回以上の繰り返し
                seq( // ( "*" / "/" ) と primary の連接
                    alt(string("*"), string("/")), // "*" または "/"
                    lazy(() -> primary()) // 右辺の primary
                )
            )
        ).map(p -> { // 解析結果を処理するラムダ式
            var left = p.a(); // 初期値 (最初の因子)
            var rights = p.b(); // 残りの演算子と因子のペアのリスト
            for (var rightPair : rights) {
                var op = rightPair.a(); // 演算子 ( "*" または "/" )
                var rightValue = rightPair.b(); // 因子の値
                if (op.equals("*")) {
                    left *= rightValue;
                } else { // op.equals("/")
                    if (rightValue == 0) 
                        throw new ArithmeticException("Division by zero"); 
                    left /= rightValue;
                }
            }
            return left; // 計算結果
        });
    }

    // primary <- number / "(" expression ")"
    public static JParser<Integer> primary() {
        return alt( // number または "(" expression ")" の選択
            number, // 数値パーサー
            seq( // "(" expression の連接
                string("("), // 開き括弧
                // 括弧内の式 (expressionを再帰呼び出し)
                seq(
                    lazy(() -> expression()),
                    string(")") // 閉じ括弧
                )
            ).map(p -> p.b().a()) // ((, (expression, ))) の中央の値を
                                  // 取り出す
        );
    }
    
    // number <- [0-9]+ (PEGの正規表現リテラルに対応)
    private static JParser<Integer> number = 
        regex("[0-9]+").map(Integer::parseInt);
}
\end{codescreen}

表記は冗長なもののほぼPEGに一対一に対応しているのがわかるのではないでしょうか？

これに対してJUnitを使って以下のようなテストコードを記述してみます。無事、意図通りに解釈されていることがわかります。

\begin{codescreen}{}
assertEquals(
    new Result<>(7, ""), Calculator.expression().parse("1+2*3")
); // テストをパス (実際には multitive で処理される)
assertEquals(
    new Result<>(0, ""), Calculator.expression().parse("1+2-3")
); // テストをパス
\end{codescreen}

\subsection{自作パーサーコンビネータのススメ}

DSL（ドメイン特化言語）に向いたScalaに比べれば冗長になったものの、手書きで再帰下降パーサーを組み立てるのに比べると大幅に簡潔な記述を実現できました。しかも、JComb全体を通しても500行にすら満たないのは特筆すべきところです。Javaがユーザー定義の中置演算子（\texttt{+}や\texttt{*}のような演算子を自分で定義できる機能）をサポートしていればもっと簡潔にできたのですが、そこは向き不向きといったところでしょうか。

手書きでパーサーを書く必要に迫られた場合でも、パーサーコンビネータを自作すれば、気軽にパーサーを組み立てるためのDSL（Domain Specific Language）を定義できるのです。それだけでなく、Javaのような静的型付き言語を使った場合、IDEによる支援も受けられますし、BNFやPEGにはない便利な演算子を自分で導入することもできます。

パーサーコンビネータはお手軽なだけあって各種プログラミング言語に実装されています。たとえば、Java用なら\href{https://github.com/jparsec/jparsec}{jparsec}があります。しかし、筆者としては、パーサーコンビネータが動作する仕組みを理解するために、ぜひとも\textgt{自分だけの}パーサーコンビネータを実装してみてほしいと思います。

\section{まとめ}

この章では、構文解析器生成系（パーサージェネレータ）という、文法定義から構文解析器を自動生成するソフトウェアについて学びました。

まず、PEGから構文解析器を生成する基本的な仕組みを、Dyck言語を例に具体的に見ていきました。文法規則を機械的にJavaコードに変換することで、手書きの煩雑さから解放され、文法定義に集中できることを確認しました。

次に、代表的な構文解析器生成系として以下を紹介しました。

\begin{itemize}
\item
  \textbf{\sffamily JavaCC}：\texttt{LL(k)}法を採用し、Javaに特化した構文がJavaユーザーに親しみやすい
\item
  \textbf{\sffamily Yacc}：\texttt{LALR(1)}法を採用する歴史ある構文解析器生成系で、C言語との親和性が高い
\item
  \textbf{\sffamily ANTLR}：\texttt{ALL(*)}アルゴリズムにより強力な解析能力を持ち、左再帰も扱える多言語対応ツール
\end{itemize}

最後に、パーサーコンビネータという、プログラミング言語の機能を活用して直接パーサーを組み立てる手法を学び、実際にJavaでパーサーコンビネータ「JComb」を実装しました。わずか500行に満たないコードで、PEGに対応した実用的なパーサーコンビネータを作れることを示しました。

構文解析器生成系とパーサーコンビネータは、それぞれ異なる強みを持ちます。構文解析器生成系は宣言的で可読性が高く、パーサーコンビネータはプログラミング言語の型システムやIDEの支援を受けられ、柔軟な拡張が可能です。どちらを選ぶかは、プロジェクトの要件や開発チームの好みによって決めることになるでしょう。

\chapter{現実の構文解析}\label{ch06}
\index{げんじつのこうぶんかいせき@現実の構文解析}
\index{ぶんみゃくじゆうげんご@文脈自由言語}
\index{ぶんみゃくいぞんせい@文脈依存性}
\index{えらーりかばり@エラーリカバリ}

LL法やLR法、Packrat Parsingといった、これまでに知られているメジャーな構文解析アルゴリズムを一通り取り上げてきました。これらの構文解析アルゴリズムは文脈自由言語あるいはそのサブセットを取り扱うことができ、一般的なプログラミング言語の構文解析を行うのに必要十分な能力を持っています。

しかし、構文解析を専門としている人や実用的な構文解析器を書いている人は直感的に理解していることなのですが、実のところ、既存の構文解析アルゴリズムだけではうまく取り扱えない類の構文があります。一言でいうと、それらの構文は文脈自由言語から逸脱しているために、文脈自由言語を取り扱う既存の手法だけではうまくいかないのです。

このような既存の構文解析アルゴリズムだけでは扱えない要素は多数あります。たとえば、Cのtypedefはその典型ですし、RubyやPerlのヒアドキュメントと呼ばれる構文もそうです。他には、Scalaのプレースホルダ構文やC++のテンプレート、Pythonのインデント文法など、文脈自由言語を逸脱しているがゆえに人間が特別に配慮しなければいけない構文は多く見かけられます。

また、これまでの章では、主に構文解析を行う手法を取り扱っていましたが、現実問題としては抽象構文木をうまく作る方法やエラーメッセージを適切に出す方法も重要になってきます。

この章では、巷の書籍ではあまり扱われない、しかし現実の構文解析では対処しなくてはならない構文や問題について取り上げます。皆さんが何かしらの構文解析器を作るとき、やはり理想どおりにはいかないことが多いと思います。そのような現実の構文解析の\textgt{泥臭さ}に遭遇する読者の方々の助けになれば幸いです。

\section{トークンが「文脈自由」なケース}
\index{とーくん@トークン}
\index{じくかいせき@字句解析}
\index{もじれつほかん@文字列補間}
\index{String Interpolation}
\index{Dyck言語}
\index{PEG}
\index{すきゃなれす@スキャナレス構文解析}

最近の多くの言語は文字列補間(String
Interpolation)と呼ばれる機能を持っています。

たとえば、Rubyでは以下の文字列を評価すると、\texttt{"x\ +\ y\ =\ x\ +\ y"}ではなく\texttt{"x\ +\ y\ =\ 3"}になります。

\begin{codescreen}{}
x = 1
y = 2
"x + y = #{x + y}" # "x + y = 3"
\end{codescreen}

\texttt{\#\{}と\texttt{\}}で囲まれた範囲をRubyの式として評価した結果を文字列として埋め込んでくれるわけです。

Scalaでも同じことを次のように書けます。

\begin{codescreen}{}
val x = 1; val y = 2
s"x + y = ${x + y}" // "x + y = 3"
\end{codescreen}

Swiftだと次のようになります。

\begin{codescreen}{}
let x = 1
let y = 2
"x + y = \(x + y)"
\end{codescreen}

同様の機能はKotlin、Python、JavaScript、TypeScriptなどさまざまな言語に採用されています。比較的新しい言語や、既存言語の新機能として採用するのが普通になった機能と言えるでしょう。

文字列補間は便利な機能ですが、構文解析という観点からは少々やっかいな存在です。文字列リテラルは従来、字句解析で扱うトークンでした。そして、字句解析では処理を単純化するためにトークンを正規言語の範囲に収まるようにするのがベストプラクティスでした。字句解析と構文解析を分離し、かつ、字句解析を可能な限り単純化するという観点で言えばある意味当然とも言えますが、文字列補間は正規表現で表現できたもの（正規言語に収まるもの）を文脈自由文法で取り扱わなければいけない存在にしてしまいました。

たとえば、Rubyでは以下のように\texttt{\#\{\}}の中にさらに文字列リテラルを書くことができ、その中には\texttt{\#\{\}}を……といった具合に無限にネストできるのです。これまでの章を振り返ればわかるように、これはDyck言語的な構造があり、明らかに正規言語では扱えず、文脈自由言語の能力が必要なものです。

\begin{codescreen}{}
x = 1
y = 2
"expr1 (#{"expr2 (#{x + y})"})" # "expr1 (expr2 (3))"
\end{codescreen}

しかし、従来の手法では文字列リテラルはトークンとして扱わなければいけないため、各言語処理系の実装者はad
hocな形で構文解析器に手を加えています。たとえば、Rubyの構文解析器はYaccを使って書かれていますが、字句解析器に状態を持たせることでこの問題に対処しています。文字列リテラル内に\texttt{\#\{}が出現したら状態を式モードに切り替えて、その中で文字列リテラルがあらわれたら文字列モードに切り替えるといった具合です。

一方、PEGでは字句解析と構文解析が分離されていないため、特別な工夫をすることなく文字列補間を実装することができます。以下はRubyの文字列補間と同じようなものをPEGで記述する例です。

\begin{codescreen}{}
string <- "\"" ("#{" expression "}" / !"\"" .)* "\""
expression <- 式の定義
\end{codescreen}

文字列補間を含む文字列リテラルは分解可能という意味で厳密な意味ではトークンと言えないわけですが、PEGは字句解析を分離しないおかげで文字列リテラルをことさら特別扱いする必要がないわけです。

PEGの利用例が近年増えてきているのは、言語に対してこのようにアドホックに構文を追加したいというニーズがあるためではないかと筆者は考えています。

\begin{center}
\begin{tikzpicture}[
  node distance=3cm,
  every node/.style={font=\small},
  % 状態ノードのスタイル
  state/.style={
    circle,
    draw,
    minimum size=1.8cm,
    font=\footnotesize
  },
  % 初期状態のスタイル
  initial state/.style={
    state,
    fill=green!10
  },
  % 遷移ラベルのスタイル
  transition/.style={
    font=\footnotesize\ttfamily,
    auto
  }
]

  % 状態ノード
  \node[initial state] (string) {文字列モード};
  \node[state, right=of string, fill=orange!10] (expr) {式モード};

  % 初期状態への矢印
  \draw[->, thick] ([xshift=-1cm]string.west) -- (string);

  % 状態遷移
  \draw[->, thick, bend left=25] (string) to node[transition, above] {\#\{} (expr);
  \draw[->, thick, bend left=25] (expr) to node[transition, below] {\}} (string);

  % 自己ループ
  \draw[->, thick] (string) to[loop above] node[transition] {通常文字} (string);
  \draw[->, thick] (expr) to[loop above] node[transition] {式トークン} (expr);

  % 例の説明
  \node[below=2cm of string, text width=9cm, align=center] {
    \footnotesize 例: \texttt{"abc\#\{1+2\}def"} の解析\\[0.2em]
    \begin{tabular}{cl}
      1. & 文字列モード: \texttt{"abc"} を消費 \\
      2. & \texttt{\#\{} を認識 → 式モードへ遷移 \\
      3. & 式モード: \texttt{1+2} を式として解析 \\
      4. & \texttt{\}} を認識 → 文字列モードへ復帰 \\
      5. & 文字列モード: \texttt{def"} を消費
    \end{tabular}
  };

  % 図のラベル
  \node[below=5cm of string, font=\footnotesize] {図6.1: 文字列補間の解析状態遷移};
\end{tikzpicture}
\end{center}

\section{インデント文法}
\index{いんでんとぶんぽう@インデント文法}
\index{Python}
\index{YAML}
\index{Scala 3}
\index{INDENT@\texttt{INDENT}}
\index{DEDENT@\texttt{DEDENT}}
\index{いんでんとすたっく@インデントスタック}

Pythonではインデントによってプログラムの構造を表現します。たとえば、次のPythonプログラムを考えます。

\begin{codescreen}{}
class Point:
  def __init__(self, x, y):
    self.x = x
    self.y = y
\end{codescreen}

このPythonプログラムは次のような抽象構文木に変換されると考えられます。

\begin{center}
\begin{tikzpicture}[
  scale=0.8, transform shape,
  node distance=1.2cm and 1.5cm,
  every node/.style={font=\small},
  % ノードスタイル定義
  astnode/.style={
    rectangle,
    draw,
    rounded corners=3pt,
    minimum width=2.2cm,
    minimum height=0.7cm,
    font=\ttfamily
  },
  classnode/.style={astnode, fill=blue!15},
  defnode/.style={astnode, fill=orange!15},
  namenode/.style={astnode, fill=green!15, minimum width=3cm},
  argsnode/.style={astnode, fill=purple!15},
  bodynode/.style={astnode, fill=red!15},
  paramnode/.style={astnode, fill=gray!10, minimum width=1.2cm, minimum height=0.5cm},
  stmtnode/.style={astnode, fill=yellow!10, minimum width=2.8cm, minimum height=0.5cm},
  edge/.style={->, >=stealth, thick}
]

  % ルートノード
  \node[classnode] (class) {class};

  % 第2レベル
  \node[namenode, below left=of class] (classname) {name: Point};
  \node[defnode, below right=of class] (def) {def};

  % 第3レベル（def の子ノード）
  \node[namenode, below=of def] (funcname) {name: \_\_init\_\_};
  \node[argsnode, left=of funcname] (args) {arguments};
  \node[bodynode, right=of funcname] (body) {body};

  % 第4レベル（パラメータ）
  \node[paramnode, below=0.8cm of args, xshift=-1.2cm] (self) {self};
  \node[paramnode, below=0.8cm of args] (x) {x};
  \node[paramnode, below=0.8cm of args, xshift=1.2cm] (y) {y};

  % 第4レベル（ステートメント）
  \node[stmtnode, below=0.8cm of body, xshift=-1.5cm] (stmt1) {self.x = x};
  \node[stmtnode, below=0.8cm of body, xshift=1.5cm] (stmt2) {self.y = y};

  % エッジの描画
  \draw[edge] (class) -- (classname);
  \draw[edge] (class) -- (def);
  \draw[edge] (def) -- (funcname);
  \draw[edge] (def) -- (args);
  \draw[edge] (def) -- (body);
  \draw[edge] (args) -- (self);
  \draw[edge] (args) -- (x);
  \draw[edge] (args) -- (y);
  \draw[edge] (body) -- (stmt1);
  \draw[edge] (body) -- (stmt2);
\end{tikzpicture}
\end{center}

インデントによってプログラムの構造を表現するというアイデアは秀逸ですが、インデントによる構造の表現は明らかに文脈自由言語を超えます。インデントによってブロックや文法の構造を表現する手法は、HaskellやScala 3、YAMLなどでも採用されており、広く使われているアイデアです。

Pythonでは字句解析のときにインデントを\texttt{\textless{}INDENT\textgreater{}}（インデントレベル増加）、インデントを「外す」のを\texttt{\textless{}DEDENT\textgreater{}}（インデントレベル減少）という特別なトークンに変換します。これにより、構文解析器自体はこれらのトークンをブロックの開始と終了のように扱うことができ、文脈自由文法の範囲で処理しやすくなります。

たとえば、先の \texttt{Point}
クラスの定義は、字句解析後には（簡略化すると）以下のようなトークン列として扱われるイメージです。

\begin{codescreen}{}
<CLASS> <NAME:Point> <COLON> <NEWLINE>
<INDENT>
  <DEF> <NAME:__init__>
  <LPAREN> <NAME:self> <COMMA> <NAME:x> <COMMA> <NAME:y> <RPAREN>
  <COLON> <NEWLINE>
  <INDENT>
    <NAME:self> <DOT> <NAME:x> <ASSIGN> <NAME:x> <NEWLINE>
    <NAME:self> <DOT> <NAME:y> <ASSIGN> <NAME:y> <NEWLINE>
  <DEDENT>
<DEDENT>
\end{codescreen}

\texttt{\textless{}NAME:A\textgreater{}}は識別子\texttt{A}、\texttt{\textless{}COLON\textgreater{}}はコロン、\texttt{\textless{}NEWLINE\textgreater{}}は改行、\texttt{\textless{}ASSIGN\textgreater{}}は代入演算子を表すトークンとします。実際にはさらに詳細なトークン分割が行われます。このように、インデント/デデントトークンがブロック構造を示すため、構文解析器は括弧の対応付けと同様の形で処理できます。

しかし、よくよく考えればわかるのですが、\texttt{\textless{}INDENT\textgreater{}}トークンと\texttt{\textless{}DEDENT\textgreater{}}トークンを切り出す時点で、文脈自由ではありません。つまり、字句解析時に\texttt{\textless{}INDENT\textgreater{}}と\texttt{\textless{}DEDENT\textgreater{}}トークンを切り出すために特殊な処理をしていることになります。\texttt{\textless{}DEDENT\textgreater{}}トークンは\texttt{\textless{}INDENT\textgreater{}}トークンとスペースの数が対応していなければいけないため、切り出すためには正規表現でも文脈自由文法でも手に余ります。

\begin{center}
\begin{tikzpicture}[
  node distance=0.3cm,
  every node/.style={font=\footnotesize},
  % コードラインのスタイル
  codeline/.style={
    rectangle,
    anchor=west,
    font=\ttfamily\footnotesize
  },
  % スタック表示のスタイル
  stacknode/.style={
    rectangle,
    draw,
    minimum width=0.8cm,
    minimum height=0.5cm,
    font=\footnotesize
  }
]

  % 左側: ソースコード表示
  \node[codeline] (line1) at (0,0) {class Point:};
  \node[codeline] (line2) at (0.5,-0.6) {def \_\_init\_\_(self, x, y):};
  \node[codeline] (line3) at (1,-1.2) {self.x = x};
  \node[codeline] (line4) at (1,-1.8) {self.y = y};

  % 右側: インデントスタックの状態
  \node[right=5.5cm of line1, anchor=west] (stack_label) {\textgt{インデントスタック}};

  % スタック状態1: [0]
  \node[stacknode, below=0.3cm of stack_label] (s1_0) {0};
  \node[right=0.2cm of s1_0] {← 初期状態};

  % スタック状態2: [0, 2]
  \node[stacknode, below=0.7cm of s1_0] (s2_0) {0};
  \node[stacknode, right=0 of s2_0] (s2_2) {2};
  \node[right=0.2cm of s2_2] {← class定義後};

  % スタック状態3: [0, 2, 4]
  \node[stacknode, below=0.7cm of s2_0] (s3_0) {0};
  \node[stacknode, right=0 of s3_0] (s3_2) {2};
  \node[stacknode, right=0 of s3_2] (s3_4) {4};
  \node[right=0.2cm of s3_4] {← def定義後};

  % スタック状態4: [0, 2]
  \node[stacknode, below=0.7cm of s3_0] (s4_0) {0};
  \node[stacknode, right=0 of s4_0] (s4_2) {2};
  \node[right=0.2cm of s4_2] {← defブロック終了};

  % スタック状態5: [0]
  \node[stacknode, below=0.7cm of s4_0] (s5_0) {0};
  \node[right=0.2cm of s5_0] {← classブロック終了};
\end{tikzpicture}
\end{center}

スタックを使っていることから一見文脈自由文法で処理できるように思えるかもしれませんが、スタックの各要素が数字であり、この数字が対応していない時点でエラーになるので、文脈自由文法ではありません。Pythonの字句解析器はこのようなインデント/デデントトークンをプログラムで処理していますが、その字句解析器は明らかに文脈自由でない言語を処理しています。

いずれにせよ、インデント文法が文脈自由言語に収まらないために、字句解析器あるいは構文解析器のどちらかにしわ寄せがきます。インデント文法を処理できるもっと強力な形式文法はありますが、プログラムでad hocに対応する方が早いため、そのような方式が一般的です。

\section{ヒアドキュメント}
\index{ひあどきゅめんと@ヒアドキュメント}
\index{Ruby}
\index{Perl}
\index{XML}
\index{HTML}
\index{たぐたいおう@タグ対応}
\index{もじれつほかん@文字列補間}

ヒアドキュメントは複数行に渡る文字列を記述するための文法で、従来はbashなどのシェル言語で採用されていましたが、PerlやRubyはヒアドキュメントを採用しました。たとえば、RubyでHTMLの文字列をヒアドキュメントで以下のように書けます。

\begin{codescreen}{}
html = <<HTML
<html>
  <head><title>Title</title></head>
  <body><p>Hello</p></body>
</html>
HTML
\end{codescreen}

特筆すべきは、\texttt{\textless{}\textless{}HTML}と\texttt{HTML}のように対応している間だけが文字列として解釈されることです。これだけなら文脈自由言語の範囲内です。実際には問題はもっと複雑です。ヒアドキュメントは\textgt{ネストが可能}なのです。たとえば、以下のようなヒアドキュメントは正しいRubyプログラムです。

\begin{codescreen}{}
here = <<E1 + <<E2
ここはE1です
E1
ここはE2です
E2
\end{codescreen}

これは以下の内容の文字列として解釈されます。

\begin{codescreen}{}
ここはE1です
ここはE2です
\end{codescreen}

ヒアドキュメント内では文字列補間が使えるのでさらに複雑です。以下のようなヒアドキュメントもOKなのです。

\begin{codescreen}{}
a = 100
b = 200
here = <<A + <<B
aは#{a}です
A
bは#{b}です
B
\end{codescreen}

これは次の文字列として解釈されます。

\begin{codescreen}{}
aは100です
bは200です
\end{codescreen}

読者の方々はおそらく「確かに凄いけど、普通はこのような書き方をすることはほぼないのでは」と思われたのではないでしょうか。実際問題そうなのですが、Rubyはこのような複雑怪奇なプログラムもうまく構文解析できなければいけないのも事実です。

Rubyのヒアドキュメントを適切に構文解析するにはHTMLやXMLにおけるタグ名の対応付けと同じ処理が必要になりますが、これは文脈自由言語の範囲を超えています。Rubyのヒアドキュメントがこのような振る舞いをすることを初めて知ったのは筆者が大学院生の頃ですが、あまりに予想外の振る舞いに目眩がする思いだったのを覚えています。

Rubyのヒアドキュメントが実際にどのように実装されているかはさておき、筆者はかつて中田育男先生が実装されたISO
RubyのScalaによる試験的な構文解析器の実装のお手伝いをした際に、ヒアドキュメントの扱いに非常に苦労しました。

実装の詳細は、中田先生のruby\_scalaリポジトリ\footnote{\url{https://github.com/inakata/ruby\_scala/blob/3f54cc6f80678e30a211fb1374280246f08182ed/src/main/scala/com/github/inakata/ruby\_scala/Ruby.scala\#L1383}}で確認できますが、ヒアドキュメントの開始デリミタを記憶し、対応する終了デリミタが現れるまでを特別に処理する、といった複雑なロジックが必要でした。

このときはScalaのパーサーコンビネータを使ってヒアドキュメントを再現したのですが、引数を取ってコンビネータを返すメソッドを定義することで問題を解決しました。形式言語の文脈でいうのなら、PEGの規則が引数を持てるように拡張することでヒアドキュメントを解釈できるようになったと言えます。

PEGを拡張して規則が引数を持てるようにするという試みは複数ありますが、筆者もMacro
PEGというPEGを拡張したものを提案しました\footnote{実装と仕様は\url{https://github.com/macro-peg/macro_peg}
  を参照してください。}。ヒアドキュメントという当たり前に使われている言語機能ですら、構文解析を正しく行うためには厄介な処理をする必要があるのです。

\usetikzlibrary{decorations.pathreplacing}
\begin{center}
\begin{tikzpicture}[
  node distance=0.4cm,
  every node/.style={font=\footnotesize},
  % ヒアドキュメントのスタイル
  heredoc/.style={
    rectangle,
    anchor=west,
    font=\ttfamily\footnotesize
  },
  % デリミタのスタイル
  delimiter/.style={
    rectangle,
    draw,
    rounded corners=2pt,
    fill=yellow!20,
    font=\ttfamily\footnotesize,
    minimum height=0.5cm,
    minimum width=1cm
  },
  % スタック要素のスタイル
  stackitem/.style={
    rectangle,
    draw,
    fill=orange!10,
    minimum width=1cm,
    minimum height=0.5cm,
    font=\footnotesize\ttfamily
  },
  % ブレースのスタイル
  brace/.style={
    decorate,
    decoration={brace, amplitude=5pt,mirror},
    thick
  }
]

  % 左側: ヒアドキュメントの構造
  \node[heredoc] (line1) at (0,0) {<<E1};
  \node[heredoc] (line2) at (0.3,-0.5) {...};
  \node[heredoc] (line3) at (0.3,-1) {<<E2};
  \node[heredoc] (line4) at (0.6,-1.5) {...};
  \node[heredoc] (line5) at (0.3,-2) {E2};
  \node[heredoc] (line6) at (0.3,-2.5) {...};
  \node[heredoc] (line7) at (0,-3) {E1};

  % デリミタのハイライト
  \node[delimiter, right=2cm of line1] (delim1_start) {E1開始};
  \node[delimiter, right=2cm of line3] (delim2_start) {E2開始};
  \node[delimiter, right=2cm of line5] (delim2_end) {E2終了};
  \node[delimiter, right=2cm of line7] (delim1_end) {E1終了};

  % 矢印
  \draw[->, gray] (line1.east) -- (delim1_start.west);
  \draw[->, gray] (line3.east) -- (delim2_start.west);
  \draw[->, gray] (line5.east) -- (delim2_end.west);
  \draw[->, gray] (line7.east) -- (delim1_end.west);

  % ブレースで範囲を示す
  \draw[brace] ([xshift=-9mm]line1.west) -- ([xshift=-9mm]line7.west)
    node[midway, left=1mm, text width=1cm, align=center] {E1の\\範囲};
  \draw[brace] ([xshift=-1mm]line3.west) -- ([xshift=-1mm]line5.west)
    node[midway, left=0.5mm, text width=1cm, align=center] {E2の\\範囲};
\end{tikzpicture}
\end{center}

\section{改行終端可能文法}
\index{かいぎょうしゅうたんかのうぶんぽう@改行終端可能文法}
\index{かいぎょう@改行}
\index{Scala}
\index{Kotlin}
\index{nl@\texttt{nl}トークン}

C、C++、Java、C\#などの言語では、解釈・実行の基本単位は\textgt{文}(Statement)と呼ばれるものになります。また、文はセミコロンなどの終端子と呼ばれるもので終わるか、区切り文字で区切られるのが一般的です。セミコロンが終端子でなく文の区切りになるのがPascalなどの言語です。厳密には違いますが、関数型プログラミング言語Standard
MLのセミコロンも似たような扱いです。

Javaでは次のように書くことで、A、B、Cを順に出力できます。

\begin{codescreen}{}
System.out.println("A");
System.out.println("B");
System.out.println("C");
\end{codescreen}

このように文が終端子（Terminator）で終わる文法には、文の途中に改行が挟まっても単なるスペースと同様に取り扱えるという利点があります。先ほどのプログラムを次のように書き換えても意味は変わりません。

\begin{codescreen}{}
System.out.println(
  "A");
System.out.println(
  "B");
System.out.println(
  "C");
\end{codescreen}

大抵の場合、文は一行で終わるのですから、毎回セミコロンをつけなければいけないのも面倒くさいものです。

そういったニーズを反映してか、Scala、Kotlin、Swift、Goなどの比較的新しい言語では（Scalaの初期バージョンが2003ですから、新しいかという話もありますが）、文はセミコロンで終わることもできるが、改行でも終われます。より古い言語であってもPython、Ruby、JavaScriptなども改行で文を終われます。

たとえば、先ほどのJavaプログラムに相当するScalaプログラムは次のようになります。

\begin{codescreen}{}
println("A")
println("B")
println("C")
\end{codescreen}

見た目にもすっきりしますし、文が改行で終わるコーディングスタイルが大半であることを考えても、無駄なタイピングが減るしでいいことずくめです。それでいて、次のように文の途中で改行が入っても問題なく解釈・実行できます。

\begin{codescreen}{}
println(
  "A")
println(
  "B")
println(
  "C")
\end{codescreen}

Scalaでも一行の文字数が増えれば分割したくなりますから、このような機能があるのは自然でしょう。

ここで一つの疑問が湧きます。「文は改行で終わる」という規則なら改行がきたときに「文の終わり」とみなせばよいですし、「文はセミコロンで終わる」という規則なら、セミコロンがきたときに「文の終わり」とみなせば問題ありません。しかしながら、Scalaのような文法を実現するためには「セミコロンがくれば文が終わるが、改行で文が終わることもある」というややこしい規則に基づいて構文解析をしなければいけません。

このような文法を実現するのは案外ややこしいものです。Javaの\texttt{System.out.println("A");}という文は正確には「式文」と呼ばれますが、この式文は次のように定義されます。

\begin{lstlisting}
expression_statement = expression <SEMICOLON>
\end{lstlisting}

\texttt{\textless{}SEMICOLON\textgreater{}}はセミコロンを表すトークンです。では、Scala式の文法を「改行でもセミコロンでも終わることができる」と考えて次のように記述しても大丈夫でしょうか。

\begin{lstlisting}
expression_statement = expression (<SEMICOLON> | <LINE_TERMINATOR>)
\end{lstlisting}

\texttt{\textless{}LINE\_TERMINATOR\textgreater{}}は改行を表すトークンです。プラットフォームによって改行コードは異なるので、このように定義しておくと楽でしょう。このような規則でうまく先ほどの例すべてをうまく取り扱えるかといえば、端的に言って無理です。たとえば、次のようなコードを考えてみましょう。

\begin{codescreen}{}
println(
  "A")
\end{codescreen}

\texttt{(}の後に改行がきていますが、ここでは単なる空白文字として扱って無視してほしいわけです。しかし、上記の規則では改行がきた時点では\texttt{println(}までしか読み込んでいません。しかも、すでに改行文字は単なる空白と同列視できませんから、式の途中で不正なトークンがきたという扱いになり、構文解析エラーになってしまいます。

C系の言語では原則的には、字句解析時に改行もスペースも同じ扱いで処理しているので、このような「式の途中で改行がくる」ケースも特に工夫する必要がありませんでした。しかし、Scalaなどの言語における「改行」は式の途中では無視されるが文末では終端子にもなり得るという複雑な存在です。

言い換えると「文脈」を考慮して改行を取り扱う必要がでてきたのです。これ自体は文脈自由文法の範囲で取り扱うことが可能ですが、字句解析器が素朴に空白文字を無視するだけでは対応できません。なぜなら、通常の字句解析器は文脈を持たないからです。字句解析器は入力をトークンに分解するだけのもので、文脈情報を持たないため、改行が文の終わりかどうかを判断できません。

このような文法はどのようにすれば取り扱えるでしょうか。なかなか難しい問題ですが、大きく分けて二つの戦略があります。

一つ目は字句解析器に文脈情報を持たせる方法です。たとえば、「式」モードでは改行は無視されるが、「文」モードだと無視されないというふうにした上で、式が終わったら「文」モードに切り替えを行い、式が開始したら「式」モードに切り替えを行います。この方式を採用している典型的な言語がRubyで、Cで書かれた字句解析器には実に多数の文脈情報を持たせています。

\begin{codescreen}{}
// https://github.com/ruby/ruby/blob/v3_2_0/parse.y#L161-L181
/* examine combinations */
enum lex_state_e {
#define DEF_EXPR(n) EXPR_##n = (1 << EXPR_##n##_bit)
    DEF_EXPR(BEG),
    DEF_EXPR(END),
    DEF_EXPR(ENDARG),
    DEF_EXPR(ENDFN),
    DEF_EXPR(ARG),
    DEF_EXPR(CMDARG),
    DEF_EXPR(MID),
    DEF_EXPR(FNAME),
    DEF_EXPR(DOT),
    DEF_EXPR(CLASS),
    DEF_EXPR(LABEL),
    DEF_EXPR(LABELED),
    DEF_EXPR(FITEM),
    EXPR_VALUE = EXPR_BEG,
    EXPR_BEG_ANY  =  (EXPR_BEG | EXPR_MID | EXPR_CLASS),
    EXPR_ARG_ANY  =  (EXPR_ARG | EXPR_CMDARG),
    EXPR_END_ANY  =  (EXPR_END | EXPR_ENDARG | EXPR_ENDFN),
    EXPR_NONE = 0
};
\end{codescreen}

「改行で文が終わる」以外にもRubyはかなり複雑な構文解析を行っているため、このように多数の状態を字句解析器に持たせる必要があります。Rubyの文法は私が知る限りかなり複雑なものの一つなのでやや極端ですが、字句解析器に状態を持たせるアプローチは他の言語も採用していることが多いようです。

別のアプローチとしてPEGを使うという方法があります。PEGでは字句解析という概念自体がありませんから、式の途中に改行が入るというのも構文解析レベルで処理できます。Pythonでは3.9からPEGベースのパーサーが導入されました。これにより、Pythonのパーサーは改行の扱いなど、柔軟な構文規則を以前よりも直接的に記述しやすくなりました。

ともあれ、私達が普通に使っている「改行で文が終わる」ようにできる処理一つとっても厄介な問題だということがポイントです。

\section{Cのtypedef}
\index{C言語}
\index{typedef@\texttt{typedef}}
\index{しんぼるてーぶる@シンボルテーブル}
\index{ぶんみゃくいぞんせい@文脈依存性}

C言語の\texttt{typedef}文は既存の型に別名をつける機能です。C言語やC++言語のプログラムをバリバリ書いているプログラマならおなじみの機能です。\texttt{typedef}は

\begin{itemize}
\item
  移植性を高める
\item
  関数ポインタを使った読みづらい宣言を読みやすくする
\end{itemize}

\noindent といった目的で使われますが、この\texttt{typedef}文が意外に曲者だったりします。以下は\texttt{i}をint型の別名として定義するものですが、同時にローカル変数\texttt{i}を\texttt{i}型として定義しています。

\begin{codescreen}{}
typedef int i;
int main(void) {
        i i = 100; // OK
        int x = (i)'a'; // ERROR
        return 0;
}
\end{codescreen}

現実にこのようなコードの書き方をするかはともかく、

\begin{codescreen}{}
i i = 100;
\end{codescreen}

は明らかにOKな表現として解析してあげなければいけません。一方で、

\begin{codescreen}{}
int x = (i)'a';
\end{codescreen}

は構文解析エラーになります。\texttt{i\ i\ =\ 100;}という宣言がなければこの文も通るのですが、合わせて書けば構文解析エラーです。それ以前の文脈である識別子がtypedefされたかどうかで構文解析の結果が変わるのですからとてもややこしいです。

C言語ではこのようなややこしい構文を解析するために、構文解析器がtypedefした識別子のテーブルを持っておいて、それを使って構文解析を行うという手法を採用しています。

\section{Scala 3での「文頭に演算子がくる場合の処理」}
\index{Scala}
\index{えんざんし@演算子}
\index{かいぎょう@改行}
\index{せみころん@セミコロン}

6.4節で改行で文を終端する文法について説明しましたが、Scala 3はさらにややこしい入力を処理できなければいけません。たとえば、以下のような文を解釈できる必要があります。

\begin{codescreen}{}
val x = 1
      + 2
println(x) // 3
\end{codescreen}

ここで、xを3とちゃんと解釈するには\texttt{val\ x\ =\ 1}が改行で終わったから「文が終わった」と解釈せず、次のトークンである\texttt{+}まで見てから文が終わるか判定する必要があります。これまで試した限り、同じことができるのはJavaScriptくらいで、Ruby、Python、Kotlin、Go、Swiftなどの言語ではエラーになるか、\texttt{x\ =\ 1}で文が終わったと解釈され、\texttt{+\ 2}は別の文として解釈されるケースばかりでした。

この処理について、Scala 3言語仕様内の\href{https://www.scala-lang.org/files/archive/spec/3.4/01-lexical-syntax.html}{1.2
Newline Characters}に関連する記述があります。

\begin{quote}
Scala is a line-oriented language where statements may be terminated by
semi-colons or newlines. A newline in a Scala source text is treated as
the special token ``nl'' if the three following criteria are satisfied:
1. The token immediately preceding the newline can terminate a
statement. 2. The token immediately following the newline can begin a
statement and is not a leading infix operator. 3. The token appears in a region where newlines are enabled.
The tokens that can terminate a statement are: literals, identifiers and
the following delimiters and reserved words:
\end{quote}

これを意訳すると、通常の場合はScalaの文はセミコロンまたは改行で終わることができるが、次の三つの条件すべてを満たしたときのみ、特別なトークン\texttt{nl}として扱われる、ということになります。

\begin{enumerate}
\item
  改行の直前のトークンが「文を終わらせられる」ものである場合
\item
  改行の直後のトークンが「文を始められる」ものであり、行頭の中置演算子ではない場合
\item
  改行が「利用可能」になっている箇所にあらわれたものである場合
\end{enumerate}

たとえば、以下のScalaプログラムについていうと、最初の改行は\texttt{nl}トークンになりません。\texttt{+}が行頭の中置演算子なので、条件2が満たされないからです。

\begin{codescreen}{}
val x = 1
      + 2
\end{codescreen}

ちなみに、調査を開始する時点ではScalaの文法の基本文法を継承したKotlinでも同じようになっていると思っていたのですが、一行目で文が終わると解釈されてしまいました。

\begin{codescreen}{}
val x = 1
      + 2 // Kotlinでは+ 2は単独の式として解釈されてしまう
println(x) // 1
\end{codescreen}

これらに対する体系的な調査は行っていませんが、Scalaは文法上、ところどころにトリッキーな機能があり、改行の扱いも特殊なものになっています。

\section{C++のテンプレート構文}
\index{C++}
\index{てんぷれーと@テンプレート}
\index{じぇねりくす@ジェネリクス}
\index{みぎしふとえんざんし@右シフト演算子}
\index{ぶんみゃくいぞんせい@文脈依存性}

C++のテンプレートは型パラメータを使ったジェネリックプログラミングを可能にする強力な機能ですが、その構文は構文解析の観点から見ると興味深い問題を含んでいます。特に有名なのが「\texttt{\textgreater{}\textgreater{}}の曖昧性」です。

Javaのジェネリクスでいうと
\texttt{List\textless{}List\textless{}String\textgreater{}\textgreater{}}
のような書き方に相当しますが、C++では歴史的な経緯から特別な問題がありました。

C++03までは、次のようなネストしたテンプレートの宣言はコンパイルエラーになっていました。

\begin{codescreen}{}
std::vector<std::vector<int>> matrix;  // C++03ではエラー
\end{codescreen}

なぜかというと、\texttt{\textgreater{}\textgreater{}}が右シフト演算子として解釈されてしまうからです。そのため、C++03では次のようにスペースを入れる必要がありました。

\begin{codescreen}{}
std::vector<std::vector<int> > matrix;  // C++03ではOK
\end{codescreen}

この問題は単純な字句解析の問題のように見えますが、実はそうではありません。次の例を考えてみましょう。

\begin{codescreen}{}
template<int N>
struct A {
    static const int value = N;
};

// これは A<1>::value >> 2 （右シフト）
int x = A<1>::value >> 2;

// これは vector<vector<int>> （テンプレートの終端）
std::vector<std::vector<int>> v;
\end{codescreen}

同じ\texttt{\textgreater{}\textgreater{}}という文字列が、文脈によって右シフト演算子として解釈されたり、テンプレート引数リストの終端記号2つとして解釈されたりする必要があるのです。これは明らかに文脈自由文法の範囲を超えており、構文解析器は現在の文脈（テンプレート引数リストの中にいるかどうか）を追跡する必要があります。

C++11以降では、この問題に対処するため、構文解析器がテンプレート引数リストの文脈では\texttt{\textgreater{}\textgreater{}}を\texttt{\textgreater{}\ \textgreater{}}として解釈するように標準が変更されました。しかし、これは構文解析器の実装をより複雑にします。たとえば、次のような入れ子になったケースも正しく処理する必要があります。

\begin{codescreen}{}
std::map<int, std::vector<std::vector<int>>> nested;  // C++11以降はOK
\end{codescreen}

さらに複雑なのは、\texttt{\textgreater{}\textgreater{}=}（右シフト代入演算子）や\texttt{\textgreater{}\textgreater{}\textgreater{}}（一部の拡張で使われる）なども考慮する必要があることです。構文解析器は、テンプレート引数リストの文脈を正確に追跡し、適切にトークンを分割する必要があります。

\section{正規表現リテラルの曖昧性}
\index{せいきひょうげんりてらる@正規表現リテラル}
\index{JavaScript}
\index{じょざんえんざんし@除算演算子}
\index{あいまいせい@曖昧性}
\index{じくかいせき@字句解析}
\index{こうぶんかいせき@構文解析}

JavaScriptには正規表現リテラルという便利な記法があります。Javaでは
\texttt{Pattern.compile("{[}a-z{]}+",\ Pattern.CASE\_INSENSITIVE)}
のように書くところを、JavaScriptでは以下のように簡潔に書けます。

\begin{codescreen}{}
const pattern = /[a-z]+/i;  // 大文字小文字を区別しない英字のパターン
\end{codescreen}

しかし、この\texttt{/}記号は除算演算子としても使われるため、構文解析上の曖昧性が生じます。

\begin{codescreen}{}
// これは正規表現リテラル
const regex = /ab+c/;

// これは除算演算
const result = a / b + c / d;
\end{codescreen}

さらに厄介なのは、次のようなケースです。

\begin{codescreen}{}
// これは何でしょう？
return /regex/i.test(str);  // 正規表現リテラル
return a / b / c;            // 除算演算
\end{codescreen}

この曖昧性を解決するには、構文解析器は現在の文脈を理解する必要があります。JavaScriptの仕様では、正規表現リテラルが出現できる位置と、除算演算子が出現できる位置を厳密に定義しています。

一般的なルールとしては以下のようになります。

\begin{itemize}
\item
  式の開始位置では\texttt{/}は正規表現リテラルの開始
\item
  式の途中では\texttt{/}は除算演算子
\end{itemize}

しかし、実際にはもっと複雑で、直前の構文要素が式を完結しているかどうか（識別子・リテラル・後置インクリメント・括弧で囲まれた式など）で判断する必要があります。

\begin{codescreen}{}
// 関数呼び出し ( 式を完結する ) の直後 → 除算
func() / 2

// 代入演算子の直後 ( 式の開始位置 ) → 正規表現
const x = /pattern/

// return の直後 → 正規表現
return /pattern/g
\end{codescreen}

ややこしいのは、同じ\texttt{)}でも\texttt{if\ (cond)}のような制御構文の閉じ括弧の直後では\texttt{/}が正規表現として解釈される点です（例:
\texttt{if\ (a)\ /pat/.test(s)}）。つまり、字句解析器は単に前のトークンの種類だけでなく、その\texttt{)}が式の終端か制御構文の構成要素かまで把握する必要があり、字句解析と構文解析が密接に連携しなければなりません。

\section{エラーリカバリ}
\index{えらーりかばり@エラーリカバリ}
\index{ぱにっくもーど@パニックモード}
\index{ふれーずれべるりかばり@フレーズレベルリカバリ}
\index{えらーせいせいきそく@エラー生成規則}
\index{ぐろーばるていせい@グローバル訂正}
\index{どうきとーくん@同期トークン}
\index{IDE}
\index{いんくりめんたるかいせき@インクリメンタル解析}

構文解析の途中でエラーが起きることは普通にあります。構文解析中のエラーリカバリについては多くの研究があるものの、コンパイラの教科書で構文解析アルゴリズムでのエラーリカバリについて詳しく言及されることはあまりありません。推測ですが、構文解析において２つ目以降のエラーは大抵最初のエラーに誘発されて起こるということや、どうしても経験則に頼った記述になりがちなため、教科書で言及されることが少ないのでしょう。大抵の言語処理系で構文解析中のエラーリカバリについては大したことをしていなかったという歴史的事情もあるかもしれません。

しかし、現在は別の観点から構文解析中のエラーリカバリが重要性を増してきています。IDEあるいはテキストエディタの拡張としてIDE類似の機能を提供することが一般的になったため、「構文解析エラーになるが、それっぽくなんとか構文解析をしてほしい」という強いニーズがあるからです。

ユーザーがコードを書いている最中は、一時的に文法的に正しくない状態になることが頻繁にあります。IDEがそのような状況でも構文ハイライト、コード補完、リアルタイムのエラー表示などの機能を提供し続けるためには、エラーが発生しても即座に解析を中断するのではなく、可能な限り解析を継続し、後続のエラーも検出できるような仕組み、すなわちエラーリカバリ機構が不可欠です。

エラーリカバリにはいくつかの代表的な手法があります。

\begin{itemize}
\item
  \textgt{パニックモード }\textbf{\sffamily (Panic Mode):}
  最も単純な手法の一つです。エラーを検出したら、セミコロン（\texttt{;}）や閉じ波括弧（\texttt{\}}）のような、文やブロックの区切りとなる「同期トークン」が見つかるまで、入力トークンを読み飛ばします。同期トークンが見つかったら、そこから解析を再開します。

  C言語ふうのコード \texttt{x\ =\ a\ +\ *\ b;\ y\ =\ c;} で、\texttt{*}
  が予期しないトークンとしてエラーになった場合、パニックモードでは\texttt{*}と\texttt{b}を読み飛ばし、次の同期トークンである\texttt{;}を見つけて解析を再開します。これにより、\texttt{y\ =\ c;}
  の解析は行われますが、\texttt{* b}
  に関するエラーの詳細は失われる可能性があります。
\item
  \textgt{フレーズレベルリカバリ}\textbf{\sffamily (Phrase-Level Recovery):}
  エラー箇所の周辺で、局所的な修正を試みる手法です。たとえば、不足しているセミコロンを補ったり、予期しないトークンを削除したり、期待されるトークンに置き換えたりします。

\begin{codescreen}{}
x = a + b y = c;
\end{codescreen}

  というコードで、\texttt{b} と \texttt{y}
  の間にセミコロンが欠落している場合、パーサーは \texttt{y}
  が予期しないトークンであると判断します。フレーズレベルリカバリでは「文の終わりにはセミコロンが期待される」という知識に基づき、\texttt{b}
  の後にセミコロンを挿入して \texttt{x\ =\ a\ +\ b;\ y\ =\ c;}
  として解析を試みるかもしれません。

  あるいは、Javaのメソッド呼び出しで

\begin{codescreen}{}
myObject.method(arg1 arg2)
\end{codescreen}

のようにカンマが抜けている場合、\texttt{arg1} と \texttt{arg2}
の間にカンマを補って \texttt{myObject.method(arg1,\ arg2)}
として解釈を試みる、といった具合です。

JSONの配列 \texttt{{[}1,\ ,\ 2{]}} で余分なカンマがある場合、削除して
\texttt{{[}1,\ 2{]}}
として解析を続けるかもしれません。パニックモードよりは洗練されていますが、どのような修正を行うかの判断が難しく、実装が複雑になりがちです。

\item
  \textgt{エラー生成規則}\textbf{\sffamily (Error Productions):}
  文法にあらかじめよくあるエラーパターンに対応する生成規則を追加しておく手法です。たとえば、「\texttt{if\ (condition)\ statement}」という正しい規則に加えて、「\texttt{if\ condition)\ statement}」（開き括弧が欠落）のようなエラー用の規則を定義しておきます。これにより、特定のエラーを「受理」し、解析を継続できます。

  Yaccのようなツールでは、\texttt{error}
  トークンを使ってエラー規則を定義できます。

\begin{codescreen}{}
statement: IF '(' expr ')' statement
         | IF error ')' statement {
             yyerror("Missing opening parenthesis in if statement");
           }
         | /* ... other rules ... */
         ;
\end{codescreen}

  この例では、\texttt{if} の後に開き括弧 \texttt{(} がない場合に
  \texttt{error}
  トークンがマッチし、エラーメッセージを出力しつつ、\texttt{)}
  以降の解析を継続しようとします。多くのエラーパターンを網羅しようとすると文法が複雑になりますが、特定のエラーに対しては効果的です。

\item
  \textgt{グローバル訂正}\textbf{\sffamily (Global Correction):}
  理論的には最も強力な手法で、入力文字列全体に対して、最小限の修正（挿入、削除、置換）で文法的に正しい文字列に変換する方法を探します。

  \texttt{if\ x\ \textgreater{}\ 0)\ \{\ ...\ \}}
  という入力に対し、グローバル訂正は開き括弧 \texttt{(}
  を挿入するのが最小の修正であると判断するかもしれません。しかし、入力全体を考慮して最適な修正を見つけるのは計算量的に非常に困難であり、実用的なコンパイラやIDEでこの手法が全面的に採用されることはまれです。
\end{itemize}

IDEのような環境では、これらの手法を組み合わせたり、部分的な構文木（エラー箇所を含むかもしれないが、解析できた部分）を構築したり、インクリメンタルな解析（変更箇所だけを再解析する）を行ったりすることで、ユーザーが編集中でも可能な限り正確な情報を提供しようと試みています。

たとえば、Javaのクラスを書いている途中で

\begin{codescreen}{}
class MyClass { 
  public void myMethod() { 
    ... 
  }
// ここで閉じ括弧 } を入力し忘れた場合
\end{codescreen}

\noindent と入力し、最後の閉じ括弧 \texttt{\}} を入力し忘れている場合でも、IDEは
\texttt{myMethod}
の本体部分についてはある程度解析を試み、メソッド内の変数に対するコード補完や型チェックを行おうとします。エラーリカバリ機構が、エラー箇所を特定しつつも、それ以外の部分については解析を継続し、部分的な構文情報を抽出しているためです。

エラーリカバリは、単にエラーを見つけるだけでなく、その後の解析をどう継続し、ユーザーにどのようなフィードバックを与えるかという、実践的で奥深い問題領域なのです。

\section{まとめ}

本章では現実の構文解析で遭遇する問題について、いくつかの例を挙げて説明しました。

筆者が大学院博士後期課程に進学した頃「構文解析は終わった問題」と言われたのを覚えています。特にその頃だと主流はLALRなどのボトムアップ構文解析だという認識のある人が研究者にも多かったのを覚えています。実際にはその後もANTLRの\texttt{LL(*)}、それに続く\texttt{ALL(*)}アルゴリズムのような革新が起きており、ANTLRはいろいろな言語のパーサーに採用されています。Python 3.9におけるPEGパーサーの採用のような例もあります。

また、本章の例のように従来の構文解析法単体では取り扱えない部分をアドホックに各プログラミング言語が補っている部分は未だに多くありますし、IDEのようなケースでわかるように、エラーリカバリの重要性も増してきています。

その上で「構文解析が終わった問題」にならなかった理由について改めて考えてみます。

一番大きいのは、当初の想定と違って「プログラミング言語は文脈自由言語で表しきれなかった」ということでしょう。もちろん、文脈自由言語の範囲に納めるように文法を制限することは可能ですが、便利な表記を許していくとどうしても文脈自由言語から「はみ出て」しまいます。

近年重要性が増してきている構文解析におけるエラーリカバリについても、エラーリカバリの研究が「実質的に」長年停滞することになったのも、当時は今ほど高度なIDEがなかったことも大きいように思います。翻って、現代では、開発者の多くがIDEを使ったリッチな環境で作業しており「全体を正しく構文解析できないが、だいたい正しく構文解析してほしい」というニーズが広くあります。

ともあれ、構文解析のような古典的な分野でも、時代の変遷によって新しい課題が出てきているわけで、古典的なコンパイラの教科書には書かれない、本章のような「泥臭い」側面を知っておくのは有益だと筆者は考えます。

\backmatter
\chapter*{おわりに}\label{afterword}\addcontentsline{toc}{chapter}{おわりに}\markboth{おわりに}{おわりに}

ここまでで構文解析の世界を概観してみましたがいかがでしたか？構文解析、特に非自然言語の構文解析というのは地味なもので、パーサージェネレータの発展などもあり、20世紀末には「構文解析はもう終わった問題だ」という人もいました。その一方で2000年代以降になってもPEGの提案があり、Pythonの構文解析器に採用されるまでになりましたし、\texttt{LL(*)}や\texttt{ALL(*)}のような革新的なアルゴリズムが生み出されています。それも、主流であった上向き型の構文解析でなく下向き型の構文解析で、です。

とはいえ地味なものは地味であり、プログラミング言語処理系を構成するコンポーネントという観点からは「脇役」という印象は否めません。ただ、私たちはプログラミング言語を書いているときは、コンパイラの内部表現や抽象構文木と対話しているわけではありません。プログラマが直接対話する相手はプログラミング言語の具象構文であり、具象構文はプログラミング言語の「UI」を担当する部分といえるでしょう。通常のアプリケーションでUIが軽視されるべきでないのと同様にやはり具象構文も軽視されるべきでないと私は思いますし、よりよい具象構文の設計には構文解析の知識が助けになると信じています。

ここまで本書では、LL法やLR法といった古典的なアルゴリズムから、PEGとPackrat Parsing、構文解析器生成系、そして現実のプログラミング言語が抱える泥臭い問題までを見てきました。しかし、構文解析の世界は広く、本書で十分に扱えなかったテーマも数多く残っています。最後に、その一端を紹介しておきましょう。この本を読み終えた皆さんが次へ進むための、ささやかな地図になれば幸いです。

\begin{itemize}
\item
  \textgt{曖昧な文法とgeneralized parsing：}

  本書で扱った手法は、いずれも「1つの解析結果を効率よく求める」ことを前提としていました。これに対して、曖昧な文脈自由文法をそのまま受け取り、あり得る解析結果をすべて求める系譜があります。Earley法、GLR法、GLL法などがその代表で、generalized parsingと総称されます。第4章と第5章で名前だけ登場したGLR法も、元々は自然言語を解析するために生まれたものでした。プログラミング言語の世界では曖昧さは避けるべきものとして扱われますが、自然言語ではむしろ曖昧さこそが本質であり、それを正面から引き受けるための道具立てがこの系譜だ、と見ることもできます。プログラミング言語の枠を越えて構文解析を眺めてみたい方には、この系譜をおすすめします。参考文献の「Parsing Techniques」は、この世界まで含めて扱っている数少ない本です。

\item
  \textgt{インクリメンタル構文解析：}

  IDEやエディタは、キー入力のたびにファイル全体を解析し直しているわけにはいきません。編集された箇所だけを再解析して構文木を更新する技術がインクリメンタル構文解析です。本書では序章のコラムで触れるに留めましたが、構文解析器生成系tree-sitterのようにこの技術を核としたツールも普及しており、LSPが当たり前になった現在、実用上の重要性はますます高まっています。差分だけを解析し直すには、前回の解析結果をどう安全に再利用するかという、本書では扱わなかった種類の工夫が必要になります。

\item
  \textgt{エラーリカバリ：}

  第6章で駆け足で紹介しましたが、本来これだけで一冊の本が書ける主題です。正しい入力だけを相手にする構文解析器は比較的きれいに書けますが、IDEなどの開発環境が相手にするのは、ほとんどの場合「書きかけの壊れたプログラム」だからです。壊れた入力から、どれだけ意味のある構文木を回収できるかが、実用の構文解析器の品質では効いてきます。

\item
  \textgt{形式言語理論：}

  第3章で概観した正規言語や文脈自由言語といった言語クラスには、対応する計算モデル（機械）があることが知られています。正規言語は有限オートマトンで、文脈自由言語は有限オートマトンにスタックを追加した\textgt{プッシュダウンオートマトン}で認識できます。有限オートマトンで括弧の対応が扱えなかったのは、まさにこのスタックの有無の差なのです。さらに外側には、線形拘束オートマトンやチューリングマシンといった、より強力な計算モデルの世界が続いています。「なぜ\(a^n b^n\)は文脈自由言語なのに\(a^n b^n c^n\)はそうではないのか」といった問いに正面から答えたくなったら、形式言語理論の教科書に進むときです。

\end{itemize}

本書で学んだ文脈自由文法や各種アルゴリズムの理解が、これらのテーマに挑むときの土台になります。

以下では、こうしたテーマをさらに学ぶための文献を紹介します。

\section*{古典的名著・専門書}\label{a.1}\addcontentsline{toc}{section}{古典的名著・専門書}\markboth{おわりに}{古典的名著・専門書}

\begin{itemize}
\item
  A.V.エイホ、M.S.ラム、R.セシィ、J.D.ウルマン，『コンパイラ -
  原理・技法・ツール』（第2版）、サイエンス社、2009年（通称ドラゴンブック）

  \begin{itemize}
  \item
    コンパイラ構築に関する標準的な教科書です。字句解析、構文解析（LL、LR）、意味解析、コード生成など、コンパイラの全般的なトピックを網羅。理論的背景もしっかり解説されています。中級者以上向けです。
  \end{itemize}
\item
  J.ホップクロフト、R.モトワニ、J.ウルマン、『オートマトン 言語理論 計算論 Ⅰ』(第2版)
  、サイエンス社、2003年

  \begin{itemize}
  \item
    オートマトンと形式言語の理論に関する大学生レベル以上向け教科書。正規言語、文脈自由言語など、計算理論の基礎をしっかり学べます。数学的な厳密さを求める方向け。
  \end{itemize}
\item
  J.ホップクロフト、R.モトワニ、J.ウルマン、『オートマトン 言語理論 計算論 Ⅱ』(第2版)
  、サイエンス社、2003年

  \begin{itemize}
  \item
    同上。Ⅱではチューリングマシン、決定不能性、計算複雑性などを取り扱っています。Ⅱは計算理論のより深い部分に踏み込んでいます。
  \end{itemize}
\item
  Dick Grune, Ceriel J.H. Jacobs, \emph{Parsing Techniques: A Practical
  Guide} (2nd Edition), Springer, 2008年

  \begin{itemize}
  \item
    書名どおり、さまざまな構文解析技術に特化した書籍です。LL、LRだけでなく、アーリー法、GLR、CYK法など、高度なアルゴリズムや曖昧性のある文法の扱いについても詳しく解説されています。構文解析を専門的に極めたい方向け。
  \end{itemize}
\end{itemize}

\section*{特定の技術に関する論文・資料}\label{a.2}\addcontentsline{toc}{section}{特定の技術に関する論文・資料}\markboth{おわりに}{特定の技術に関する論文・資料}

\begin{itemize}
\item
  Bryan Ford, ``Parsing Expression Grammars: A Recognition-Based
  Syntactic Foundation'', 2004

  \begin{itemize}
  \item
    PEGを提案したオリジナルの論文です。PEGの形式的な定義、操作的意味論などについて説明があります。理論的な背景を深く理解したい方向け。
  \end{itemize}
\item
  Tim A. Wagner, Susan L. Graham, ``Efficient and Flexible Incremental
  Parsing'', ACM Transactions on Programming Languages and Systems,
  20(5), 1998

  \begin{itemize}
  \item
    インクリメンタル構文解析の代表的な論文です。編集された箇所だけを安全に再解析するための理論と実装を詳述しており、tree-sitterをはじめとする現代のツールの基礎となっています。
  \end{itemize}
\item
  Lukas Diekmann, Laurence Tratt, ``Don't Panic! Better, Fewer, Syntax
  Errors for LR Parsers'', ECOOP 2020

  \begin{itemize}
  \item
    LR構文解析のエラーリカバリを実用的な品質へ引き上げた比較的新しい論文です。エラーリカバリ研究の現在地を知る入口として適しています。オープンアクセスで公開されています。
  \end{itemize}
\item
  Terence Parr, \emph{The Definitive ANTLR 4 Reference}, Pragmatic
  Bookshelf (2nd Edition), 2013年

  \begin{itemize}
  \item
    ANTLR v4の作者自身による解説書です。ANTLRの文法定義、使い方、\texttt{ALL(*)}アルゴリズムの概要、実践的なパーサー構築のテクニックが豊富。ANTLRを使いこなしたいなら必読。
  \end{itemize}
\item
  ANTLR公式サイト、\url{https://www.antlr.org/}

  \begin{itemize}
  \item
    ANTLRのドキュメント、チュートリアル、文法リポジトリなど。最新情報やコミュニティのサポートも得られます。
  \end{itemize}
\item
  JavaCC公式サイト、\url{https://javacc.github.io/javacc/}

  \begin{itemize}
  \item
    JavaCCのドキュメント、チュートリアル、FAQなど。
  \end{itemize}
\item
  GNU Bisonマニュアル、\url{https://www.gnu.org/software/bison/manual/}

  \begin{itemize}
  \item
    Bison (Yacc互換)の詳細なマニュアル。LALR(1)やGLRパーサーの生成方法、文法定義の書き方などが学べます。
  \end{itemize}
\item
  tree-sitter公式サイト、\url{https://tree-sitter.github.io/tree-sitter/}

  \begin{itemize}
  \item
    インクリメンタル構文解析とエラーリカバリを核とした、エディタ向けの構文解析器生成系です。多くのエディタやコード解析ツールで採用されています。
  \end{itemize}
\end{itemize}

これらの資料を通じて構文解析の世界への探求をさらに深めていただければ幸いです。

% \chapter*{著者プロフィール}\label{author}\markboth{著者プロフィール}{著者プロフィール}

\noindent
\textbf{水島 宏太（みずしま こうた）}

株式会社ネクストビート所属。筑波大学に在学中からプログラミング言語や構文解析に興味を持ち、大学院ではParsing Expression Grammarの研究を行い、博士（工学）を取得。修了後はアカデミアと接点を持ちつつも、Web開発を中心とした会社で開発や教育業務に携わる。情報処理学会プログラミング研究会において編集委員を務める。無類のプログラミング言語マニアでトイ言語を含めると1ダース以上のプログラミング言語を作った経験がある。

関わった書籍として、共著の『実践Scala入門』（技術評論社）、監訳に『Scala スケーラブルプログラミング 第4版』（インプレス）など。Twitter/GitHubでは@kmizuのIDで活動。

\printindex
% \thispagestyle{empty}

\begin{spacing}{0.7}
\noindent \textgt{\footnotesize ●本書に対するお問い合わせは、電子メール（\texttt{zup-gr@zen.ac.jp}）にてお願いいたします。\\
但し、本書の記述内容を越えるご質問にはお答えできませんので、ご了承ください。}
\end{spacing}


\null\vfill

\noindent \textgt{\large \ruby{構文}{こうぶん}\ruby{解析}{かいせき}のしくみ}

\vspace{1mm}

\noindent \textgt{\ruby{電卓}{でんたく}から\ruby{言語}{げんご}\ruby{処理系}{しょりけい}まで\ruby{自分}{じぶん}で\ruby{書}{か}く}

\vspace{2mm}

\noindent {\small 2026年9月28日　初版発行}\\

\vspace{2mm}

\begin{spacing}{1.1}
\noindent 著／\ruby{水島}{みずしま}\ruby{宏太}{こうた}\\
\end{spacing}

\vspace{1mm}

\noindent 発行／ZEN大学出版会
\begin{spacing}{0.7}
\noindent {\footnotesize 〒 249--0007} \\
\noindent {\footnotesize 神奈川県逗子市新宿3丁目12-11} \\
\noindent {\footnotesize 電子メール \texttt{zup-gr@zen.ac.jp}} \\
\noindent {\footnotesize \texttt{https://zen.ac.jp/}} \\
\end{spacing}
\noindent 発売／株式会社KADOKAWA
\begin{spacing}{0.7}
\noindent {\footnotesize 〒 102--8177} \\
\noindent {\footnotesize 東京都千代田区富士見2--13--3} \\
\noindent {\footnotesize KADOKAWA購入窓口　0570--002--008（ナビダイヤル）}\\
\noindent {\footnotesize \texttt{https://www.kadokawa.co.jp/}} \\
\end{spacing}

\noindent 印刷・製本／株式会社リーブルテック \\
\noindent {\footnotesize Printed in Japan} \\

\vspace{2mm}

\begin{spacing}{0.8}
\noindent {\footnotesize 本書（ソフトウェア／プログラム含む）の無断複製（コピー、スキャン、デジタル化等）並びに無断複製物の譲渡および配信は、著作権法上での例外を除き禁じられています。また、本書を代行業者などの第三者に依頼して複製する行為は、たとえ個人や家庭内での利用であっても一切認められておりません。}\\
\noindent {\footnotesize 定価はカバーに表示してあります。}
\end{spacing}

\vspace{2mm}


\noindent {\footnotesize ISBN978--4--04--901207--1 C3004} \\

\begin{spacing}{0.7}
\noindent {\footnotesize ZEN大学出版会}\\
\noindent {\footnotesize 編集／鈴木嘉平}
\end{spacing}

%\include{okuduke_ebook}

\end{document}
